\chapter{C)~上流社会}

\section{VII.~研讨会——或Mebkhout的层与反常性}

\subsection{不义——或归来的意义}

\notehead{75}{}

\dmark{}\origpage{285}(5月2日)看来我学无止境!我刚得知两份文本,它们(至少对我而言)以料想不到的光,照亮了前文已谈过的那场“藏掖”(对Mebkhout著作的藏掖)(“当值无名之辈与好上帝定理”一注,即注48$'$)。文本所揭示的,正是那两位显赫的同行兼前学生所扮演的角色——我曾确认他们心存轻蔑的漠然,却不曾怀疑其职业诚意。两份文本都属于\textbf{吕米尼研讨会}(1981年7月6日至11日)的会议录;研讨会题为\textbf{《奇异空间上的分析与拓扑》},刊于Astérisque~n\textsuperscript{o}~100(1982)。

第一份文本是研讨会的导言,署名\textbf{B.~Teissier}与\textbf{J.L.~Verdier}(即那位曾充当 Z.~Mebkhout正式的论文导师的同一位)。这份一页半的文本,开篇先解释某条“所谓Riemann--Hilbert对应”,它分明注定要在研讨会中扮演头等重要的角色(而它不是别的,正是“好上帝定理”,又名Mebkhout定理)。在这一对应里(这正是它的魅力与深度所在,也是必须引入导出范畴的原因):一个正则完整\textbf{模}(即集中在零次的一个正则完整复形),对应着一个由$\mathbb{C}$-向量空间的层构成的可构造复形;据称,它可由纯拓扑的性质来刻画,而这些性质,对一个任意域上定义的、不必光滑的簇上的平展层的可构造复形,同样保有意义。据解释,这就是研讨会“主旋律”——“\textbf{反常性、相交复形、纯粹性}”——的出发点:所谓“\textbf{反常}”的(\dmark{}\origpage{286}层)复形\fnote{1}{(*)}{(5月4日)关于这一奇怪的应用,见注n\textsuperscript{o}~76“反常性”。}不是别的,正是那些在“实质上”对应于(“Mebkhout式地”)最简单的正则完整微分算子复形——后者借一个单独的$\mathcal{D}$-模来表达——的层复形。

第二份文本,是\textbf{A.A.~Beilinson、J.~Bernstein与P.~Deligne}关于反常层的长文的一部分\fnote{2}{(**)}{(5月4日)我此后收到了全文;它确认了我手头这一部分已经向我表明的东西。};导言里称之为研讨会的中心工作。如同目录与我手头的其他页码所见证的,这篇文章宣告了导出范畴与三角范畴以骤然的强力重返公共广场——追随着Mebkhout那些鲜为人知的著作,与那道著名的“所谓Riemann--Hilbert”定理。

难以置信,却千真万确:两份文本里,Z.~Mebkhout的名字都缺席,正如它在参考文献里也缺席。我要说明:不仅 J.L.~Verdier对Mebkhout的工作了如指掌(这是理所当然的!),Deligne也同样了如指掌;而且,对这样一个对数学时事如此灵通的人——况且这还是与他切身相关的主题——甚至难以设想会是别样\fnote{3}{(***)}{我尤其提请注意:Mebkhout的著作与其“好上帝定理”,相对于Deligne(1969年的)较早的工作,是一个决定性的进展——而Deligne没有发表那些工作。对此可参见已引用的注n\textsuperscript{o}~48$'$。}。

B.~Teissier\fnote{4}{(****)}{(6月12日)B.~Teissier久已关注Mebkhout的工作,并因此是极少数对他持鼓励态度的人之一。因此,他对这桩巧取了如指掌,而且是明知故犯地给予了配合。他曾向Mebkhout辩白说,反正他“什么也改变不了”。}以及吕米尼研讨会的其他与会者情况如何——尤其是与Deligne共同署名前引文章的那两位\fnotevariant{5}{(****)}{(*****)}{(5月28日)我此后得知:A.A.~Beilinson与 J.~Bernstein是通过 P.~Deligne(1980年10月)与Mebkhout本人(1980年11月,在莫斯科的一场报告里,内容极为详尽)得知Mebkhout的结果的。早在1981年6月(原书此处如此)的吕米尼研讨会之前,这两位作者就在他们关于著名的Kazhdan--Lusztig猜想的证明里,本质上使用了好上帝定理。可与“一种不义与无力的感觉”一注(注n\textsuperscript{o}~44$''$)里所引的Zoghman~Mebkhout的信对读。(6月3日)关于研讨会全体与会者之团结的更多细节,见下一条注“研讨会”,n\textsuperscript{o}~75$'$。}——我不知道。似乎与会者中没有一个人曾如此好奇,想去了解把他们动员起来的那些观念与关键定理的父系渊源。

我料想,事情不言自明——有点(很有点)像次年将为Motif重返同一“公共广场”揭幕的Lecture~Notes的LN~900那样\fnote{6}{(******)}{对此可参见注n\textsuperscript{o}~51、52、59。}:父权属于发起并主持研讨会的那些光彩夺目的数学家中最光彩的那一位。对所有人来说,可以肯定的至少是:那绝不是Riemann,也不是Hilbert——否则,那场光辉的研讨会\dmark{}\origpage{287}就该在1900年召开,而不是1981年,即在无名的学生于Jean-Louis~Verdier门下答辩论文的两年之后。

我在这里得以确认的这类战役,如今也许已是家常便饭\fnote{7}{(*)}{我还想到另外两场同方向的“战役”,它们落实为LN~900(参见前一条脚注)的出版,以及五年前的SGA~4½(对此可参见注n\textsuperscript{o}~67、67$'$、68、68$'$)。(5月9日)关于与前面几场紧密相连的第三场这样的战役,见“好的参考文献”一注(n\textsuperscript{o}~82)——这次是 J.L.~Verdier手笔的另一篇“值得纪念的文章”。},并且完全得到认可,只要实践它的是占据高位的数学家,而付出代价的那位只是个无名之辈(尽管人们好心地邀他来,让他高兴)。这些实践者中的一位,倘若凭其手段也凭其作品,俨然一位大数学家(这使他立刻高悬于一切怀疑之上),也丝毫不改变事情的本性。肯定,我太守旧了——在我那个年代,这类战役叫作\textbf{巧取};而这一桩,在我看来,是容忍它的那一代数学家的一场\textbf{耻辱}。

天才的光辉并不减轻这样的耻辱;它给它添上一个新的维度——在我们的科学史上或许是独一无二的维度\fnote{8}{(**)}{在数学之外的任何一门科学或艺术的历史里,我也从未听说过类似的事情。}。它\dmark{}\origpage{288}可以使人隐约看到:在这一举动的荒谬与无端的表象背后(它出自一个命运已对他奢侈到无边、却仍然耽于巧取的人……),或许有别种力量在起作用,而不只是出风头的欲望,或羞辱、或使那个自觉无援无声的人绝望的无谓欲望。

既然我分明已深陷“风俗画”之中,我顺带指出(几乎当作理所当然):我的名字,在前引文本里同样缺席。然而,我能欣然确认:前引文章里没有一页(在我手头的那些页码里\fnote{9}{(*)}{(5月4日)还有别的页码,我此后也知道了。})不是深深植根于我的著作、带着它的印记的——直到我引入的那些记号,直到那些随处可见的观念所用的名字——那些名字,是我在它们尚未得名时就认识它们的时候,给它们取的。当然,合乎规矩的调整是有的:比如我在五十年代提炼的双对偶定理\fnote{10}{(**)}{平展对偶理论同样如此:在其慷慨的朋友Deligne的笔下,它也改姓“Verdier对偶”!},就为这场合改名为“Verdier对偶”——还是那个Verdier,没有搞错……\fnote{11}{(***)}{(5月5日)可与注n\textsuperscript{o}~48$'$、63$''$对读。在这场绵延将近十五年的漫长的葬礼中,也在这场刚刚于上个月、由主要的“被预葬者”所发现的葬礼中, J.L.~Verdier分明始终与他显赫的朋友不可分割:在这可悲的场合,他不计其数地为他撒下那些必不可少的万千花束。}

不过,我的名字不可能连间接出现都没有:经由对几部至今无可替代的文本(尽管有SGA~4½——它还不太胜任其使命)的偶尔引证,即EGA与SGA。(在“SGA~=~Bois~Marie代数几何研讨班”这个缩写的解释里,我的名字当然不出现;但在EGA里,人诚实不诚实总藏不住:全称是给出的,作者名单里包括我……)另一个打动我的细节,它见证了下葬综合征的强迫之力——在一位分明毫无“强迫症倾向”的人身上:我看到的仅有的两处对SGA的引证,每次都恪尽职守地特别写明“SGA~4里M.~Artin的定理……”,生怕哪位多心的读者会想:该定理说不定出自那位被小心隐去的人——尽管天知道,那场报告分明是做过了的,而且出自一位叫得出名字的作者!(注77)

\dmark{}\origpage{289}这一切,想必人们会认为,在今天的“上流社会”里,本是兵家常事。这场小规模交战虽然不能让我高兴(它也不是为此而设的……),却并不真正损害那位被预葬者:他那象征性的遗骸,就这样被交付给这场抢夺集市的种种偶然——而我不过两个星期前,才惊叹地发现这场集市。它并不以无力之中遭受的“不义”之感啃噬我的生活。它没有折断把我推向数学事物、推向周遭世界之事物的喜悦与冲劲;它没有在我身上烧掉这些事物的精微的美。我可以自认为幸运——而我也\textbf{确实}幸运……

我也为我那意味曾向我滑脱的意外的“归来”而喜悦。倘若这归来只需教给我这些天里已经学到的东西,它就不曾虚度——它已经使我满盈。($\Rightarrow$注76)

\subsection{研讨会}

\notehead{75$'$}{}

(6月3日)关于研讨会的其他与会者,我得知了一些细节,它们驱散了一切疑问。尽管官方日程里没有安排Mebkhout的任何报告,Verdier却不得不当场、在最后一刻请他做一个报告,以补某场官方报告的缺漏(那场报告曾托付给布日林斯基——他对D-模理论不甚了解)。Mebkhout由此得以陈述其观念与结果,尤其是好上帝定理,其陈述方式使人们对这条定理——以及与之相随的那套哲学——的父系渊源不留任何疑问;而正是它们,使得代数簇上同调壮观地重新启动,并落实为这场研讨会。于是,\textbf{研讨会的全体与会者,经由这场报告,都被告知了这一父系渊源}。我料想,所有人无一例外,此后也都读到了研讨会的会议录,尤其是导言与前引的Beilinson--Bernstein-Deligne文章。显然,没有一个人觉得有什么不对——或者即便觉得,也没有流露半点。Zoghman~Mebkhout在这方面的回音,一无所获。于是,研讨会的全体与会者,都满可以被视为这场在研讨会期间做成的欺瞒的共谋。

这场集体的欺瞒,早在研讨会当时就已分明,因为当Deligne在其关于所谓“反常”层的口头报告里不提Mebkhout的名字时,没有任何人觉得有什么不对。报告人只限于陈述好上帝定理,并说\dmark{}\origpage{290}他不会在报告中证明它。另一方面,他倒也(带着他那惯常的谦虚)把话说得很足:猜出他所称“反常”的层那些卓越的、事先难以预测的性质,“没有任何功劳可言”——那些性质,由他方才谈到的“Riemann--Hilbert对应”明摆着提示出来\fnote{12}{(*)}{可与所引文章的第10、11页对读。(6月7日)关于藏掖之技艺的细节,见下一条注“变戏法的人”,n\textsuperscript{o}~75$''$。}。人人都觉得正常:他避免点名那位有幸发现这条天赐对应的“有功之人”,并摆出一副父权在己的模样——而恰恰是他们刚被告知、或将在随后几天里被告知:事情绝非如此。人们想必认定:一个研讨会上的无名龙套,竟是一条如此卓越定理的作者,这必是某种不可容忍的错配;于是人人尽了一份力,拨正船头,确立一个共识,把父权归于那位分明是天造地设的人——那位\textbf{本应}是作者的人\fnote{13}{(**)}{(6月5日)况且一切都环环相扣!“学生”行列(继“研讨会”行列之后)中继续进行的反思,以及某种声调(尤其还见于最近与Deligne的一次简短的通信往来——见“葬礼”一注(n\textsuperscript{o}~70)的第一条脚注),向我表明:对Deligne、对我的其他上同调学家学生们来说,久已分明:平展上同调的发现与其掌控,那位本该的作者,同样是Deligne;而在某个层面上(即支配行为与态度的层面),他们深深相信:\textbf{归根结底}那是他,而我在旁边,不过扮演一个毛糙而笨拙的助手的角色——与其说有用,不如说妨碍一套理论(通向“Deligne定理——原Weil猜想”)的和谐展开,妨碍一种令所有相关者满意的角色的分配……}。

有一处很有代表性的细节:\textbf{Mebkhout的报告没有出现在研讨会的会议录里}。Verdier曾请Mebkhout不要成文他的报告,理由是:研讨会旨在陈述新结果,而Mebkhout的结果早已发表两年有余。

当人们不肯把自己关在技术话语里,而去看这场光辉研讨会里真正发生了什么——看推动各方的力量与胃口——人们会以为在看一部电影:某座遥远\dmark{}\origpage{291}的巨型都市的暗巢里,黑手党的天下。然而这是一幅地道我们自家的风俗画,演员们是法国与国际科学最高贵的翘楚。那位对一切发号施令的大主厨,不是别人,正是不久前还之于我扮演着谦和微笑的精神之子——或至少合法继承人(同样谦和,同样微笑)的那一位。至于那可差遣、可宰割的“软柿子”——在满是寸步不让的“硬汉”的世界里——由于一个我尚未完全参透其意味的奇怪“偶然”,也与我这个人紧密相连。他是我的“学生”,一如大主厨(而且同他一样,是带引号的“学生”……)——那是在我已被宣告死亡并下葬多年之后,进到我的门下来学艺的……

\subsection{变戏法的人}

\notehead{75$''$}{}

(6月7日)人们尽可以在那篇“值得纪念的文章”(前两条注里谈到的)里,欣赏那套漫不经心的藏掖的炉火纯青。作为全部工作之本质动机的那个范畴等价,头一回被引进,是在导言第4页(第10页,第9--15行)的一句话的转角处,不给它任何名字,随即转入“反常”层观念的一连串推论(第10、11页)。此后直到第16页末,它再未被提起;在那里我们读到\fnote{14}{(*)}{下面的引文里,着重号是我加的。}:

\begin{quote}
“应当指出,在以下几点上——\textbf{它们本应在这份笔记中占有一席之地}——我们未能完成任务。

——反常层与完整模的关系。如本导言所述,它曾起过重要的启发作用。关键陈述是4.1.9(\textbf{此处未证明})……”
\end{quote}

(随即转入别的“本应占有一席之地”的要点。)

我赶紧去看:作者们没肯劳神纳入其工作(或至少没肯纳入其证明)的这个“关键陈述”究竟是什么。找n\textsuperscript{o}~4.1.9……我碰到一条“评注4.1.9”——想必不是它;我找的是一个“关键陈述”,一个正式表述的定理或附注,并带有一个出处,表明作者们\textbf{在}哪里证明过或将证明它——既然他们不\textbf{在}这里证……可我怎么找,也没有“定理4.1.9”的踪迹——只有一个段落编着4.1.9的号。我姑且把这条“评注”读下去(不抱什么信念\dmark{}\origpage{292}——想必是编号出了错……):里面说“4.1.1在复上同调中的类似是真的……”糟糕,难道我还得回到4.1.1去看这说的是什么?我跳过去,接着往下读——瞧,我本已不再指望了:十一行之后,有一句话以“人们知道……”开头,以“诱导范畴……与反常层的范畴之间的一个等价”收尾。

呼——原来就是它!可我再怎么往下找,也找不到只字片语,来澄清这句谜语般的“人们知道……”。那位本不曾“知道”的读者,只觉得自己全然配不上这局面。对他来说,反正有一点是清楚的(且不论他配不配):这个“本应在那份笔记中占有一席之地”的结果——人们在此借一条技术评注的转角处“提醒”读者应当知道的东西——分明归于那“笔记”的作者们,或其中一位;也许是其中最显赫、并执笔了这篇文章的那位(有一种“家传的笔法”是不会错的……),也就是做了那场口头报告、而其人人皆知的谦虚当然不许他说“是我!”的那位——但人人都懂了,无须他说……

这立刻让我想起最近几星期反思里的几段记忆。头一段,是Deligne1968年的第一篇文章:在4月22日的“驱逐”一注(n\textsuperscript{o}~63)里——那是在LN~900谜底揭穿之后三天——我(十六年后)终于费心把它细看了一遍。我在这里又见到同一套笔法,想必经过十三年的“磨合”而略有变体。在那篇1968年的文章里——其主要灵感来自我——他在文章结尾处顺带、谜语式地提了我的名字,权作“手续完备”。在这里,他连这份小心都不再费——经验久已教他:完全用不着了!相反,在他年轻时那篇文章里,既然他感到非提我不可,他就以完全藏掖该文的最初动机作为补偿(连权的瑜伽一起,以便六年后以一个替补的父权把它端出来,再过八年,又等着Motif的开棺出土……)。无论如何,即便藏起(并留作独家收益……)这篇文章本质的算术动机,这篇文章也“站得住”,完全可读,配得上其作者把事情做得尽善尽美\dmark{}\origpage{293}的名声。而在这里,他所发展的理论,若没有那条启发性动机,便不可理解。于是他给出了动机,却以“关键陈述”这个定语来指涉它,同时把它从腿底下打发掉——不赐给它一个名字,不赐给它一个成形而冠以定理或命题之名的陈述;连“对应”(所谓Riemann--Hilbert的)都没有——他把这份差事留给了他的朋友Verdier与Teissier。他无须给它命名(鉴于其分量之轻\fnote{15}{(*)}{(6月14日)为安放这个“轻”,我提请注意:Deligne曾在IHES主持一场研讨班,试图发展可构造离散系数向连续系数的翻译,而未获令人满意的结果。对此可参见“当值无名之辈与好上帝定理”一注,n\textsuperscript{o}~48$'$。}——他肯定五分钟就能证明它!),也无须提名任何人——别人自会代劳,并且完全合他的心意。分明有一部瑜伽、一套哲学,作者以完美的娴熟与权威挥洒自如,而无须命名任何东西——那个他装作不屑的“轻”(“本应在这份笔记中占有一席之地”),他心里清楚:只要他知道在该沉默处沉默、肯等待,它自会外加送到。他头一回成功地玩这一手,那个“轻”是“权方面的若干考虑”,借一条谜语式评注的转角处影射了一下(且等六年后,再把权的哲学大张旗鼓地端出来)。据我所知,第二回是在我1970年离去之时——那个“轻”是“Motif之梦”,它在十二年间都不配得到一言的礼遇(想想吧——不过是个梦,而且还是个亡故者的梦,外加未经发表!),且待这一回发现真正的motifs(以及拿它们能做什么),再把那无可争议的父权——一如既往地谦虚——戴在自己头上\fnotevariant{16}{(***)}{(**)}{关于这套“以轻蔑来据有”的技术的其他评注,见次日之注,n\textsuperscript{o}~59$'$。}。

\subsection{反常性}

\notehead{76}{}

(5月4日)我清楚地记得,头一回听到“反常层”这个名字(译按:法语pervers还有“邪恶”的意思)——那应当是两三年前——它曾令我不适地一震;它在我身上激起一种不安之感。这份感觉,在此后两三次再听到这个不寻常的名字时,又出现了。那是一种内心的“却步”,停留在意识的表层;倘若当时我停下来检视它,它想必会表达为类似这样的话:给一件数学的东西起这样的名字,是什么主意!甚至给\dmark{}\origpage{294}别的任何东西或生灵起这种名字——严格说,唯有对人可以破例——因为明摆着,在宇宙间的一切“物”之中,唯我们人类,是这个字眼有时可以施用的对象……

我相当有把握(不敢说全然有把握):头一回对我谈起所谓“反常”层的,不是别人,正是Deligne本人——那是在吕米尼研讨会之后,他顺路到我家来的时候\fnote{17}{(*)}{倘若确实如此(如今我已深信不疑),那就得归功于我朋友的谦虚:因为(至少在意识的层面上)我并不曾疑心,引入并命名它们的人不是别人,正是他。是读了那篇“值得纪念的文章”,我才明白过来。(5月28日)老实说,文章本身对此也同样不置一词——正如它不置一词说Deligne是Riemann--Hilbert对应的父亲。然而,对“反常层”这一命名的父权,我从未有过怀疑;而且这一父权,此后确实得到了证实。}。那甚至想必是我们之间最后几场数学谈话之一——他到我家来之后,再没有过别的。恰恰是在那次来访时,显现了那个“征兆”;几个星期或几个月之后(当这一征兆在那次会面之后的数学通信往来中得到再次确认时),正是它使我断绝了数学层面上的交流\fnote{18}{(**)}{在纯属个人的层面上,这份关系仍以过去那种亲切友好的音调延续着,表面毫无变化。我的朋友习惯大约每两年,趁某次远足之便,来看我一次。去年夏天他又来过,那是一个可喜的机缘,使我也得以认识他的妻子列娜与他们尚且幼小的女儿娜塔莎。我记得那是从又一场吕米尼研讨会回来的路上;关于那场研讨会,我几乎没有听到什么回音(只有Mebkhout的几句阴郁而含糊的暗示——人们又一次好心地邀请了他,而他也想不出比再度入局更好的事……)。他们在我家住了两三天,接触自始至终都极好。}。(关于这一情节,见“两个转折”一注,n\textsuperscript{o}~66.)

回到那些所谓(我以为是错称的)“反常”层:显然,“正常地”说,这些层本该叫作“Mebkhout的层”,那才算公道。(不止一次,我曾把一些我提炼并研究过的数学观念,冠以先行者或同行的名字——他们与观念的联系,远不及Mebkhout与这个美好观念的联系来得近——何况这观念,依我看,其音色与其说“反常”,不如说“崇高”!)Deligne在发现并命名这个出自Mebkhout著作的观念之际的心境——他自己分明早已“奢侈到无边”,却仍准备巧取他人——这份心境,满可以正当地称之为“反常的”。肯定,我的朋友自己,在其内心深处,在某个不会为所炫示的门面所欺的层面上,也曾感受到这一点。在这个名字的授予里(乍一看,它显得悖谬),我感到一记\textbf{挑衅}之举,一种在如此完全的权力之中的醉意——它甚至可以炫耀(象征性地,借一个挑衅性的名字的陈列——其分明之极的真意,\textbf{没有一个人}肯去读!)它那巧取他人的“反常”本性。

\dmark{}\origpage{295}在我看来,绝非不可能:在某个深层,我曾知觉到我朋友身上这份心境的音色,而这曾促成我方才谈到的那份不安\fnote{19}{(*)}{我甚至倾向于认为确实如此。我不止一次在我身上确认:对事情的深层知觉,其精微与敏锐,与掠过意识层面、浮于意识表层的东西,全然不可同日而语。那个完全“醒着”的人,想必正是这些知觉被持续地整合进有意识的视野与有意识的经历之中的人——也就是那个完全按其真实的手段生活、而不只是靠这些手段中微不足道的一小部分生活的人。}。这份不安,尤其表现为:对他想必曾向我作的种种解释,我一概心不在焉;而我不相信,在这次会面之前,曾有哪次我没有持续专注地听他说话——尤其当事关数学的时候。我心中对这个被(天知道为什么)称作“反常”的观念,有某种阻塞——我并不真想听人谈它,尽管它与我一直(且在某种程度上依然)十分贴近的一些问题紧密相连。

干脆说吧:Deligne等人的这整篇文章,就是典型的、如假包换的“\textbf{Grothendieck学问}”,它们满可以同样出自我手(唯独主要观念的名字除外)!这大体上就是我在前一条注(n\textsuperscript{o}~75)第二部分里已表达过的;也是我通读所引文章的那一刻就已感到的——只是那份弥散的感觉,还不曾落实为这一触目的确认——它是我刚刚才做的。这份确认,以一种惊人的方式,使我重新对这桩深层的矛盾敏感起来:那个(在某种意义上)不能自已地复制着他所否认的人、并同化于他的人——那个他要埋葬的人,同时\dmark{}\origpage{296}也正是他\textbf{想成为}的人,而且(在某种意义上)他\textbf{就是}的人。

前天,在写前一条注(“不义——或归来的意义”)时,我已被这一巧合打动:我与朋友之间关系的这一转折——它骤然失去了一段共同热情中的契合,而那契合曾是这段关系之所以存在、并最强韧的内驱——恰逢我的朋友从那场值得纪念的研讨会归来之时;而那场研讨会的意味,刚刚向我显露。1981年7月我们相会时——在某个层面上,它与我们会面的其他场合同样友好、同样亲切——令我瞠目的,是那个“征兆”:论声调与神色,隐约,却又粗暴地分明,一份蓄意的轻蔑。它像是我的朋友预先支取的一笔\textbf{定金}——这一回,是在个人关系的层面上——对着那份隐含的、同样“隐约”(而且同样“粗暴地分明”)的轻蔑:那是他刚刚在吕米尼研讨会上、在当时一场当红名家间的技术精湛的盛大演出之中,公开对作为公众人物的我表达的。也是同一份“轻蔑”,刚刚对着那个胆敢(哪怕只有一点点)自称源出我门下的人表达出来(只是这一回,带着一种全然不同、“反常”得多的粗暴);而那人因此注定在我朋友皮埃尔眼里——至少在某个层面上——只成了“另一个Grothendieck”\fnote{20}{(*)}{在我们的个人关系里,我的朋友以我名字亚历山大的爱称(源于俄语)称呼我;这也是自我童年起,我的亲人、最亲近的朋友们对我的称呼。}——而如今,当务之急是不惜一切代价碾碎他……

\subsection{暂停!(Pouce!)}

\notehead{77}{}

(5月5日)翻阅这篇值得纪念的文章\fnote{21}{(**)}{关于“值得纪念的文章”,见注n\textsuperscript{o}~75。}——它(据说)主宰了1981年6月那场同样值得纪念的吕米尼研讨会——另一个细节打动了我。最后一章,顶着“从$\mathbb{F}$到$\mathbb{C}$”这个意味深长的名目,连篇累牍地描述一条卓越的原理——我把它引入代数几何,恐怕已有二十年了;那应当还在Motif观念诞生之前(Motif经由原Weil猜想,给出它最\dmark{}\origpage{297}深刻的例证)。这条原理断言:对涉及一域上有限型概形的某些类型的命题,只需在有限基底域上(亦即在一种“算术性质”的局面里)证明它们,即可推出其在任意域上、尤其在复数域上的有效性——在后一情形,所论的代数-几何结果有时可以经由超越途径改述(比如以整系数或有理系数上同调的语言,或以Hodge结构的语言,等等)\fnote{22}{(*)}{(5月6日)我记得,运用这样一条原理的第一个例子,见于Lazard关于仿射空间$\underline{E}$上代数群律的幂零性的定理(在任意域上)。它的证明曾深深打动我;我由它取得启发,处理了许多别的命题,并把它做成一套“哲学”,它主宰了我关于Motif理论的反思。}。我的朋友知道它,不是从别人、恰恰是从我嘴里,经年累月,就着许多例子\fnote{23}{(**)}{其中一个例子,见“驱逐”一注(n\textsuperscript{o}~63)。}。这条原理的父权(其初等形式甚至在EGA~IV里有明文——请别问我哪一节哪一条……)更是众所周知的\fnote{24}{(***)}{(6月5日)我自称这条原理的“父亲”,或许言过其实:据我所知,它的第一次应用当归Lazard(见前一条脚注)。一如在别的场合,我的角色,是感到他人一个想法的一般性,并把它系统化到使之成为一种“反射”或“第二天性”的地步。在权与Motif瑜伽的框架里,头一个使用这条原理的,很可能是Serre(而非我),凭他虚拟Betti数的那个想法——正是它,把我引上一部一般的权与Motif瑜伽之路。(Serre的那个想法,见注n\textsuperscript{o}~46\textsuperscript{9}。)诚然,流行的惯例是:一条已成通例的推理“原理”,其父权不归于可找到其最初踪迹的作者,而归于头一个知觉到其一般分量、把它系统化并推广的人。在这个意义上,可以说 N.~Katz的更正(下面这句话里谈到)把这条原理的父权归于我,是有道理的。}。以至于1978年赫尔辛基大会上为我这位光彩夺目的朋友授予菲尔兹奖章时, N.~Katz禁不住在颂扬 P.~Deligne的演说里顺带提了它一句,从而(不动声色地)更正了其显赫的受奖人那有点难堪的系统性“遗忘”。我是几天前才读到这篇演说的,与那篇“值得纪念的文章”同时。

无论如何,在这篇文章里,“算术”向“几何”过渡的哲学,被以这样的措辞呈献:任何一个不知情的读者都不可能怀疑,那位光彩夺目的主笔作者(请原谅\dmark{}\origpage{298}我的失言……)刚刚才发现这条分量如此之大的奇妙原理。

诚然,我没有为这条方法申请专利,而我光彩夺目的朋友也从未在任何地方说他是它那神奇的发明者;正如他并未明言他是那条著名的“对应”的父亲(请欣赏这个词,它散发着好一股十九世纪的馨香!)——那对应被谦虚地记在Riemann与Hilbert名下(这两位,倒也配得上给一位如此显赫的后继者的孩子们当教父);正如他在那部“值得纪念的文集”(LN~900,原书此处作LH~900)里也未指明:发明Motif、MotifGalois群以及与之相随的一整套哲学(他至今才掏出其中一段)的,分明是他;SGA~4½那部名著也同样无话可说——在那里,人们甚至再度不吝赐我以该卷“合作者”的荣衔:该卷从零开始如此光辉地发展平展上同调,还肯赏脸(尽管那两卷带着可惜的“多余细节的矿渣”等等)求助于SGA~4与SGA~5这两卷注定被遗忘的卫星卷,并且宽宏大量地承认后者的功劳:提供若干补充与技术枝蔓(其中有些甚至“非常有趣”)\fnote{25}{(*)}{关于“SGA~AT之役”(原书如此)的细节,见“白板”“格格不入者”“绿灯”“颠倒”四注(n\textsuperscript{o}~67、67$'$、68、68$'$)。}。

在所有这些情形里——也在最近五六年里我得以确认的许多别的微型个案里,尽管我从没想到过要\textbf{辨明我的不安}、给我所目睹或参与的事情起个名字\fnote{26}{(**)}{在一个个案里“辨明我的不安”的第一步,是不到三个月前在《收获与播种》里迈出的:那场(事后证明十分费力的——而且不无原因!)反思“注记,或新伦理”(第33节)。这场反思在那则注“年轻人的势利,或纯洁的捍卫者”(注n\textsuperscript{o}~27)里得到继续;然后,不到两周之前(在前一天发现“值得纪念的文集”(LN~900)所带来的冲击之下),又随注n\textsuperscript{o}~59“新伦理(2)——或抢夺的集市”重新开始。写那一注时,我心中对动用“抢夺的集市”这个相当粗鲁的字眼,还留有一丝踌躇。此后接连而来的发现向我表明:任何踌躇都不合时宜。}——在所有这些情形里,我认出同一套\textbf{做派}。我的朋友永远、彻头彻尾地处于“\textbf{暂停}”的豁免状态——他可以从容取用,带着那份由同侪与下侪的仰慕(无论如何是有根有据的)所给予的坦然良心;那仰慕,担保着完全的豁免。

\subsection{中国皇帝的新衣}

\notehead{77$'$}{}

\dmark{}\origpage{299}(5月7日)当然,那些看着我朋友Deligne行事、并且多少“在局中”的人——我是说那些并非刚下船、并非刚刚从当事者本人或其他当红名家(即便不都金光闪闪)的出版物里才学到“时下数学”的人,亦即那些了解来龙去脉的人——这些同行(毕竟,他们也还没有那么稀少!),在\textbf{某个层面上},是明白发生了什么的。在那几桩“有点扎眼”的个案里,他们必定感到了那份特殊的小小不安——那是我自己在这些“微型个案”面前、比起他们的个案微小百倍的情形面前,也不止一次感到过的。但他们所感到的是如此\textbf{庞大}、如此\textbf{不可思议},以至于它想必从未浮出水面——一如它最终开始在我身上浮出水面那样,经由一项\textbf{工作}:它表达为围绕着一个微型个案的两份文本(前一条脚注里谈到的那桩个案)。我确实没有听说,在我们的科学或任何别的科学的历史里,有过堪与比肩的事。对某些人来说,它不曾“浮出水面”,反而想必\textbf{成了风气},或至少被视为\textbf{正常}——只要一位分明是天才的、为万人所仰慕的人物,以天底下最自然的神色实践它,在众目睽睽之下,而这事从来不曾(据我所知)引起过只言片语的评注。

最近这几天里,我禁不住多次重新想起那个故事:“中国皇帝的新衣”。那位皇帝,受了无耻骗子的欺、也受了自己虚荣的欺,传谕说他将举行盛大的出巡,身穿举世未见的最华美的衣裳——那衣裳,是那些号称裁缝艺术家的骗子们斥巨资为他赶制的。而当他出巡,满朝臣子盛装华彩、前呼后拥,那些“艺术家”们打躬作揖,皇室倾巢而出——无论在仪仗里,还是在聚拢来瞻仰这第七大奇迹的人群里,竟没有一个人敢相信自己眼睛的见证;人人恪尽职守地赞叹,并为这身衣裳那无与伦比的华彩锦上添花——皇帝就披着它!直到一个混在人群里走迷了路的小孩失声喊道:“可是皇帝什么也没穿呀!”——刹那间,所有的人仿佛同一个声音,跟着这孩子喊道:“皇帝是赤身裸体的!”

我觉得自己就像那个孩子:他相信自己眼睛的见证,哪怕他所看见的,是如此闻所未闻、前所未见、为众人忽视并否认的东西。

\dmark{}\origpage{300}至于孩子的声音是否足以使某些人回到他们健全官能的谦卑见证,那是另一回事。故事就是故事:它向我们诉说着实在的某样东西——但它并不是实在\fnote{27}{(*)}{(6月14日)写完这条注之后,“中国皇帝的新衣”这个名字,向我显出是这场葬礼的一个天然的副题,表达它一个格外触目的方面。随后,反思转而面向我的全体学生,乃至数学界建制的“全体教众”,这一副题显得不再那么切合。然而,我终于明白:那则首先在我想到我的朋友Deligne时来到心头的寓言,同样适用于葬礼的各个侧面与各个插曲——它们每走一步,都在不可思议之中抵达乌布王式的荒诞(那不可思议,人人都恪尽职守地装作不见),而它却是真的。这个方向上的反思,尤其可参见“进步是挡不住的!”“研讨会”“受害者——或两种沉默”“玩笑——或权复形”“欺瞒”“掘墓人——或全体教众”各注(n\textsuperscript{o}~50、75$'$、83、85$'$、97)——没有一条特别关涉我的朋友皮埃尔。}。

\subsection{墓外相会}

\notehead{78}{}

(5月6日)才过去五天,我终于得到了那份慷慨的一包文献——来自我的朋友Zoghman~Mebkhout;其中尤其是那两份已被检视过的“值得纪念的研讨会”文本——一场围绕着一桩纪念碑式的\textbf{欺瞒}搭起来的研讨会!“不义——或归来的意义”那条注——我在其中尽力消化这桩新“事件”那令人难以置信的意味——就写在我收到这包文献的当天(5月1日的次日),犹带着发现的余温\fnote{28}{(**)}{连同“注记——或新伦理(1)”一节,这条注是仅有的不得不重写数遍的两篇注或节之一:因为头一稿(乃至第二稿)所“流出”的,还带着我一贯看待事物的方式所带有的全部惰性,远在有待检视的实在之下。}。

自4月19日——我终于得知Lecture~Notes那部“值得纪念的文集”(LN~900——见注51、52)——以来,这是关于宏大葬礼之种种盛典的第三项大发现;它在我看来分量也最大:既因它照亮了一些与我关系极近之人的作为,也因其作为一个时代(显然独一无二的时代——诚然,我对历史所知寥寥……)的“风俗画”所具有的含义。

第二项发现紧随第一项而来:即“Motif”的开棺\dmark{}\origpage{301}出土——它们被埋了整整十二年。继“值得纪念的文集”之后,轮到“值得纪念的研讨班”——那场从未举行过的“研讨班”,顶着假名字(无论SGA还是4½这个号),还附赠一篇幽灵论文的“第0状态”,更不必说(真正的)SGA~5研讨班的一场中心报告(它以晚出的面目出现,尽管比它早十二年);一场为战役的需要而被“借走”、不经任何手续的报告。这场光辉的战役,以及它在降临到可怜的SGA~5研讨班(从头、从尾、从中间被拆散!)身上的那些离奇坎坷中所扮演的角色,是在4月24日至30日间持续的一场反思中渐渐显露的。(对此可参见“同伙”“白板”“格格不入者”“信号”“颠倒”五注,n\textsuperscript{o}~63$'''$、67、67$'$、68、68$'$。)

那项发现刚消化完——与此同时,我“我的朋友皮埃尔”的回顾反思也接近尾声——而4月30日,我刚以“尾声——或全体一致的赞同”这条名字双重欢腾的“终注”,在这场无休无止的葬礼下面,骄傲地画上最终的句点(这回是肯定的——我总算走到了头!)——就在这时,这包倒霉的文献到了,把句点、尾声、版式、编号……统统推翻。对文献以及随附的评注与书信匆匆一瞥,便明明白白:我的句点完蛋了,那场头等葬礼的堂皇部署——我正准备修饰其最后细节——也完蛋了;我只好重新套上司仪的挽具……

不过天知道,我的朋友Zoghman有的是时间把局势告诉我!局势以潜伏的形态,想必已持续十年;以“急性形态”(这还是委婉语)至少三年——自那场研讨会起。在会上,他分明感到了风向,而无须等到次年、以其显赫的前老板兼庇护人为教父的那份高度官方的会议录问世。

他的论文答辩(1979年2月)之后几个月,他曾到过我住了六年的那个村子,想送我一份论文。不巧,几天前我刚离开(从此没有回去过,除路过之外……),退居孤独。他只遇到我的女儿,论文是她后来转交给我的。是第二年,我想\dmark{}\origpage{302},我们才终于正式相识,在蒙彼利埃大学;我们聊了一两个钟头。那时我对数学并不怎么上心,想必再不会记得一篇我只花了几分钟翻过的论文,也记不得其作者的名字。但这并没有妨碍接触的热忱。我清楚地记得一股即刻的相互好感。数学谈得不多(据我所记),谈的更多是或多或少的个人事情。Zoghman后来告诉我(此事我已忘却):他还是向我大致解释了一点D-模的“哲学”,而且他对这次会面很满意,感到我在从他那里获悉新事物时哪怕微微“震动”了一下,虽然这些新事物(在某种意义上)又是“意料之中”的。我印象最深的,是他的为人给我的印象——一种执拗而平静的力量之感,一个“闯将”的力量之感。在那个时刻,比起去年我们相会时、或其后的通信中,我更强烈地感到性情的深切投合——尤其在“闯将”这一点上。但两次会面之间的两三年,似乎已颇损耗了他……

我不记得在我们第一次短暂的会面时,Zoghman曾向我谈起他所工作的孤立,谈起那些曾是我的学生的“权威们”没有给他任何鼓励。倘若他有所暗示,想必不曾坚持。在当时,这件事本没有什么可让我惊讶的\fnote{29}{(*)}{(5月30日)这并不完全真——我把后来产生的失望情绪投射到了过去。我还记得,去年夏天与Zoghman相会时,我曾惊讶:我的上同调学家学生们(尤其Deligne、Verdier、Berthelot、Illusie)竟没有一个在Zoghman的工作中扶持他。这份惊讶,十来天后Deligne来我家时又一度重演(我想必向他提了Zoghman一句,没有遇到回音),后来又经与Illusie的一次电话再次重演。(对此可参见“欺瞒”一注,n\textsuperscript{o}~85$'$。)}。我说不好那是在1981年6月吕米尼研讨会之前还是之后\fnote{30}{(**)}{(6月3日)那是在之前——1980年2月,他论文答辩后一年。}。倘若在之后,他心里想必还压着滚烫的事情——而他全然没有那副模样。他更像一个知道自己想做什么、知道自己要什么的人,安静地走自己的路,不寻衅,也无人寻他的衅。

\dmark{}\origpage{303}后来我们没有继续通信。但我清楚地记得他;去年年初,我给他去了个短信,姑且一问:他是否有空着手去做一项为“温和拓扑”而设的宏伟奠基工作——它(依我看)只等着一个有他这般才具的人来做。Zoghman起初没有向我明说:他对这一前景并不真感兴趣;不过他似乎很高兴借这个机会再见一面。那时我离局太深,没有看清局势:我以为D-模理论已然功成、就此收官,一如凝聚对偶理论那样(注78$_{1}$),而Mebkhout也许正缺“大任务”。直到去年夏天我们会面,我才明白:就在他亲手开创的理论里,“大任务”并不缺——有些甚至还未动工,因为连被人看见都还没有!

无论如何,那是一个现成的机会,来第二次会面——这一回不是像第一次那样旋风式的。去年夏天,6月里,Zoghman在我家大概住了一个星期。在数学层面上,这次会面主要使我得以多少弄懂D-模的瑜伽。我“解冻”得很慢:与我旧日的上同调之恋已有些失联,而且尤其一头扎在《追寻栈》的写作里——它处在相当不同的音区。Zoghman看到我一面心不在焉地听,并不气馁;他不倦地卷土重来,带着动人的耐心。我想,当我终于明白那些大名鼎鼎的D-模不是别的,正是我很久以前所称的\textbf{模的晶体}、并因此于奇异空间上仍保有意义的那一刻,我被触发了。于是,我看见我那晶体-微分过去的整片直觉之网从被遗忘的深处浮起,我那“六种运算”过去有点生锈的反射,重新咬合……

于是轮到Zoghman有点跟不上;或者不如说,是事后他才决定,不肯把手指伸进那副齿轮(一如我的朋友皮埃尔不肯——尽管只要我还在场,他曾满腔热火……)。($\Rightarrow$注78$'$)

\notehead{78$_{1}$}{}

\dmark{}\origpage{304}然而,凝聚对偶确有一定数量的“精细”结果——尤其关于“对偶化微分模”的结构、它们与“朴素”微分模的关系、以及非光滑平坦情形的迹映射与留数映射——是我在五十年代末发展的,而且据我所知从未发表。尽管如此,就其实质而言,凝聚对偶理论(至少在概形的框架里),一如平展对偶理论(及其对局部紧空间的离散系数上同调的变体——由Verdier依平展的模型发展)、又如线性代数或一般拓扑,都显现为实质上\textbf{已完成}\fnote{31}{(*)}{(6月12日)就平展对偶而言,只要纯粹性猜想与“双对偶定理”尚未在全部一般性下得到证明,这就不完全真。}的理论,因而其本性是打磨停当、随时可用的\textbf{工具},而不是多少不为人知、有待探究与消化的\textbf{实体}。

\subsection{受害者——或两种沉默}

\notehead{78$'$}{}

我们的会面,是在一种友好信赖与亲切的气氛里进行的。然而这份气氛并未兑现其许诺。我如今明白:从那一刻起,我朋友心中的信赖就远非完全。那是那场大名鼎鼎的研讨会之后两年、会议录在Astérisque里刊出之后一年\fnote{32}{(**)}{(10月9日)Zoghman提醒我:这些“会议录”实际上迟至1984年初才问世。}——也就是他正付出一场骇人听闻的巧取的代价的时刻。可他直到仅仅四天前,才肯把这件事告诉我!去年他来的时候,是从又一场吕米尼研讨会\fnote{33}{(***)}{(5月7日)这里的记忆有一点点混乱——我倒觉得,他是正准备去参加那场研讨会。当然,从那时起,我所忆起的那些“辛酸而含糊的字眼”,缘由并不缺。但这辛酸,又因他在我家小住之后途经吕米尼而再度加剧。他从吕米尼回来后给我打了个电话,我由此得了些回声。从那一刻起,我就有一种十分清晰的直觉:他赶到吕米尼去,是为了享受被“那些人”(我没有细问是哪些人)怠慢的滋味——那些人慷慨地邀请他,为的是享受把他当无足轻重之辈来对待的乐趣。我想必把这话对他说了,或至少暗示了;这想必无助于改善当时我的朋友对我的心绪。}\dmark{}\origpage{305}回来的(这回干脆就是以D-模为主题的)——人们又一次慷慨地邀请了他,而他也赶紧应召而去。他谈起它,措辞既辛酸又含糊,话里带着意思:如今他“火中取栗”已毕,“一切都是别人做的”。我完全可以想象那幅图景——尤其Verdier,忽然想起三角范畴(连导出范畴也算上,反正一不做二不休!)的父权——那父权他搁置了十来十五年,勉强容忍他的“学生”Mebkhout在工作中使用它们……(注81)。

尽管他当时不肯把事情说清,Zoghman心里对Verdier想必积怨已深——这完全可以理解,看看他前老板那副不甚鼓励的样子。然而,我的其他上同调学家学生——Deligne、Berthelot、Illusie——也同样不曾屈尊对他所做的事情发生兴趣、多少扶持他一把。但人们几乎要说:对Zoghman而言,这本是理所应当——(人们会说)他从不曾在长辈们中间体验过别的态度。倘若他当时怨谁,那唯一而且不折不扣的,是Verdier。

从Zoghman的暗示(他分明无意细说)里,我懂了:“人们”在系统地把他的工作轻描淡写——一笔带过而已。这毕竟是人世间最寻常的事。对一件事的重要性的评判,在很大程度上是主观的;给自己的工作、给自己伙伴与盟友的工作,比对别人的工作——尤其对那些由于这种或那种缘故而人们有意贬低者的工作——记上更多的功劳与重要性,是寻常而几乎普遍的事。(而在这里,那个“缘故”对我而言实在算不得什么谜!)没有什么能使我疑心:远在这些寻常态度之外,这里有一场纯粹的巧取战役——其中根本谈不上“轻描淡写”,而分明是要把Mebkhout对那些使停滞之地重获生命的观念与结果的父系渊源\textbf{整个藏掖}起来……

然而,倘若天底下有一个人,我的朋友向他倾诉是再自然不过的,那正是我:在那几年执拗的工作里——有时辛酸,逆着时下的风尚——是我的著作给了他灵感;是我,在家里亲切地接待他,并且在尽力学习他乐于教给我的东西之时,轮到自己多少\dmark{}\origpage{306}做了他的学生\fnote{34}{(*)}{而且,Zoghman不但不对我谈他自己的下葬,也不对我谈我的下葬——尽管他置身头排,注视那场下葬的进行,已然将近十年!老实说,他的那些“庇护人”(边缘上还带着几分勉强)甚至曾让他亲手抬过装着我的遗骸的棺材的一角;但他们没有原谅他:他是满座宾客中唯一敢偶尔说出那个众人皆讳莫如深的名字的人!于是,我的朋友想必在与我的关系里自觉处境尴尬,而他未能在自己身上找到那份单纯,去承担一段(像我的一样)满载暧昧的过去,并不绕弯、清清楚楚地说话。谈他的下葬,也就谈我的下葬、谈他本人在其中扮演的角色……无论如何,我最终发现这场大名鼎鼎的葬礼的全部盛况,是顶着一种“沉默的共谋”才做到的;这一共谋既包括我的朋友Zoghman,也包括我的朋友皮埃尔,想必还包括我在数学“上流社会”里的大多数朋友。(6月3日)关于更多的细节,见随后的注n\textsuperscript{o}~78$''$。}。

我朋友怀着亲切的热忱来访之后,紧接着就有了一次“反弹”。我有这样的印象:他决意把八到十年间——在他那些曾是我学生的某些人所给的冷漠与轻蔑的刺痛下——积攒起来的疑忌与辛酸,移栽到我的身上。在其后的几个月里,我们之间的通信始终没有离开酸中带甜的音域——它最终停在一张贺年卡上,没有得到回音。

直到3月底,我才重新联络Zoghman,给他寄去“一段过去的重量”以及我当时为该节补加的各注(n\textsuperscript{o}~45、46、47、50)。我是想问他:是否同意我像已经做的那样,让他出现在关于我著作的那段简短反思里(在“我的孤儿们”一注,即n\textsuperscript{o}~46中)——尽管对人人都会一清二楚:我使用的是他给我的信息,而他或许认为那些信息是保密的。我全然没有把握:我的朋友会不会(如在他之前的别人那样)宁可“把自己压扁,也不愿招人嫌”。倘若真是那样,我会难过。

等他的回信,我觉得日子漫长——十天后才收到。我多少料到它还会是模棱两可的——但\dmark{}\origpage{307}这一回,它分明是热忱的。他毫无保留地表示同意,甚至为我在其中谈到他的那些措辞而动容。

在他那封长信(八页)的第6页,他才顺带地、就其定理那“令人印象深刻的”应用之众多(“无论在平展拓扑的框架还是超越的框架里”)提了一句:这条定理在文献里始终顶着“Riemann--Hilbert对应”之名\fnote{35}{(*)}{其信件的引文,见“一种不义与无力的感觉”一注,n\textsuperscript{o}~44$''$。}。他说得那么附带,而且那一手字迹潦草得仿佛故意,险些整个从我眼皮底下溜走!但我到底记住了;那真是一件怪事。怪到几乎难以置信;再说,也许我的朋友言过其实,分明他对人人都怀恨在心,包括我在内——尽管我心无恶意,这一点本是够清楚的。于是我补写了一条注(好个Zoghman,我还以为已经收尾了呢!),名之为“当值无名之辈与好上帝定理”,外加另外两条:“本能与时尚——或强者的法则”(写它时,在别人之中我也想到了他)与“罐装的权与十二年的秘密”。这条“当值无名之辈”的注,我起初写得并无完全的信念;Zoghman在我眼里纠结之至、充满矛盾,以致我自问:我这是要驶向哪里——我只做他的回声,对事实本身并没有多少第一手了解。我压根没有想到可能有巧取,更没有想到Verdier或Deligne本人牵连在内。Zoghman对我说过的话里,没有任何东西能提示这个……

然而,这两人都与那条好上帝定理系得太紧:不经过他们至少心照不宣的同意,其父系渊源几乎无从藏掖。这事想必在其后的几天里在我身上发酵。我想起Deligne曾对那个(十年后)由Zoghman解决的问题苦苦思索过;至于Verdier,他毕竟担任过研究导师;即便他没为学生费多少力气、宁可用冷淡与气馁来打发他,他至少应当知道这项工作里两条主要定理是什么——Zoghman肯定向他解释过,在那几场大名鼎鼎的、Verdier\dmark{}\origpage{308}“恩准过”的“会谈”里!于是我给这条注增补了一则评语,论Mebkhout的工作与Deligne一次较早尝试的关系,外加一条论Verdier角色的脚注。这同时也是对我朋友Zoghman的一次试探……

人们也许会想:这样一来,Zoghman必会抓住时机,亮出他藏了三年的炮群,让分明的真相轰然炸响,让受压迫者的大业凯旋!全然不是!十五天的沉默之后,来了一封信,谈数学里的一切,唯独不谈好上帝定理——或者不如说,关于它,他只限于按我的请求,给我其论文中的精确出处。(我到底想知道,这条我如此斩钉截铁地担保过的大名鼎鼎的定理,在哪里得到证明!)

我不得不在回信里,就“对我的著作的宏大巧取”——我刚发现的那桩(随“值得纪念的文集”LN~900,外加“预支了许多乐趣”:不日将在系图书馆里结识SGA~4½)——说上几句,才在又一个十天的沉默之后,使我的朋友终于引爆!

这一回他总算“倾囊而出”——而且是一\textbf{大}包:挑选得当的文献,使我(一个不常光顾图书馆、甚至连我在大学办公室里越堆越高的抽印本也不常翻的人……)得以对一种“气候”形成均衡的印象;在那气候里,许许多多的人,至今仍未参与我那漫长而庄严的葬礼\fnote{36}{(*)}{(6月12日)不过,Katz、Manin、Langlands似乎不在其列……(1985年3月)关于Katz的另一种说法,然而可参见“把一个个 i 点清”一注(n\textsuperscript{o}~164(II5))与“调动”(n\textsuperscript{o}~169)“插曲2”。(1985年4月)Langlands亦然,见“预先出土(2)”一注,n\textsuperscript{o}~175$'_{1}$。}。除了主要的“铁证”(那场大名鼎鼎研讨会的两篇文章,它们使那桩不可思议的欺瞒原形毕露)与另一篇“值得纪念的文章”(这一回出自Verdier的手笔\fnote{37}{(**)}{关于这篇文章,见“好的参考文献”一注,n\textsuperscript{o}~82。})之外,还有 N.~Katz论“菲尔兹奖得主”Deligne的演说,外加Langlands与Manin在同一届1978年赫尔辛基大会上的报告;然后是Deligne1970年尼斯大会上的“Hodge理论I”(其中第3行又一次影射“Grothendieck的猜想性Motif理论”(注78$'_{1}$))\dmark{}\origpage{309},以及同一位Deligne1974年温哥华大会上的“代数簇上同调中的权”(其中我的名字未被提及(注78$'_{2}$));最后还有与 A.~Borel(又一位老哥们——我同时得知他已回苏黎世……)的通信,以及Mebkhout在法国科学院报告上的两条短文,其中1980年那篇是其学位论文第五章(前一年答辩)的摘要,把好上帝定理更突出了一点\fnote{38}{(*)}{这条注的精确出处——Mebkhout的论文与好上帝定理——见“大部头与上流社会(或:鱼目与珍珠……)”一注,n\textsuperscript{o}~80。}。还不算一份文件——嘘!——以保密之名传达,我在这里一个字也不多说……

两封信随这份厚实的包裹而来(4月27日与29日),一封很长,两封都很有分量。如今他终于把“底”(这一回是真的!)亮出来了,Zoghman却仍旧像自我重新联络他以来那样,力劝我极端谨慎。倘若听他的,我会小心不把我的反思笔记公开,它们将只在我们二人之间保持绝对的秘密——至少牵连任何人的那一部分,因为“他们”“大权在握”,“人人都站在他们一边”\fnote{39}{(**)}{(5月30日)我兴之所至,在这里说得有点过头。Zoghman从未在任何时刻建议我不要发表我笔记的这一部分或那一部分。近来他甚至坚持:这些笔记理应以书的形式问世,以造福“后世”——一份限量印制、类似预印本的东西,在他看来有点“竹篮打水”。}!然而,我分明告诫过Zoghman:这些笔记——我寄给他的只是与他相关的摘录——注定要公之于众,而且越快越好。

一切要素似乎终于齐备,足以使受压迫者的公义大业凯旋;但那位“受害者”似乎仍在竭尽全力,存心把水搅浑——仿佛带着一种(人们会说的)事后的懊悔,后悔亮出了那道大名鼎鼎的“底”——直到命中注定的5月2日为止,Zoghman想必是它唯一的持有人。这份暧昧,几乎每一行都透出来(我几乎没有夸张),直到我刚收到的最后几封信——包括最后那封:他以一种阴郁的凯旋神色,把“值得纪念的文章”全份寄来(而先前那“大包”里,他到底只舍得与这份关键铁证的头二十页分手\fnote{40}{(***)}{(10月9日)Zoghman向我说明:其实他手头起初并没有全文的复印件,是后来才复制的。}。)

\dmark{}\origpage{310}至于皮埃尔朋友——我是说Deligne(对世人而言,他既非皮埃尔,也非“朋友”……)——他简直恨不能含着热泪唱起赞歌——这样一来,倒仿佛不再是Zoghman是“受害者”,而是可怜的Deligne:他被周围的人带坏了,毒害至深;唯一的恶人、把他带坏的人,是Verdier(而且……不妨顺着我的目光看……):分明是我“一定得罪过Verdier”,以致他恶意到只为害人而害人;更不必说,做他导师的也是我,授予他博士头衔、荣耀以及其他一切的也是我——总而言之,“绝对权力”的手段全在我!\fnote{41}{(*)}{这已不是我头一回听到“绝对权力”这样的调门,人们想借此让自己相信自己的无能,并为之辩护。倘若有人曾把对自己本人的“绝对权力”托付给什么人,那么对Zoghman而言,那不是别人,正是Zoghman自己!}

显然,倘若我的朋友怨谁,他真正怨的,并不是他那位显赫的前老板——那个人,他十年里有幸见面“会谈”总共不过三次(照他最近写信对我说的来理解)——一个遥远得令人眩晕、全然够不着的人;他怨的,是那个他想来就来、可以同吃同住的人……\fnote{42}{(**)}{(5月8日)我朋友对我的关系中,冲突的毫不含糊的征兆,恰恰出现在那次小住之后——那次他“分享了我的面包与住处”,气氛毫无保留地亲切,消除了我们第一次短暂会面想必不曾完全抹去的“距离”感——这肯定不是偶然。我在这里遇到一个我由来已熟悉的局面;我曾在“敌父(1)、(2)”两注(第29、30节)里(以相对一般的措辞)表达过它。当年我写它们,作为对先前反思的评语时,我不曾料到:我在其中描绘的那个原型局面,竟会始终处于一场尚在到来的漫长反思的中心——而那时,我还自以为已近旅程的尽头!}

Zoghman每前进一步,披露一点新东西,使我更多地了解他身为受害者的那场巧取局面(并能多少有助于解开它),我都感到那像一次\textbf{挣脱}——一场筋疲力尽的内心斗争的终点。有一个\textbf{角色},他似乎已以身家性命相认同,死死攀住,如同攀住他最珍贵的财产——那\textbf{受害者}的角色\dmark{}\origpage{311};而要保持这个角色,他就必须围绕这个角色以及为之辩护的局面,保持最绝对的秘密\fnote{43}{(*)}{(5月30日)自这些文字写成(5月6日)以来,我朋友的态度已发生急遽的演变;近来我不再觉察到他对受害者角色的攀附。当然,下面的文字(一如上面的文字)所关涉的,是我朋友生活中的某些情节,绝不声称刻画一种性情、或描绘一种恒常的成见。}。而此刻,在他不情愿的配合之下——那配合仿佛是被一个局面的逻辑硬挤出来的;而那个局面,不是别人,正是我带着这些多事的、关于一场风平浪静的葬礼的反思造成的——这个秘密行将终结,连同也许还有那个他不知从何时起便乐于固守的角色;在这样的时刻,他确实可能痛苦欲裂,比以往任何时候都更怨我。

我朋友Zoghman的这场“下葬”,是由\textbf{两种沉默}的合力做成的:每一种回应着另一种、又反过来挑衅它,在一轮毫无罅隙的回旋里,一方的角色紧紧贴合一方的角色——巧取者们与被巧取者。倘若我不止一次被这样一个景象攫住:“下葬者”同时、而且更深地,是他自己的“被葬者”,那我同样被另一幅景象攫住:在另一位朋友的身上,一个“被葬者”同时、而且更深地,是他自己的“下葬者”——与那些他乐于为之充当自愿受害者的人,紧密共谋。而且我看得很清楚:他自己那场巧取的第一责任人,不是别人,正是我的朋友Zoghman本人——三年来,他经由沉默,默许那些对他恣意妄为的人对他的羞辱。他手里握有一切,可以一战;他却三年来选择忘记自己还有手,未曾抵抗便告败北\fnote{44}{(**)}{(5月30日)这当然是一个有斗争家气质的人的主观看法;那种气质,在这位身上似乎阙如。不过,看来这些文字写成之后,战斗的心弦在我朋友身上已经苏醒,他已决意与那桩让他付出代价的不义一战。(1985年4月18日)关于对我朋友心境的一种不同、也不那么“苛刻”的照亮,另可参见“根”一注(n\textsuperscript{o}~171$_{3}$)。}

\notehead{78$'_{1}$}{}

这份简短的预注,我从不曾亲手捧读过;我读到的只是更详尽的“Hodge理论II、III”——发表于《Publications Mathématiques》。正因如此,我一直有这样的印象:Deligne不曾认为有必要哪怕暗示一下,他关于Hodge理论的那些想法之生成中,\dmark{}\origpage{312}Motif理论所起过的角色。我对自己说:倘若他有意愿提及我可能在他身边扮演过的角色\fnote{45}{(*)}{(5月30日)而且直到几星期前,我还在系统地淡化这一角色。对此可参见5月27日的“格格不入者”一注(n\textsuperscript{o}~67$'$),我在那里头一回觉察到我身上的这份态度,并隐约瞥见其意味。},他无疑会借着构成其学位论文工作的“Hodge理论II”来做——那正是提及这类事情的天赐良机\fnote{46}{(**)}{(5月30日)我也不记得曾被邀参加其论文答辩委员会。葬礼那时已然进行得热火朝天……}。我刚刚看到:他经由这句碑文式的话\fnote{47}{(***)}{Serre也经由指涉记号[3],在同一行里含蓄地露面——好奇的读者,在“Hodge理论I”的参考文献里能找到他的名字。这条草草的对照线,想必是从1968年到今天之间唯一的一处:对它一口气提到的那些“源头”(Serre(别名[3])、motifs、Grothendieck……)——无论多么谜语式——的一点影射。(5月28日)然而此后我又撞见另一处这样的影射;鉴于那个场合之特殊,它格外有趣。对此可参见“葬礼悼词(1)——或恭维”一注(n\textsuperscript{o}~104),以及紧随其前的“掘墓人——或全体教众”一注(n\textsuperscript{o}~97)的结尾——后者安放了那个“特殊的场合”。}——影射“Grothendieck的猜想性Motif理论”——一劳永逸地尽了提及我的手续,外带一个出处(Demazure在Bourbaki研讨班上的报告)。

又是无话可说!他从未想到要指明:他学得这套理论(且别忘了,整个是猜想性的!)是通过\textbf{别的来源},而不是Demazure那份寒碜的文本——后者绝无可能给出一套如此丰富的理论(整个是猜想性的!)的图像;那套理论,以底纹的形式贯穿Deligne其后围绕权的瑜伽的全部著作——且待那部海盗文集LN~900的步步升级, MotifGalois群(这回连一行含着亡故者名字的简短出处都没有……)才终于(十五年后)在那里开棺出土。

反思之后,在这句简短的引文里,我认出同一套“暂停!”做派——一句纯礼节性的引用,为的是把手续办清;附带的那个出处,绝无照亮读者之效(在这个案子里,即照出与那些恰好需要藏掖的想法之间那些分明而深刻的关系\dmark{}\origpage{313}——而它们在随后的十二年里也确实保持着藏掖),却\textbf{有欺瞒之效}\fnote{48}{(*)}{写这些文字时,一个联想不由分说地来到:早两年发生的、围绕“权”的第一起显影事件;它在“罐装的权与十二年的秘密”一注(n\textsuperscript{o}~49)的开头被谈到,并在“驱逐”一注(n\textsuperscript{o}~63)开头处得到更详细的论述。至于一般的“暂停!”做派,见“暂停!”一注(n\textsuperscript{o}~76,原书如此,题名对应注77)的反思。这套做派,开始对我变得非常熟悉了!}。

\notehead{78$'_{2}$}{}

这份文本\fnote{49}{(**)}{“代数簇上同调中的权”,P.~Deligne,1974年温哥华大会,《会议文集》,第78--85页。}我无需亲手捧读(我是仅仅几个星期前才知道它存在的),就知道我的名字不在其中。Serre的名字同样不在——而最早隐约瞥见“权的哲学”的人,正是他;那套哲学,是我随后以极大的细节提炼出来的。

\subsection{老板}

\notehead{78$''$}{}

(6月3日)Zoghman向我解释:他只是逐渐地、起初还模模糊糊地,才意识到围绕我的著作所做的那场“巧取”。1975年Verdier给他的手稿(见“好的参考文献”一注,n\textsuperscript{o}~82)对他曾如同天赐,尤其向他引入了可构造性观念及其本质性质,以及双对偶定理——他从后者取得启发,做出了D-模语境里的双对偶定理(或“局部对偶”定理)。只是多年之后,在读SGA~5时(屠杀版诚然,但还没有被屠杀到足以瞒过一个像他那样留心的读者),他才开始觉察到什么。长年以来,他对那位遥远的长辈满怀钦佩与感激,深信他大量汲取的那些观念属于后者。甚至似乎有多年之久,他确实深信:那套所谓“Verdier对偶”的对偶理论确实归Verdier——或至少归“Serre--Verdier”;同样,他所称“Poincaré--Verdier对偶”的那个对偶想法,也确实归Verdier。直到1979年前后(他答辩之年),他才开始觉察有什么不对——但我料想,他必定小心地不对其显赫的“老板”流露半点,一如在我们1980年2月与1983年6月会面时,不对我也一样。直到“反常”研讨会,1981年6月,当他开始感觉到\dmark{}\origpage{314}自己本人的著作正在被人藏掖,他才同时开始更清楚地认识到:自己闯进的是怎样一个世界\fnote{50}{(*)}{Zoghman那时对他前老板的评价,已低落到这种地步:他甚至深信,Verdier六十年代所做的一切(我在注n\textsuperscript{o}~81“赊账的学位论文与全险”的一条脚注里评阅过),或多或少是我在口授、或至少在提示。}!肯定,对他而言,我必属于这个世界:在那里,我的旧学生们(或至少其中某些)占据高位,以对待亡师的同样漫不经心,劫掠着遗孤学生。倘若真有区别,那不过是:我是亡者,而他们活得好好的,并且以确凿的方式证明着这一点……

我可以想象,甚至在“反常”研讨会之后,Zoghman仍难以相信自身健全认识能力的见证——尽管那见证分明地教给了他发生了什么。那篇大名鼎鼎的、由B.~Teissier与其“老板(原文如此)”Verdier署名的研讨会会议录导言,他直到1984年1月才亲手拿到。在驳回那份明摆着的事实近三年之后,震动想必愈发剧烈——我自以为是这样理解的。两个月之后我才重新联络他,3月底给他寄去“我的孤儿们”与“拒受遗产——或矛盾的代价”各注;又过了一个月,他才终于决定向我“亮底”,把“反常层会议之欺瞒”告诉我。

\subsection{我的朋友们}

\notehead{79}{}

瞧,我正准备把这场反思收尾并公之于众——它将终结Zoghman自己在那场他付出代价的巧取周围保持的沉默;而那场巧取的阴暗收益,他也照单全收\fnote{51}{(**)}{(5月30日)我提请注意:这场反思所由生的,是我朋友身上的一些心境;它们如今看来已被超越。(可与注n\textsuperscript{o}~78$'$的5月30日两条脚注对读。)}。也许它会令他不快,一如也许会令我的朋友皮埃尔不快——一旦它完成、文本誊清并印出,我就会把它亲手送到他手里\fnote{52}{(***)}{我原不曾想到,在我前面尚余的年月里,还有机会回到首都住上几天。但十多年来,我的朋友皮埃尔频频移步,到我蛰居的乡野深处来看我;于是这一回,趁着这个例外场合,轮到我移步了——那同时也是对一个屡次发出、而我从未应邀的邀请的回应。}。我能献给我朋友Zoghman与朋友皮埃尔的\dmark{}\origpage{315}最好的东西,他们两位多半会把它当作最坏的东西来接受:当作一场灾祸,或一桩冒犯。更何况,我的见证是公开的——一如他们两人的沉默,都是公开的仪式性行为,使这一位和另一位都须承担责任。

拒绝或接纳我的见证,是他们的选择;对Jean-Louis亦然——我曾把他算作我的朋友,一如今天的Zoghman与皮埃尔。这些选择与我切身相关,却不是我的选择。我毫无预言它们的兴致。我不久就会知道;我怀着强烈的兴趣、一份悬念——却没有一丝焦虑的阴影——期待着到来的星期与月份将带给我什么。我唯一的关切、我唯一的责任,是:我所献出的,确实是我所能献出的最好的东西——也就是说:真。

也许有人会惊讶:我竟毫不容情地谈论我以“朋友”之名相称的人;并在这名字里看出一句套话,甚至看出其中并不存在的讥讽声调。当我把Zoghman~Mebkhout或Pierre~Deligne称作“朋友”时,那是对同情的、亲切的、敬重的感情的重申——它们在我写下这些文字的时刻,在我身上在场。敬重告诉我:我无须“容情”对待一位朋友,一如我无须“容情”对待自己——一如我,他有资格与那份寒微的真相相遇;一如我,他无须容情。

倘若我不以“朋友”之名指涉Jean-Louis~Verdier,那绝不是因为我认为他比我的朋友Zoghman与皮埃尔、或比我自己更不“好”、更不“配”,而是因为生活恰巧把我们分开了。十五年前、乃至更久以前把我和他系在一起的同情的、亲切的感情,已被时间或多或少磨蚀,不曾有机会经由哪怕一点个人的接触而复活。我为重建这种接触做过的寥寥几次尝试,没有得到回音;我不知道,这些反思的阅读,会不会使一段已然僵死的关系复活。但即便此刻他于我不再是“朋友”,我也不认为:我像对待自己、对待我的朋友们那样毫不容情地对待他,就是不敬;而且我深知,反其道而行,于他、于任何人,都无益。更不必说:无论他还是我的朋友皮埃尔,倘若他们真想“自卫”(或进攻)而不肯冒险看自己一眼,并不缺手段,也不缺后援。更不必说:在他们有条件\dmark{}\origpage{316}令人气馁、或碾压人的地方,这一位与那一位,都不止一次这样做过,毫不容情,也毫无怜悯。

\subsection{大部头与上流社会(或:鱼目与珍珠……)}

\notehead{80}{}

(5月9日)况且,也该是时候了:我到底给出那条大名鼎鼎的Riemann--Hilbert——(不具名的Deligne)——亚当与夏娃——上帝——(尤其不是Mebkhout)定理的出处:人人都在大量引用它(包括我自己),而显然还没有人想到要问它在哪里得到证明。我自以为从我的朋友Zoghman那里懂得了:那条“值得纪念的定理”就在他的论文里;我果然在其目录里找到了它,顶着(确实土气、确实配得上莽汉的)名字“一个范畴等价”,第三章,第3节,第75页。更倒霉的是,它甚至没资格叫“定理”,而是叫“命题3.3”(而且更糟的是,我的名字出现在同一页上,还带着着重号!)。我承认,由于没有读过前面那75页来弄清脉络,我甚至不完全有把握是不是它——Zoghman向我确认是的;我信他\fnote{53}{(*)}{(1985年4月17日)最终显明:人们通常使用的“好上帝定理”的形式,并非此处所引定理的形式,而是一个相近的形式,可以用同样的方法证明。见“一个视野的绽放——或闯入者”一注(n\textsuperscript{o}~171$_{1}$),尤其其中今日所写的那条脚注。}。证明(据说)是同一篇论文第五章的对象——论文于1979年2月15日在巴黎第七大学答辩,答辩委员会由 D.~Bertrand、R.~Godement、G.~Houzel、Le~Dung~Tráng、J.L.~Verdier组成。有兴趣的相关人士,倘若尚未收到作者寄赠的副本(凡他无论对错地疑心可能感兴趣的人,他都寄了论文),只消向他开口,他将乐于从命……他当然给每一位我的前上同调学家学生都寄了副本;没有一位给过活着的音讯。想必他们其间都改了行,真不巧……

得承认:Zoghman做买卖的本事,确实不怎么样,不会以清楚而诱人的方式呈献他的货色——这是学得来的手艺;而他不像我的前学生们有那份运气,得以跟一位行业高手、且从不吝惜时间的高手学徒。可他也无从抱怨:他那“三次会谈”是得过了;说不定哪天,哪位“权威”会想起来,甚至为他那份难啃的大部头回一句“已收到”。他自己想必也发觉了:这份大部头很难啃(即便它对Riemann与Hilbert都并非白费……):他在法国科学院报告上发了一条短文——那毕竟短些——以引起人们对其大名鼎鼎定理的注意;题目嘛,您猜一百次:《论Hilbert--Riemann问题》!我早就知道,我的朋友Pierre~Deligne的历史修养并不比我强:他只需把年代顺序拨正\dmark{}\origpage{317},再贡献出那个漂亮的民间传说式的称呼“对应”,戏法就完成了——Zoghman,这可是他自找的……这条短文是1980年3月3日的,Série~A,第415--417页。

Verdier想必是在他“恩准”其学生的那“三次会谈”之一里(或在答辩时)得知这条定理的;但倘若真要问,他想必什么也没察觉。Deligne呢,他到底何时察觉到了什么,我说不好;但可以肯定:1980年10月他已知情,而据他自己说,Bernstein与Beilinson也是。Mebkhout还亲自去了莫斯科,向Beilinson与Bernstein(唯恐他们读起来费劲)详详细细解释其结果。我不知道是他们还是Deligne读过那篇论文或其后那篇科学院短文;但想必他们到底弄懂了里面的东西,因为次年那场大名鼎鼎的吕米尼“研讨会”恰好围着它转——出于天底下最大的偶然。

总结一下,并把我情报部门最新惠予的讯息算在内:对局势完全知情、并参与了那场所谓“反常层会议”之欺瞒的,至少有五人,即(按演员姓名的字母顺序)A.A.~Beilinson、J.~Bernstein、P.~Deligne、J.L.~Verdier与 Z.~Mebkhout——外加整整一场研讨会的与会者(原文称他们为acultes,疑为adultes“成年人”):肯定还都是光彩夺目的数学家,他们分明不求别的,只求被欺瞒、把鱼目当珍珠\fnote{54}{(*)}{(6月3日)实际上,看来研讨会的全体与会者,无一例外,都已在现场被告知了局势。对此可参见今天写成的“研讨会”一注,n\textsuperscript{o}~75$'$。}。这再次证明:我们数学家——从显赫的奖章得主到无名的学生——并不比张三李四多一丁点聪明或明智。

\section{VIII.~学生——别名老板}

\subsection{赊账的学位论文与全险}

\notehead{81}{}

\dmark{}\origpage{319}(5月8日)我觉得是时候了,该对“幽灵论文”这桩事作更周详的表述了;此前我只在两条注里(注48与63$'''$)“顺带”谈过。一位不细心的或不怀好意的读者,也许会说,我同时责备我的前学生 J.L.~Verdier两件自相矛盾的事:既“埋葬”了导出范畴,又把它们“发表”了(在SGA~4½里)并自居其父;正如同一读者会说,我责备 P.~Deligne,既“埋葬”了Motif,又把它们开棺出土(在LN~900里)。因此,或许并非多余:对1960年至今的局势作一番回顾。

1960或1961年前后,我向Verdier提议一项可能的论文工作:发展同调代数的新奠基,基于导出范畴的形式主义——那是我在此前几年里,为概形语境中的凝聚对偶形式主义之需而提炼并使用的。当时已经说清楚:在我向他提出的纲领里,前景中并没有什么严肃的技术困难;主要是概念性的工作,其出发点已然在手,或许需要相当规模的展开——规模堪比Cartan--Eilenberg那部奠基之作。Verdier接受了这个题目。他的奠基工作进展顺利,于1963年落实为一份关于导出范畴与三角范畴的“第0状态”,由IHES油印。那是一份50页的文本,1977年作为附录复制进SGA~4½(如注63$'''$所说)\fnote{55}{(*)}{单是这份文本,对一位天资出众的青年研究者两三年的工作而言,似乎嫌瘦。但Verdier当时的大部分精力,用于修习同调代数与代数几何不可或缺的基础——尤其通过听我的研讨班,以及两两对谈的工作。他对对偶形式主义的贡献(见下文)来得较晚:在我与Artin于SGA(1963/64)里详细发展了平展对偶的形式主义之后,当时我建议他(在其导出范畴的奠基工作之旁),在“寻常”拓扑空间、及这类空间的可光滑化态射的框架里,发展同一形式主义。大约在我以SGA~1开创“代数几何研讨班”丛书之际(1960年),Verdier联络了我,与Jean~Giraud、Michel~Demazure同时——问我有没有工作给他们;而他们敲的正是对的门!一件曾打动我的巧合:打从我写“我的孤儿们”一注(n\textsuperscript{o}~46)起,我记得他们三位联络我的时候,恰好刚组成一个小研讨班,名为“孤儿研讨班”(主题:自守函数,计算进路(原文approche calculs à brin de zinc,其中à brin de zinc措辞待考)),因为他们的老板(或CNRS的教父?)不打招呼就走了一年,让他们意犹未尽,又有些茫然。这片空,很快就被填上了……}

\dmark{}\origpage{320}答辩之所以没有在1963年、而是1967年举行,是因为难以想象:这份50页、尚待来临的奠基工作的胚芽,能构成一篇国家博士论文——这个问题当然根本不会被提出。也是出于同一理由,1967年6月14日答辩时(评审委员会由 C.~Chevalley、R.~Godement与我——主席——组成),根本不打算把这项工作作为论文呈上。提交委员会的文本,17页(加参考文献),以一项正在撰写的大工程的\textbf{导言}的面目出现。它勾勒这项工作的主要想法,并把这些想法安置在其众多使用的语境之中。第10、11页,对这项奠基工作所计划的章与节作了周详的描述。

倘若 J.L.~Verdier的理学博士头衔,是凭这份17页、勾勒着连他自己都说并非出于己手的想法的文本授予的\fnote{56}{(*)}{论文开头写道:“本论文是在 A.~Grothendieck的指导下写成的。它所包含的本质观念出自他。没有他最初的启发、他始终如一的帮助、他富有成果的批评,我不可能把这项工作进行到底。愿他在这里找到我深深的感激的表达。\par 我感谢Claude~Chevalley肯担任我论文答辩委员会主席,并有耐心读完这份文本。\par 我感谢 R.~Godement与 N.~Bourbaki,把我领进数学之门。”——按:“本论文”一词,恐怕只能指那项着手了的奠基工作的全体——所提交的文本构成其导言;而严格说来,那项工作在答辩之时并未“进行到底”。(5月30日)这一不自洽,恰好映出那份局势的暧昧;而那份局势,我作为论文导师——并且(照我手头那份论文封面上的说法)作为委员会主席——是第一责任人。我对一位光彩夺目的学生不够“严格”,有所纵容;这一纵容的方向,与我对Deligne所表现的(见“格格不入者”一注,n\textsuperscript{o}~67$'$)如出一辙,并且出了一份力,使同样的果实结了出来。},那显然是委员会与他之间的一纸本着相互信任订立的契约\dmark{}\origpage{321}:他承诺把这项工作做到底、并公之于众——对这项工作,他呈献了一篇光辉的导言。这份契约,候选人没有履行\fnote{57}{(*)}{越发引人注目的是:1983年12月, J.L.~Verdier竟拒绝了我请他加入Contou-Carrère论文答辩委员会的提议——委员会里还有 J.~Giraud,而我任研究导师;他的理由是:那篇论文(分明已全部写成,并经 J.~Giraud细心读过)与那个委员会,不经\textbf{巴黎}各大学论文委员会(原文如此)的核准,不具足够的严肃性的担保。}:他宣布的那份文本——一份按照一个已被证明有效的新观点写成的同调代数奠基文本——从未发表。

显然,倘若Verdier1961至1967年间的工作,仅限于写下1963年那份骨架式的“第0状态”,委员会绝不会想到接受这份“赊账的论文”。当时,其工作的成文想必已有相当的进展,可望一两年内完成;而出于实际的缘由,让Verdier不必等到那项本该为之奠基的工作完成就取得头衔,显得合宜。

还得补上一句:1964至1967年间,Verdier对对偶形式主义做出过几项有趣的贡献(注81$_{1}$);这些贡献,连同他本该继续的那项奠基工作,可以佐证人们给他赊的这笔账。单是他的那些对偶贡献,严格说,本就足以构成一篇像样的博士论文。然而这样一篇论文,绝不合我素来命题的风格:那些题目,都是对一个我感到其需要与紧迫的理论,做系统的、贯彻到底的发展(注81$_{2}$,原书误作82$_{2}$)。我不记得Verdier曾想过提出这样一篇“凭既有成果立论”的论文;我怀疑我是否会接受——因为这样一篇论文,与他我之间立下的“契约”毫无对应:我把导出范畴这个美好的题目托付给他,责成他发展大面积的奠基。

我承认我全部的责任,作为 J.L.~Verdier的论文导师与委员会主席:我轻率地(与 C.~Chevalley、R.~Godement一同——他们信任我所给的担保)把博士头衔授予一项尚未做成的工作\fnote{58}{(**)}{在这份责任之外,我还该补上另一份:在其后的两年里(在我退出数学舞台之前),我未曾督察Verdier是否确实履行了他所立的契约。得承认,我的精力当时全部投在我自己承担的奠基工作上——更不必说Motif的与其他的种种反思——以致那桩提醒他人尽其义务的不快差事,我顾不上多想。Verdier放弃出版那项预定工作的决定,我是七十年代初才知道的——那时我已全然不再关心数学,“反应”的念头根本无从起。}。

\dmark{}\origpage{322}倘若我今天确认这份轻率的某些果实,我原无权抱怨。但这并不妨碍我公开地做这份确认,并且声明:我的前学生 J.L.~Verdier的所作所为,只关他自己的责任,不关任何别人。

不履行对我、对那位曾信任他的委员会所立的契约,是一种埋葬的方式:埋葬我引入、并由他承揽以一项大工程去奠基的导出范畴的观点。这项工作也许做过,却从未交到使用者手里。这是对一整组他自己也曾协助发展的想法“画一道叉”。

Mebkhout的著作重新拾起导出范畴的观念,没有从Verdier那里得到任何鼓励(同样也没有从我其他任何一位身为“上同调权威”的学生那里)。对导出范畴的事实抵制,在我看来,直到1981年前后都是彻底的\fnote{59}{(*)}{(5月30日)这些带几分存疑语气的措辞,其实并不合宜。Zoghman~Mebkhout向我确认了(他是付出代价才知道的):我对“Grothendieck风格”的同调代数所受待遇的存疑表述,与实情完全相符。}——就在那年,它们随着需求的突然推挤,在那场值得纪念的吕米尼“研讨会”里强力重返(见注75)。

然而,Verdier那份“论文”的“第0状态”,早在此四年之前,即1977年,就已作为附录刊于SGA~4½卷(见注n\textsuperscript{o}~63$'''$)——也就是其论文答辩十年之后,而在(据我所知)Mebkhout是唯一在其工作中使用导出范畴的人的时刻,逆着此前七年的风尚。如无差错,他一直是唯一的一个,直到那场已提到的研讨会上围绕那条大名鼎鼎的“Riemann--Hilbert对应”的大“抢滩”——在那里,Deligne,别名Riemann--Hilbert,以这条“对应”之父的面目出现——原文如此;而Verdier(连同他那被其慷慨朋友大量引用的天赐“第0状态”),以导出范畴与两千年风格同调代数之父的面目出现——我这个人则绝口不提,更不必提Mebkhout\fnote{60}{(**)}{可与“同伙”与“不义——或归来的意义”两注(n\textsuperscript{o}~63$'''$与75)中的评语对读。}。

\dmark{}\origpage{323}在这些事件的光照下,我想我懂得了这份“第0状态”出人意表的发表缘由——SGA~4½的导言(还是那位一以贯之的朋友执笔的)说它“已经无处可寻”,而当时除了(也许)Zoghman~Mebkhout,谁也不关心去“找”它\fnote{61}{(*)}{不过,恰恰是4月底翻阅我刚收到的一份 Z.~Mebkhout著作的参考文献时,我才知道这份“第0状态”业已发表——那份上了年纪的文本的存在,我甚至已经忘了……}。于是恰恰有这么一个倒霉人,在他自己的角落里,逆着所有人;执拗地使用着这些过了时的观念,却没人弄得清他想达到什么——他终究那么执拗,以至于一个疑心开始冒头:万一哪天这个家伙端出几样够分量的东西呢,谁知道……毕竟,他斗胆引为灵感之源(在大师的著作旁边)的那个人,当年可是用这一切证明过、发现过东西的;那些东西,人们无法统统装作忘记,即便忘了它们的作者——而且大师本人,Jean-Louis~Verdier亲自,不正是靠着那条“Lefschetz--Verdier公式”起家奔向荣耀的吗?而没有这些该进纸篓的观念,那条公式他连写下来都费劲,更遑论证明……

我那影响巨大的前学生在这将近十年里(自从他摆脱了那道烦人的手续……)\textbf{押注导出范畴没有前途}、并且还将这样押注到关键时刻(那场大名鼎鼎的研讨会)的钟点,他想必认定谨慎为上(谁知道呢……),抢先于可能到来的事件——总之,来一份“全险”:发表的(当然不是那项本该终有一日构成一篇论文的大工程,而是)一份“证人文本”、一件“以备不时之需……”的证据;一份能证明他对一个\textbf{孤儿}之父权的文本——那孤儿,他曾乐于厌弃,并在等待事件来临之际,继续否认\fnote{62}{(**)}{倘若 J.L.~Verdier真有让埋葬了七年之久的导出范畴瑜伽为人所知的愿望,他本该选择发表构成其论文的那篇导言文本,而不是一份无人问津、只有放在那部瑜伽及其众多使用的底色上才见趣味的技术文本。但人们懂得:他绝无意愿在那份50页的证人文本旁边,附上其论文的那17页——里面有如今叫人难堪的断语,论的是那个绝不可提名的人的角色……}。

\notehead{81$_{1}$}{}

\dmark{}\origpage{324}所说的贡献是:
\begingroup
\renewcommand{\labelenumi}{\arabic{enumi})}
\begin{enumerate}
\item 局部紧空间语境中对偶形式主义的奠基;
\item Galois模语境中的相应形式主义(与J.~Tate合作);
\item 那条所谓Lefschetz--Verdier的\textbf{不动点公式};
\item 局部紧空间中的对偶。
\end{enumerate}
\endgroup

贡献2)与3)相对已知的东西,构成一个“意外”。最重要的贡献,依我看是3)。它的证明,从对偶形式主义(无论“离散”还是“连续”系数)轻易得出;这并不妨碍它构成我们上同调武库中“万用”公式的重要一味。这条公式的存在,是Verdier发现的;对我而言,那是一次(愉快的!)意外\fnoteplainmark{63}{(*)}{(1985年4月19日)关于这条漂亮的公式、它的角色、以及它在葬礼中的离奇坎坷,我在《收获与播种》第四部分的“真正的数学……”“……与‘无意义’”“把戏与创造”三注(n\textsuperscript{o}~169$_{5}$、169$_{6}$、169$'_{6}$)里还要谈到。}。

局部紧空间语境中的对偶形式主义,就其实质而言,是对我在概形平展上同调语境里所做之事的“势在必行的”移植(而且没有那种一切尚待去做的局面所固有的困难)。不过它带来一个有趣的新想法:把函子$f^!$(不先对$f$作可光滑化处理)直接构造为$Rf_!$的右伴随,外带一个存在性定理。这一办法后来被Deligne在平展上同调里采用,使他不借可光滑化假设,在那个框架里定义$f^!$。

这些评语,我想,足以说明:到1967年,Verdier已证明了他做原创数学工作的能力;这当然是为他赊账的决定性因素。

\notehead{81$_{2}$}{}

作为另一个例子,我指出:在局部紧空间的语境里,依“六种运算”与导出范畴那套“万用”形式主义的精神,对偶形式主义的详尽发展——Verdier在Bourbaki研讨班上的报告,可算它的胚芽。即便只就拓扑\textbf{流形}的语境而言,据我所知,Poincaré对偶形式主义至今也没有一部令人满意的参考文本。

\dmark{}\origpage{325}(6月5日)还有两个方向,我遗憾地看到,Verdier不曾认为值得把工作做到底——那工作他起手的势头已足够强,本可\textbf{收取其信誉}(我是指:在离散系数与局部紧拓扑空间的语境里,开启一套对偶形式主义),尽管那些本质观念并非出自他,而且他(一如对导出范畴)全无心肝去做一项任务的\textbf{仆人}、把一套完整的形式主义交到使用者手里(一如我在SGA~4、SGA~5、SGA~7三场研讨班里竭力所做的)。我所预见并建议他发展的对偶纲领,置于一般拓扑空间(不必局部紧)的框架、以及其间“分离的”、局部“可光滑化”的态射(即局部地,源嵌入$Y\times\mathbb{R}^n$,$Y$为靶空间)之中。这显然是同\textbf{任意}概形的平展上同调框架的类比所明摆着提示的。Verdier看得出:在局部紧空间的框架里,态射局部可光滑化的假设是多余的(这曾是一件意外)。但这并不妨碍:局部紧空间的语境(从而排除并非局部紧的“参数空间”)分明碍手碍脚。一个更令人满意的语境,应当同时涵盖Verdier所选的情形和我预见的情形:拓扑空间(乃至topos?)多少任意,而态射$f:X\to Y$须满足两条限制:1)分离;2)“局部可紧化”,即$X$局部嵌入$Y\times K$,$K$紧。

在这个语境里,一个“被允许的”态射的纤维,是任意的局部紧空间。下一步,则是允许$X$与$Y$不是拓扑空间,而是“拓扑多重簇”(即“局部如同一个拓扑空间”的topos),乃至任意的topos,同时对态射作合宜的限制(有待明确),使纤维成为\textbf{局部紧的多重簇},必要时附加补充条件(也许近于佐竹的$G$-流形的观点),比如(而且在万不得已时!)局部地取$(X,G)$之形,$X$为带一个\textbf{有限}算子群$G$的紧空间。

据我所知,连“寻常的”\dmark{}\origpage{326}Poincaré对偶,也还不曾在紧光滑拓扑多重簇(光滑:即局部如同一个拓扑流形)的情形里得到发展。有限群的分类空间的情形似乎表明:人们几乎只能指望在模去挠性的意义下得到一条(绝对整体的)对偶定理;更确切地说,须取一个$\mathbb{Q}$-代数作为系数环。除这一限制外,倘若Poincaré对偶(“六种运算”风格)在这个语境里照原样成立,我不会惊讶。从来无人看过它一眼(除了那些不可救药的微分几何学家——他们装模作样地看叶状结构的“叶空间”的上同调),这并不奇怪:看看我的上同调学家学生们——Deligne与Verdier领衔——对多重簇观念本身所设置的普遍抵制。

老实说,缺的是这样一场奠基的反思:在任意topos及其上的“离散”系数层的语境里,描述(倘能做到)topos态射的“真性”“光滑性”“局部真性”“分离性”诸观念,从而提炼出一个“被允许的”topos态射$f:X\to Y$的观念——对其,$Rf_!$与$Lf^!$两种运算(前者为后者的左伴随)各有意义,从而取得六种运算形式主义的惯常性质。这里,topos被视为不带环层的,或也许带着环层(必要时假定为常值或局部常值);先假定(至少起初)带环层的topos态射$f:(X,\mathcal{A})\to(Y,\mathcal{B})$满足$f^{-1}(\mathcal{B})\to\mathcal{A}$\textbf{是一个同构}(注81$_{3}$)。上面的反思提示:当我们把自己限于特征零的系数环层(即$\mathbb{Q}$-代数)时(译按:特征零本身不足以保证环为$\mathbb{Q}$-代数,例如$\mathbb{Z}$),“被允许的态射”的观念可以明显放宽,以涵盖这样的“纤维”:比如多重簇(拓扑的或概形的),而非寻常的“空间”(拓扑的或概形的)。

这个方向的最初尝试(除我处理过、随后Verdier依同一模型处理过的情形之外),归Tate与Verdier:在离散群或投射有限群的语境里。这一初步尝试的记忆,曾鼓励我去年继续这个方向上的反思:在小范畴(它们推广离散群、充当同伦模型)的语境里。虽然没有走远,这场反思仍足以使我相信:必须存在一套完整的六种运算形式主义,在小范畴的范畴(Cat)的语境里。(对此可参见\dmark{}\origpage{327}《追寻栈》,第VII章,第136、137节。)这样一种理论在(Cat)、乃至在Pro(Cat)里的发展,一如这类理论在拓扑或概形的空间与多重簇的语境里的发展,对我而言,其主要趣味在于:它是朝着在一般topos的语境里更好地理解“离散对偶”迈出的一步。

去年Illusie向我透露:他在半单纯空间(或概形)的情形里,曾同对偶理论的种种疑难苦斗。那听上去分明还是同一回事——在一个个案里察觉到一套六种运算形式主义的存在,并理解它。但似乎,一场奠基反思的前景本身,就有本事使我的前学生们——至少我的上同调学家学生们——人人浑身冰凉。我当年同他们费过力气,是怀着这样的信念:他们不会(在概念工作的层面上)恰好在与我同行的那个地点戛然而止;每当一个新的局面显出他们与其伙伴同我做的工作不够用时,不会只是搓手发愁。人们所做的概念工作,从长远看\textbf{总是}不够用的;数学的进步,恰恰在于重新拾起它、超越它——除此之外别无他途。1955到1970年间,年复一年,我发现前些年所做的已不敷需要,便立刻重新伏案——至少当别人(比如Mike~Artin,以他自己意义下的“代数簇”观点)还没有抢先把这件事做掉的时候。但似乎,我的学生们把我给他们做的榜样也一并埋葬了,连同我这个人、我的著作。

\notehead{81$_{3}$}{}

我记得,在(比如说)平展上同调的六种运算形式主义里,充当系数的环层局部常值的假设是多余的——本质的假设是:它们是素于各剩余特征的挠层,\textbf{并且}$f^{-1}(\mathcal{B})\to\mathcal{A}$是一个同构。一旦放弃后一假设,人们便进入一门(据我所知从未明言的)理论:它把“空间离散”对偶与“凝聚”对偶(关涉系数环层及其同态)混合起来。于是,人们设想:在概形(或更一般的topos)$X$、$Y$上,把系数环层$\mathcal{A}$、$\mathcal{B}$换成$X$、$Y$上的相对概形(不必仿射)$X'$、$Y'$,把带环层的topos态射\dmark{}\origpage{328}$(X,\mathcal{A})\to(Y,\mathcal{B})$换成如下类型的交换图表
\[
\begin{tikzcd}[row sep=2em, column sep=3em]
X' \arrow[r] \arrow[d] & X \arrow[d] \\
Y' \arrow[r] & Y
\end{tikzcd}
\]
并带一套这一类型语境中的“六种运算”形式主义。当$X$、$Y$等为点的topos时,人们应当找回惯常的凝聚对偶。

\subsection{好的参考文献}

\notehead{82}{}

(5月8日)说的是 J.L.~Verdier的文章“与一个闭链相配的同调类”,1976年刊于Astérisque~n\textsuperscript{o}~36(SMF),第101--151页。从某种意义上说,这篇相当不可思议的文章(尽管如今已没有什么应当再使我惊讶……)与Deligne等人的“反常文章”互为表里。除一点保留外,它实际上等于花五十页,在一个略有不同的语境里,\textbf{照抄}我在十年或十五年前纵横发展过的观念、构造与论证——术语、记号,应有尽有,一字不差!我几乎以为自己回到了1965/66年间SGA~5研讨班的一场:那些东西在那里被讲明过(显然,详尽到与会者吃不消\fnote{64}{(*)}{这一方向上的评语,见注n\textsuperscript{o}~68、68$'$“信号”与“颠倒”,我在其中检视这场研讨班成文工作的离奇坎坷,以及这些坎坷与Deligne“SGA~4½之役”的关系。随后的反思,向我显出这些坎坷、以及母研讨班经Verdier与Deligne合谋肢解的又一个未曾预料的侧面。两人的出版物为这场肢解揭幕,分别刊于1976与1977年——它们构成给Illusie的“绿灯”,让他去筹备(十一年之后……)SGA~5的出版(那一卷,如Deligne在SGA~4½里所说,“可以看作一系列枝蔓,其中一些还非常有趣”)。})整整一年。至少在那场研讨班之后,对所有多少“在局中”的人而言\fnote{65}{(**)}{关于我回到这份“仓促”印象的一场反思,见“沉默”一注(n\textsuperscript{o}~84)。},这一切都属于“众所周知”的领域。Verdier当然到场了,一如Deligne(唯一一个从不掉队的人——尽管那是他头一回把脚迈\dmark{}\origpage{329}进我的研讨班\fnote{66}{(*)}{这场研讨班的那一年(我想),正是我结识Deligne的一年;他当时该有十九岁。他很快“入了局”,甚至揽下成文我前一年平展对偶报告的差事(那些报告,他想必经由我的讲解与我的笔记而知),以及关于与一个闭链相配的上同调类的那场报告——后文在引过的“颠倒”一注(n\textsuperscript{o}~68$'$)里谈到,在本注里也还要谈到一点。以他那般手段、对主题的那份完全驾驭,他竟等了十一年才做成文,然后不经告知,把它收进其SGA~4½——如今回顾,这向我表明:早在1966年(而不仅是我曾以为的1968年——见“驱逐”一注,n\textsuperscript{o}~63)——即我们相识的头一年,我的朋友对我的关系中,便已存在一种深层的暧昧,而且从那一刻起就以全然清晰的方式表达着;而我一直回避去面对它,直到今天!}——总得有人这样……)。倒是有一件事:到1976年,那场大名鼎鼎研讨班的“成文(原文如此)”,由那些早已厌烦的“志愿者(原文如此)”拖着,已经拖了十年——我现在明白:其中一位“志愿者”到底还是按他自己的方式揽下了“成文”的差事,早在SGA~5于1977年出版之前!看来这位倒霉研讨班的坎坷,方便的还不止Deligne一人——他以自己的方式,从一盘散沙的局面里渔利。不过在那个时刻,Deligne倒还小心:一边把SGA~5的一场关键报告拆下来、名正言顺地并入其SGA~4½,一边到底还在其成文(关于与一个闭链相配的上同调类)里提一句“依Grothendieck的一次报告”。(诚然,他从中得到一份补偿:得以据此把我列为他的“合作者”!——见“颠倒”一注,n\textsuperscript{o}~68$'$。)

回到与一个闭链相配的\textbf{同调}类(别混淆!)——按标题,它是Verdier文章的对象——这个形式主义,我在口头研讨班上,分几次报告,以穷尽的细节发展过;而且面对的听众分明在求饶(照例只有Deligne一人,始终神采奕奕、毫无倦容……)。那是那一年我在平展框架的对偶形式主义上发展过的无数“长练习”之一;我感到需要完全驾驭一切在我看来必须彻底领会的要点。这里的趣味在于:拥有一个在不必正则的环绕概形上成立的形式主义——到正则情形的\textbf{上同调}类的过渡,以及与我那套使用具支撑上同调的老构造的联系(它立即给出与上积的相容性),都是直接的。我还发现:\dmark{}\origpage{330}研讨班的这一部分,属于未收入已刊版本的那一堆——想必Illusie(出一版拿得出手的(哼)本子的全部准备工作,到头来都压在他身上)该庆幸Verdier揽下了这件事,依样画葫芦(此处即:一字不改!)。

按照如今已经功德圆满的套话,“几乎无须说”:我的名字既不在正文里,也不在参考文献里(除了那条永恒的SGA~4引证——它含蓄地带着我,尽管人们早晚得找到替代……)。没有任何暗示提到一场与SGA~5这个缩写相应的“代数几何研讨班”——作者想必听说过它——尽管我分明记得见过他,殷勤地、懂事地做着笔记(像所有人一样;当然,除了Deligne……)。

况且,我说我的名字不见于正文,只夸张了一丁点:它有一次露面,神秘而碑文式,在第38页,第3.5节“基本的上同调类,相交”(人们终于抵达问题的死结了!)。那个出处是一句谜语式的话,其意我承认参不透:“系统地使用权复形(???又是这些该死的权!)的想法归Grothendieck,并经Deligne整理成型”——对这些神秘的“权复形”——据说我出过想法、而我此刻是头一回听说——再无任何解释。此后全书再未提及(此前37页里也从未提及)。谁懂谁懂!至于该节的内容,它一字不差地抄自十年前举行的那场SGA~5研讨班(而那时这一构造已有五六年的历史,见注n\textsuperscript{o}~68$'$);那场研讨班,他自然绝口不引。对Deligne的引证(他竟“整理成型”一条想法——当我的朋友还是中学生时,那条想法就已然现成!)是一朵“花”:这一想法想必之所以来到作者心头,是因为年轻的新来者Deligne确实揽下了成文我那场报告的差事(而且一搁十一年,为了人们知道的那些收益,见前引注)。这朵“花”,属于这对形影不离的朋友之间礼尚往来的交换。

不过,文章里有一个(想必)新颖而极有趣的结果(定理3.3.1,第9页):解析可构造离散层经由一个解析真态射的高阶直像的稳定性。Verdier约十五年前从我嘴里学得各种语境中的可构造性观念,以及\dmark{}\origpage{331}那条稳定性猜想——我在五十年代末提出过它(并对凡愿听的人谈起),那时我还没有结识他的荣幸。读这篇文章,一个不知情的读者(但这样的人已开始稀少……我恐怕又在重复自己)绝不会想到:作者并不是正在趁热端上他刚刚才发现的观念与陈述。他无须说那是他——不言自明。这就是那套大名鼎鼎的“暂停”做派,分明已经蔚然成风。

除了这个细节(依我之见,它合乎行业的新规范)之外,这篇文章想必总有十页上下(五十页里),围绕这个有趣的结果,呈献作者本人的工作。比例放到相应的大小,在Verdier身上、一如在Deligne身上,最打动我的,是他完全有能力做出漂亮的数学。即使在这篇令人痛心的文章里,前引定理也透出这一点。但他(一如其朋友)把自己保持在掘墓人的心境里,便与他显赫的朋友一样,只靠其手段中微不足道的一部分运转。一个(曾令我瞠目的)表面平庸的征兆,见于一位分明给出过机巧与嗅觉之证明的数学家:对他的“学生(原文如此)”Mebkhout的工作的分量,全无直觉——他乐于以自己的伟大俯视他,却从不曾自己做出深度与原创性堪比的工作\fnote{67}{(*)}{同一个令人瞠目的嗅觉迟钝,在这同一场合也见于Deligne:他仿佛到1980年才“感到风向”(Mebkhout那些想法的重要性),而Mebkhout在这一方向上自1974年起就工作着。我的朋友身上,自然嗅觉被自满所堵塞,我已有不止一次机会观察到;尤其自1977年(或1978年)起——那一年似乎构成第一个“转折”(对此可参见“两个转折”与“葬礼”两注,n\textsuperscript{o}~66、70)。}。这不是说他也许没有能力做到Mebkhout或我的地步。而是他从不曾给自己留一点机会去做大事——也就是说,放开缰绳,让一股热情牵引,而不是把数学及其天赋当作\textbf{工具},用以炫目、支配或碾压。直到此刻,他一直满足于照原样拾取那些现成而丰产的观念与观点。确实,他似乎已全然丧失了何为\textbf{数学创造}的感觉。

然而我记得,他与我一起工作的时候,这份感觉还在。没有任何外在于他的东西,妨碍这份感觉\dmark{}\origpage{332}重新浮出。一如其朋友——在那人身上,我常常感到那件精微而活泼之物的同样的隐没,被同样的自负所堵塞。

这篇50页的不可思议的文章,刊于一份有分量的杂志,为我照亮了“注记——或新伦理”那桩事故(第33节)新的一面:一篇\textbf{几页}的科学院短文——概括一项扎实而\textbf{原创}的工作,题材重要(依我浅见),是一位天资极高的青年数学家\textbf{两年工作}的果实——竟被两位权威以“毫无趣味”为由拒稿\fnote{68}{(*)}{细节见“棺材4——或没有鲜花也没有花冠的topos”一注,n\textsuperscript{o}~96。}。其中一位权威,不是别人,正是Pierre~Deligne——同一个Deligne,却肯一字不差地、亲自照抄我的一个学生那份不起眼的博士论文(他还恪尽职守地引用它)。(这份复制件,经一项显赫的署名拔高,成了那部同样显赫的丛书里“值得纪念的文集”LN~900中最厚重的一篇!对此可参见注52、67的末尾。)

分明,这幅“风俗画”日甚一日地丰满起来;而我为此甚至无须走出隐退、走街串巷去混迹“上流社会”。在几篇挑选得当的“大文本”里翻阅的零零星星几个钟头,已足以使我大开眼界……

\subsection{玩笑——或“权复形”}

\notehead{83}{}

(5月8--9日)我重新想到了那个“权复形”——它出现于Verdier值得纪念的文章里那条“出处——暂停”\fnote{69}{(**)}{见前一条注“好的参考文献”。}之中——一条俨然滑稽戏、纯粹无意义的出处。这条离奇的出处刚一入眼,一个联想便来到心头,并且一直在脑际萦回。我面对看似离奇、似乎挑战一切合理解释的东西——而其意味却又清清楚楚、并且被清晰地知觉到,只是在一个不同于常规逻辑的层面上——这已远不是头一回。常规逻辑,是我一生中几乎全部时间在意识层面运转的唯一逻辑——结果是:我不断地被“离奇”的、不可理解的事件甩在身后——它们那不可还原的离奇,令人焦虑!我的生活变化很大,自那一刻起(不到十年之前):我开始在我官能的一个更宽广的音区上生活。我懂得了:一切离奇、一切所谓的“无意义”都有\textbf{意味}——单单\textbf{知道}这一点\dmark{}\origpage{333},从此对无意义背后的意味保持好奇,就常常向我敞开其明摆着的意义。

在这“权复形”的无意义里,我感到一记与“反常层”的命名同类的\textbf{挑衅}\fnote{70}{(*)}{见“反常性”一注,n\textsuperscript{o}~76。}——此处的乐趣在于向自己证明:在一份有分量的杂志里、在一个自命为标准参考文献的文本里\fnote{71}{(**)}{而且,这个文本看来今天确实是一份标准参考文献——至少多年来,它一直是Zoghman的枕边书之一(他最近寄给了我)。正是在那里,他学得了可构造性观念(它在其定理里起本质作用);而且长年以来,他深信Verdier是这条对他至关重要的观念的神奇发明者。},人们\textbf{可以放胆}说一句与此相类的离奇话,而\textbf{没有一个人}会想到哪怕提出一个问题!而且我确信:这一挑衅里所含的赌注——文章发表八年以来的赌注——\textbf{赢了},直到今天:今天我是头一个向作者提出那个天真的问题的人。

当然,一桩离奇出现的时间(或地点)——在这里,恰是全文唯一一次指涉我个人的当口——绝非偶然;其形态也绝非偶然:这里借对一类观念——“权”——的影射(它与整篇文章的主题全然无关),并借对一个从未存在过的复合观念“权复形”的即兴,来实现。立时来到我心头的那个联想,很可能提供这桩离奇——在挑衅、权力示威之外——的更精确意味的钥匙。那是与另一处谜语式、同样纯礼节性的影射(但尚不带离奇的附加维度!)的联想——在注49开头所引的Deligne文章里\fnote{72}{(***)}{即“罐装的权与十二年的秘密”一注。从我们此处关心的角度对这篇Deligne文章作更周详的检视,见后文所引的“驱逐”一注(n\textsuperscript{o}~63)。}。那正是一处晦暗的影射,在一篇“权”一词严格缺席的文章里——除Serre与我,谁也无从看出——影射某些“权方面的考虑”,它们使我(文中特意写明:以一个欠一般的形式)猜出了该文的主结果。如我在更详尽的“驱逐”一注(n\textsuperscript{o}~63)里所解释的:在这份纯礼节性的影射背后,透出\textbf{藏掖}的意图——既藏我的角色,也藏那些他打算独占其收益的观念(关于“权”及其与一般上同调、尤其与Hodge上同调的关系)。这一意图,Verdier想必领会得格外透,因为他自己\dmark{}\origpage{334}就在同一个音高上“运转”(至少就他对我的关系而言;而依我看,这正是那对形影不离的朋友之间主要的黏合剂)。在两种情形里,诚实的做法本应是:开篇就明白指出主要观念、或促成文章的那个或那些问题的源头。

把这一点说清之后,这桩表面无意义的象征语言背后,我所知觉的意味便是:我可以毫无顾忌地当众陈列一桩明摆着的\textbf{无意义},同时借这个无意义、借对“权复形”那句荒谬的影射-出处,表达我真正的意图:我对Grothendieck在这项工作中的角色,无意透露分毫,一如Deligne当年对“权方面的\textbf{考虑}”的凭空影射——那条影射,对当时的读者,其意味不多一分,恰如我此刻发明的“\textbf{权}复形”,为着事情的需要、也随我的高兴!

我刚刚把昨天写的这条注誊清——其间我被Verdier的一个电话打断:那天我本想找他,正是为了当面问这个问题。我向他解释:我正在晚年勉力学一点上同调——这门学问我从来没弄懂过,他知道得很清楚——Mebkhout为供我学习,给了我一篇Verdier的旧文章;Mebkhout曾长期拿它当枕边书。我如今正尽力读它,但那里有这条谜语式的出处——他当然好意引我——可我全然不懂他指的是什么。

他非常高兴,甚至有点受宠若惊——可不是嘛,宽大的笑容从一副慈父般的欢容后面透出来——我竟然在我这把年纪,靠他那页旧纸学起上同调来了。我没有料到:当我说他分明知道我对上同调从未入门时,他会想到反驳我——那分明是久已有之的定论……至于那大名鼎鼎的“权复形”,我在听筒里再次感到他的宽大笑容(人们会说我编故事!),为终于有人(而且正是收信人本人)注意到那件被放过了这么久的事而欣喜。同时,又带着一丝窘迫——我想,那不是不知如何作答的窘迫,而是没能把一份乐趣\dmark{}\origpage{335}藏住的窘迫(一如人们对一桩略有荤味的故事所感到的乐趣……)。掉队如我,他实在无须在这一头费心!他不假思索地把话岔到Deligne身上(我没有提这个名字):Deligne在一篇文章里做过一个证明,其中还引了我——他记不太清在哪里;总之那里谈的是权,可不是嘛——当然他有点忘了——但那确实不是算术的权;这一点我全对,不是一回事……

声调欢容而不容置辩,他示意已经为我花了不少时间——略显匆忙,却不失那副从容、略带庇护的口吻。我为用这样一个有点愚蠢的问题打扰他表示歉意,并为他的解释道谢。我的歉意是真诚的,我的道谢也是——他确实把我想知道的一切都教给了我\fnote{73}{(*)}{即便摆出那副掉队的神气,我也并不真觉得自己在演戏(我没有这份天分);那全然自然——说实话,对这套我已近十五年不曾操持的玩意,我确实有点掉队!但我相信,即便老糊涂到该上灵车,我仍闻得出空核桃与满核桃的区别……}。

\section{IX.~我的学生们}

\subsection{沉默}

\notehead{84}{}

\dmark{}(5月9日)昨天我也许写得急了点:我在“好的参考文献”(见注82)里写道,作者兼前学生毫无顾忌地照抄的内容,“属于多少‘在局中’的人‘众所周知’的领域”。我曾试着为自己弄清,这些“多少在局中的人”究竟是谁——结论是:\textbf{不多不少,正是1965/66年间SGA~5那场研讨班的亲爱的听众们}——而且,如我已有机会说过的,这些听众多半或多或少掉队;而从这场研讨班的成文工作落在一群志愿者手里、我不肯觉察其信念不足的坎坷来判断,常常是“多”于“少”(当然,照例是Deligne例外)。事实上,只要SGA~5尚未成文出版、以便人们经由阅读而“入局”,就绝不会有别的“局中人”!这场研讨班实际上(天意弄人!)是在我最亲爱的两位学生与战友的两大“值得纪念的出版物”之后才出版的:其一,Verdier1976年的那篇文章——其中对由他执笔首次发表的观念之来源只字不提;那些观念,已被大量谈过\fnote{74}{(*)}{尤其见注n\textsuperscript{o}~67、67$'$、58、68$'$。};其二,Deligne随SGA~4½\origpage{338}。此后,人们才好心地邀请Illusie,来张罗其余部分的出版!

我已记不太清那场研讨班的与会者——比如Artin在不在场。我想,无论如何,我第一期的学生多半都在——当然Sinh~夫人与Saavedra除外(当时我尚未结识他们),或许还有Hakim夫人。此外还有Bucur(已故)、Houzel、Ferrand——我不算Serre——他对上同调的重型装备素来没有胃口,只是远远地、谨慎地偶尔露一脚。虽然除Deligne外,也许没有人十分清楚这一切将通向哪里,但我想,总该有十来位听众(不算很投入)至少跟得足够紧,可以被算作“在局中”\dmark{}。从昨天起在我脑际萦回的念头是:在所有这些“在局中”的人中间——他们因此俨然以上同调行家自居(即便不都像Illusie与Berthelot那样以“上同调”论文成名——那些论文分明够分量)——即便把Verdier与Deligne除外,想必仍有相当一批人,亲手捧读过Verdier这篇文章!Verdier身上某种神气,使我深信:从没有人向他暗示过,或许有什么地方不对劲。我也清楚:同样从没有人向我提过这件事——我是5月2日,即恰好一周前的今天,经由Mebkhout才知道这篇文章的存在的;而他对那场巧取,当然已知情多年。

这给十天前“全体一致的赞同”(为埋葬我寒微的身量而达成的)那句兴高采烈的确认,赋予一个十分具体的意味(注74)!这份赞同,把我“1970年以前”的学生中的许多人(乃至全体)——也就是今天在数学界引领风骚的许多人——都包括在内;它还把(或曾把)我的朋友Zoghman本人包括在内:他被当作上流社会的灰姑娘,逆着所有人死死攀住某种“对我的著作的忠诚”(借用他自己的说法\fnote{75}{(*)}{(6月7日)最近来访时通读关于葬礼的各注,Zoghman向我指出:他那个“对我的著作的忠诚”的说法,并不真正表达他的想法。他心里想的,毋宁是对自己的判断力与数学本能的信任——它们向他表明,我的著作为他带来了他所需要的某些想法。因此,那是一种对\textbf{他自己}的忠诚;而要做真正革新的工作,这确实是本质的东西。}),并有胆量与执拗,时不时以此自命,后果则有目共睹。谁能懂,就去懂吧!

总之,我暗示那份有分量的杂志刊登了一篇照抄“众所周知”的空壳文章,是我错了。作者当着(虽非所有、却)众多见证者的面照抄的东西,既未发表、也非“众所周知”(唯有一个闭链在凝聚框架里的上同调类除外——那个我久已发表);而那些观念,我再没脸去贬低:我并不曾认为,花一年工夫、在众多听众面前在一个研讨班里发展这些观念与其他观念,是虚掷光阴。想必Verdier的文章,是我所发展的观念与技术的一小部分的“摘要”,有用而做得好:恰好使它们得以进入“众所周知”的领域,成为把上同调(或同调)用于那些多少配得上“簇”之名的对象者的日课。\dmark{}\origpage{339}从这个观点看,Verdier做了有用的事\fnote{76}{(*)}{他做了,诚然,以原初SGA~5研讨班的“拆解”为代价——那场拆解,他与Deligne是主要的操办者与“受益者”。\par (6月7日)三天后(5月12日)的反思(见“屠杀”一注,n\textsuperscript{o}~87)显明:与这场拆解相比,所发生的更是一场“屠杀”,而Illusie的关联比Verdier还要直接——即便他并非“受益者”,而是代人行事。},我最终并无不满之处。然而,根据我从我的前学生、今天仍是朋友的人那里——在电话里——以及从他这个人身上我得以感到的许多别的事情(其中最“重”的、或至少最“轰动”的,是反常层会议之欺瞒)所感到的——我分明感到\textbf{有什么地方不对劲}。那场值得纪念的研讨会,从数学上说,想必在许多方面都极其辉煌。而“不对劲”之处,处在全然另一个层面。我可以试着以言词勾勒它,但我分明感到,那没有多大意义。在那个不感到这场研讨会不对劲的人那里——也在许多别的研讨会里,毫无欺瞒——当我做了这场“勾勒”的努力、哪怕完全成功地做到令我满意——他也不会多感到一丁点……

对我仍然悬而未决的问题是:这个“征兆”——它表现为一件如今想必相对寻常的琐事(一位作者把他人的未发表观念呈献为自己的)——这个征兆,是世风普遍败坏的征兆,因而只是今天数学界某种“时代精神”的典型征兆呢,还是它更想就我个人、就那个曾经是我、如今又回到我身上的人(经由昔日学生们对我的态度)教给我点什么。

两种可能的意义绝不互相排斥。倘若没有某种世风的状况怂恿,我的前学生们对我的关系,本无从循这条路径表达。况且,早在这个“征兆”之前,我已看到别的许多征兆;在我看来,它们在“风俗画”的层面上更雄辩。而在这个征兆里打动我的,是它区别于所有其他征兆的特殊之处:它似乎\textbf{同时牵连我的大部分昔日学生}。

这样的情势绝非偶然。把它干脆记在“世风败坏”(千真万确的败坏)的账上,便是回避其更个人的意味——那意味牵连我,一如牵连我的每一位前学生。倘若我说“每一位”,似乎超出了这个征兆的实际幅度,那是字斟句酌。因为这个征兆适时地提醒我:我的某一位昔日学生,竟不曾至少面对过这类局面,几乎是不可设想的。多年来,我感到一种关于我的某种“风向”,吹拂在我离开的数学家世界里(如今我分明看清这风的来源与缘由,我想)。不可能他们中的任何一位,从不曾感到这阵风的吹拂——无论是经由像这篇掘墓文章的发表这样的“事故”,还是经由任何别的场合。无论当事人愿意与否,这样一场遭逢,必然向他提出(或重新提出)他与我的关系的问题——我曾教他手艺。而我确认的征兆,在把我引到这件事的那一个之外,是:\textbf{那些曾是我的学生的人里,没有一个就此事给过我任何回音}\fnote{77}{(*)}{(5月31日)有意思的是,唯一一个曾向我暗示存在一场下葬的人,是一位非洲朋友——十年前,他跟我做过一篇第三阶段博士论文(即“1970年以后的学生”,身份寒微);我一直与他保持友好的交往。他有所暗示的那封信,应当是两三年前;在那个时刻,这毫无使我惊讶之处。我当时没有就他的印象问细节;他只是最近,才重新谈起它们。}。这是一桩“巧合”,其意味仍向我滑脱——但它不可能没有意味(注84$_{1}$)。

天开始破晓——我感到是停笔的时候了。我拿不准:在《收获与播种》里,此刻此地继续深挖这桩触目巧合的意味,是否合宜。这也许是留给别的明天的收成——倘若我今夜的反思,在我昔日学生中的某一位那里遇到回音。($\Rightarrow$注85)

\notehead{84$_{1}$}{}

(5月16日)我的昔日学生们之间的这份完美的协调——在对我的全然沉默之中——与其他征兆指向同一方向。其一是迎接“外来者”情节(见第24节)的全然同样的沉默——关于那份沉默,我在注n\textsuperscript{o}~23v里已多少追问过。另一方面,除Berthelot(他给我寄过许多抽印本)、Deligne(寄过四份,在其约五十项出版物里)与Illusie的一份之外,我的任何一位前学生的抽印本,我都没有收到过。关于他们对我之关系的两价性,这说明了许多。寄送抽印本——即便人们怀疑我是否会在自己的工作中使用它们\fnote{78}{(*)}{(5月31日)直到1976年,这甚至想必可以算是被排除的:七十年代初,我曾相当明确地说过,我不打算再重操数学。1976年在IHES的那场关于带除幂的De~Rham复形的报告,则相当清楚地表明我仍在关注数学。}——本会是最明显不过的方式,向那位教了他们手艺的人表明:这门手艺在他们手里并未僵死,而是活生生、在运转。但同样真确的是:至少对其中某些人,他们的出版物也见证了他们参与一场心照不宣的下葬;对那场下葬,最好还是别告知被预葬者——手艺也罢,非手艺也罢……相反,我却收到过好几位从事晶体上同调的作者的许多抽印本\fnote{79}{(**)}{(5月31日)说的是我不认识的年轻作者;我料想他们效法了Berthelot——在他们眼里,他想必是长辈。这里稍显奇怪的是:至少两年来(自1982年9月6--10日的吕米尼研讨会起),Berthelot正亲力亲为地埋葬我(对此可参见“共同继承人……”一注(n\textsuperscript{o}~91)5月22日的脚注)——莫非他对我的关系近来起了转折?我不记得收到过那篇晶体上同调及其盟友的综述文章的抽印本——文中对我的名字讳莫如深;他肯定小心没寄给我!},甚至收到过许多我只闻其名的分析家同行们的抽印本——当他们的工作重提(有时解决)我在三十年前或更早提出的问题之时;尽管明摆着我不会回到那个我已告别的题目,而且从“实用”的角度看,那些抽印本是白寄了。但这些同行想必感到了某种我的学生们不屑于感到的东西。——当然,在六十年代,我的学生们在我的一切出版物里都是头等的受赠者:无论我的文章,还是EGA与SGA两大系列;他们每一位(除Sinh~夫人,或许还有Saavedra)都应当拥有我1955至1970年间出版的全部著作(想必近一万页)。

诚然,我的前学生们有很好的同伴:我在“上流”数学界的任何一位旧友——包括那些著作与我的著作血肉相连的人、或在六十年代我的工作纲领的展开中扮演过角色的人——都不曾认为值得在我离开共同环境之后继续给我寄抽印本\fnote{80}{(***)}{(5月31日)当然,促使他们给我寄的心理缘由,远不及对我的学生们那样强;但,人们也许天真地想,总该比我的分析家同行们、甚至比我收到过抽印本的众多代数几何学家——我并不认识或只略识面——强些。分明,自我离开共同环境之后,“曾是朋友”这件事,在我数学界的旧友们身上造成或强化了我已有机会观察到的那种拒斥的自动机制。(关于这些态度——《收获与播种》里此处彼处顺带影射过——见5月24日的“掘墓人——或全体教众”一注,n\textsuperscript{o}~97.)}。最近,在我寄出《纲领概要》的十五二十位旧友(包括几位学生)之中——那份概要(除别的以外)向他们宣告:中断十四年之后,紧张的研究活动重新开始,且研究主题与我们当年共同推进的血脉相连——只有两位(Malgrange与Demazure)费心寄来几行致谢。我所收到的较周详(而且热忱)的些许回音,来自我近来结识的青年数学家,以及我的老朋友Nico~Kuiper——尽管他对我所做的这一类事情全然外行。他是经第三者得知这份文本的,并为我那不期然的“重返”由衷高兴\fnote{81}{(*)}{(5月31日)这几乎是唯一的、来自我的一位旧友(或前学生)的、对“重返”表示首肯的回音。这当然毫无奇怪之处:亡者现身,总归不合时宜地打断葬礼的顺畅进行……\par (6月17日)然而最近,我荣幸地收到Mumford一封热忱的信:他自称对《概要》里勾勒的想法感到“激动不已”(thrilled)且“非常兴奋”(very excited),并向我确认:我为其Teichmüller塔的组合描述所需要的那个关键技术结果,确实已得到证明。这是自1978年以来,头一回有我的一位旧朋友对我的“阿那贝尔式”想法产生共鸣——那想法的非凡分量(堪比Motif瑜伽),在我看来,自始便是明摆着的。\par (1985年3月28日)在这些文字写成之后,我还收到 I.M.~Gelfand一封非常热忱的信(注1984年9月3日),回复《概要》。}。

\subsection{团结}

\notehead{85}{}

(5月11日)那场倒霉的SGA~5研讨班的事,继续在我脑际萦回。“好的参考文献”\fnote{82}{(**)}{见注n\textsuperscript{o}~82。}分明给这段往事投下一束新的光,并同时赋予那场光辉的“SGA~4½之役”一个新的意味。

越想,SGA~5的事在我眼里分量越\textbf{重}。我的第一印象——那是仅仅几个星期前我刚“下船”时(见注68、68$'$)——是:65/66年那场研讨班的穷听众中间的一盘散沙,被我的朋友皮埃尔以其大名鼎鼎的战役加以利用;而在那场战役里,别人都无关。至于SGA~5的诸种不幸,既不怪他也不怪任何人,而宁可怪我自己:我没能——唉!——点燃我的志愿者-执笔者听众的热情,也没能替他们做那项他们执拗着不做、嘴上却说着“马上就动手”\dmark{}\origpage{343}的工作。然而近几天显明:毕竟有一位,其热情十年后才苏醒,把他在那里乐于取用的东西发表出来(对研讨班不置一词),从而在别的“志愿者”们还没下决心动手的时刻,为自己立了一份“好的参考文献”。

自昨天起,对我越来越清楚的,是:直接卷入这场研讨班之藏掖的,不止两个“恶人”,而是\textbf{我的每一位“上同调学家”学生}。如无差错,他们每一位都听过这场研讨班——即(按我的“上同调学家”学生们出现的先后为序):Verdier、Berthelot、Illusie、Deligne、Jouanolou。(Jean~Giraud不算;他在与SGA~5或其前驱SGA~4的主题相当不同的音区上工作。)

这场研讨班,是我首先\textbf{为了我的学生们的收益}而做的——即便有时他们哀求饶命——\textbf{我认为它不是一堆狗屎}。那一年里,他们每一位都学到了一大包“使用上同调的数学家”这门手艺!我在平展框架里、以周详得多的方式,重做我最初在凝聚框架里发展的那些观念——这些东西,除了这场专为他们的收益而办的研讨班,他们在别处无处可寻:在我之前,没有人费心做过;除我之外,没有人甚至感到有什么要做、为什么做。(照例,Deligne除外:他在这个研讨班本身里经月习得,比别人悟性快。)正是听过这场研讨班(及其前一场)并在自己身上多少消化了它,而没有别的,才使他们从此对对偶形式主义“在局中”——而且他们是\textbf{仅有的}在局者。这份\textbf{特权},依我看,为他们造成一份\textbf{义务}:即留意使这份特权不致滞留在他们独家手里;使那些从我嘴里学到、并在他们直到今天的全部后续工作中构成不可或缺行装的东西,在合理而惯常的期限内——至多一年,宽限两年——交到人人手中。

\dmark{}\origpage{344}人们会不无道理地说:首先轮到我来督察这件事。但我当年是真心接受学生及其他听众自荐承担成文工作的(对肯认真投入的人,成文只会带给他们最大的益处)——并非为了享受一项本属于我的工作由别人代劳、而我袖手旁观的好处。我继续(有Dieudonné的协助与其他人——其中1966/67年间还有Berthelot与Illusie)发展另一些在我看来同样紧迫的奠基文本,而那些,当年没有别人能代我做、或在没有我的协助下做\fnote{83}{(*)}{1960到1970年间,我必须以平均每年约一千页文本(EGA、SGA、文章)的节奏工作,其中几乎全部都成了通行参考文献(这一点,我在写它们时、或鼓励某位合作者在我的协助下写它们时,就已经很清楚)。}。这些文本本身,也成了不可或缺的参考文献——包括对我的“上同调学家学生们”:他们同所有人一样,很高兴在需要时找到现成的它们。

凭着经由与我的接触、经由听我的或参与我的研讨班而获得的对上同调观念与技术的驾驭,合力成文这场研讨班,与对“大名鼎鼎的数学共同体”所提供的服务相比——也许还有,稍后,对我的一份忠诚义务相比——是一件微不足道的小事。我说过,对我这个有手感的人,写定整场研讨班是几个月的活。五人分工,加上他们在那些年里各自积累的成文经验,手头又有我的详细手写笔记,每人的投入至多一两个月。他们比别的执笔者(比如Bucur——他分明巴不得把一件明显超出他能力的差事,交给更年轻、动机更直接的手)装备好得多。

只要我还在场(也就是其后的三年里),我理解“靠我”的反射可能起作用:理应由我统筹一切、跟“志愿者们”周旋。很可能,倘若我当时要求他们每人限期做两三场报告,我同样承担,以便了结此事,他们不至于推却。\dmark{}\origpage{345}自我退出数学世界的那一刻起,局面全然改观。他们发现自己成了\textbf{某份遗产的独家保管人}:既是默示的(没有遗嘱),又非常具体。诚然,从实际看,我的离去等于一场\textbf{消失}——我确实“亡故”了,意思是:除他们之外,无人知道这份遗产,无人能用它,无人(无论为善还是为恶……)关心它的命运。

倘若在我离去后的七年里,这份遗产始终隐没(除1976年的“好的参考文献”!),那是因为\textbf{我的学生们不曾在整段时光里让它变成公共的}。按比例说,这局面在我看来,与“Motif瑜伽”的局面相去不远:那部瑜伽,除我之外,唯有Deligne洞悉入微,而他判定把它留在自己怀里、独享其益。若说乍看有区别,那是:此处的“受益人”是五个而非一个;而且一个人所藏之物的深度,与五人共同藏起之物的深度,不可同日而语。

每个人的深层动机,我当然不知道——即便是Deligne的情形,我的领会仍模糊,而且想必仍将模糊。但在“实际”的层面上,Deligne的棋(连同“SGA~4½之役”以及其余)再清楚不过。同样清楚的,是这些战役\textbf{没有全体的团结便无以进行}。依我看,Jouanolou却不太在局中——他不俨然“权威”,我的印象是他久已离开上同调的泥沼(注85$_{1}$)。但我难以想象,Illusie与Berthelot不曾亲手捧读过SGA~4½与“好的参考文献”两者;他们跟我一样识字,也不比我笨。

倘若Illusie突然张罗起SGA~5的出版——恰在Verdier取用、Deligne取用、Deligne需要一座后勤基地以成其大名鼎鼎的SGA~4½(从中顺手牵走孕育了这份文本与其全部著作的那两场研讨班)之时——而Illusie手里已有十年功夫可以做这件事,这肯定不是偶然。倘若我1966年所做的收官报告——论未决问题与猜想——“不幸没有成文,其非常漂亮的开篇报告也没有[sic]:后者巡视了各语境(拓扑的、复解析的、代数的)中的Euler--Poincaré公式与Lefschetz公式”,这同样肯定不是偶然——而那是一场我不在行的下葬。同样不是偶然:把研讨班的一场关键报告截肢、不经任何手续送进SGA~4½,在Illusie与Deligne看来竟如此自然(而且只配在“细节变更”里顺带一提)\fnote{84}{(*)}{(5月16日)实际上,正如次日我就发现的(见注n\textsuperscript{o}~87):对母研讨班(或父研讨班!)SGA~5,在Verdier、Deligne与Illusie的手里,发生了一场真正的“屠杀”。}。

Luc~Illusie——我对他怀有如对Pierre~Deligne一样的感情,而他(一如Deligne)在我面前始终极其亲切\fnote{85}{(**)}{甚至在我1970年离去之后,Illusie对我仍有过细致的体贴——比如,长年以来,岁末节庆之际,他都寄给我非常漂亮的贺年卡。恐怕我不常回信致谢、报个平安——这些忠实的友情的信使,来自一个显得无限遥远的过去,而我已经与它失联。\par (5月16日)相反,Illusie从无继续或重建数学层面接触的念头;去年(我为数学问题联络他时),我感到他的迟疑。自我离去以来的十四年里,我只收到过他一份、唯一的一份抽印本,日期1979年。}——他的(有意识的与无意识的)意图,我不知道。但我确认:他在Deligne身边,成了一场\textbf{毫不脸红的欺瞒}的共同演员:那场欺瞒,把1965/66年的母研讨班SGA~5——恰恰是Deligne头一回听说概形、平展上同调、对偶及其他“枝蔓”的地方——说成一份顶着欺眼之名SGA~4½的文本集(晚八年写成,却佯装在前:既凭标题里的编号,也凭Lecture~Notes里的出版号,更凭作者那句不寻常的评语“它的存在(SGA~4½的)不久将使人得以\textbf{照原样}出版SGA~5”——着重号是我加的)的一种无形而近乎滑稽的附录——并且还摆出对那份寒碜的合集所自出的全部工作的毫不掩饰的轻蔑。

没有这些被如此潇洒地轻慢的工作,奠定其应得名望的Deligne的大部头著作,眼下没有一部写得出来——倘若真要问,一百年后也写不出来(Illusie与我的其他上同调学家学生们想必同样)。在这场“SGA~4½之役”的精神里,有一种\textbf{无耻}\dmark{}\origpage{347};Illusie(想必连自己都不觉察)为它作保;而它之所以能如此招摇,只因有一个\textbf{共识}的心照不宣的首肯。除Deligne本人之外,卷入这一共识的头一批人,恰恰曾是我的学生、某份遗产的主要受益人:那份遗产,在他们眼前,被交付给抢夺集市的风险与轻蔑。

还有那份武断的自满神气,那份慈父般的庇护神气——前天在我们通电话时,我还在我的前学生身上领教过\fnote{86}{(*)}{关于这次通话,见“玩笑——或‘权复形’”一注(n\textsuperscript{o}~83)。};还有我的朋友皮埃尔身上那份较隐约的居高临下——自“SGA~4½--SGA~5”双重大捷的翌日起我就领教过(而当时,乃至其后七年,我对它连一丝疑心都没有)——这些神气,绝不是孤独的产物,而分明仍是一个\textbf{从未被质疑过的共识}的征兆。这些神气,告诉我的不仅是Verdier与Deligne,还有我曾经的学生们全体;而首当其冲,是那些(就其工作主题与日常操持的工具而言)最直接相关的人。

“欺瞒”这个词——它不召自来——适时地使我想起另一场欺瞒,其中摊开着同样的犬儒:即所谓“反常”研讨会的那一场。如今,这两场在我眼前\textbf{血肉相连、不可分解——使两者皆成为可能的是同一个精神}。也许除Jouanolou(他已不太混迹“上流社会”)之外,我认为这些前学生是这场耻辱的共同责任人、团结一致的上同调学家。至于Berthelot与Illusie,没有任何东西允许我预设恶意或不诚(在Verdier与Deligne的个案里,那是不容置疑的)。但我至少确认一种盲目、一种健全官能之使用的堵塞;其深层缘由当然向我滑脱。倘若他们心中没有那份冷漠与轻蔑的蓄意,那么Zoghman~Mebkhout——七十年代唯一公开自称出自我的著作的人,而且题材恰与两人切身相关(尽管他们不屑于觉察)——本会享有最低限度的“善意推定”,使他们至少稍微看一看他所做的事情,从而觉察他自1974年起所投入的方向的趣味;那趣味是\textbf{明摆着的}\dmark{}\origpage{348}!两人都不屑于从一个装作还出自Grothendieck的无名小卒那里觉察任何东西。无名小卒的论文,是他亲手寄给他们的;我不知道他们开没开封,或读没读那些更短、更易消化的说明文字——总之,他们不屑于哪怕寄一张收条(Deligne亦然——分明是他定调子的\fnote{87}{(*)}{(6月12日)其间我得知:两人都不曾参加那场研讨会(1981年6月的吕米尼)。不过,可参见“欺瞒”一注,n\textsuperscript{o}~85$'$。})。

这当然不曾妨碍:与那场值得纪念研讨会的其他与会者一起,他们带着兴趣了解了那条卓越的“Riemann--Hilbert对应”,却想不到就其来源或父权提出丝毫问题,或至少(作为扎实的数学家)问一问它在哪里得到证明(注85$'$)。不过,在这一点上,我信得过Deligne:他自会乐于优雅地向他们讲解那份证明——对像他们那样的人,想必明白之至——恰恰是那种借Hironaka式奇点消解的证明,他们久已学得,而且不是从别人、正是从我这里(注85$_{2}$)。Riemann--Hilbert,Hironaka,阿布拉卡达布拉——戏法告成!

显然,一如Verdier、一如Deligne,他们已全然忘了何为\textbf{数学创造}:一个经月历年缓缓澄清的视野,揭示那件“明摆着”却无人看见的东西,在一个无人想到过的“明摆着”的陈述中成形(而在本案里,一位Deligne曾徒劳地试过整整一年……)——随后随便什么人都可以用现成的技术在五分钟内证明它,靠着他在某场遥远研讨班的长凳上有幸学得、如今却不屑(或不记得)想起的技术……

倘若我谈论Berthelot与Illusie时毫不容情,那并不是我尤其想给他们抹黑(在与他们的两位朋友清账之后)。我知道:他们既不比他们亲爱的同行们“更坏”、更蠢,也不比我更坏更蠢;而我在他们身上、在这一场合所察觉的那份嗅觉与健全判断的欠缺(有时,还有那份对他人必要的尊重的欠缺……)绝非痼疾,而是一个\textbf{选择}的效果。无疑,这一选择曾给他们带来中意的\textbf{收益}\dmark{}\origpage{349};而随着我的反思而来的另一份“收益”,对这一位或那一位,也许不合时宜。倘若如此,那只说明:它再度重演同一个选择——即在自身官能的微不足道的一小部分上运转的选择,不惜把鱼目当珍珠、把珍珠当鱼目,并无望地把空核桃(好哥们的)与满核桃(无名外人的)混为一谈。各人自便!($\Rightarrow$注86、87)

\notehead{85$_{1}$}{}

Jouanolou是我学生中除Verdier外唯一一个不曾把发表其论文放在心上的人。在我看来,这是对他所发展的奠基工作——即导出范畴观点下的$\chi$-进上同调——冷淡的征兆。由于他这一主题的工作,大部分是在我离去\textbf{之后}进行的——即在同学们以Deligne与Verdier为首、对我引入同调代数的观念(尤其导出范畴观念)之普遍冷淡已然亮出信号的时刻——气候并不鼓励Jouanolou认同于自己的工作、给它(应得的)发表之荣。而当这同一些Deligne与Verdier,追随着Zoghman~Mebkhout(别名Verdier的无名学生、别名Grothendieck的遗孤学生)的著作,终于(在大张旗鼓、互相吹捧之中)发现导出范畴的重要性时(见注n\textsuperscript{o}~75、77、81),那份遭冷落的Jouanolou论文,自“反常”研讨会起,重新恢复了全部的现实性;倘若1970年我离去之后,概形上同调理论的发展照常进行,那份现实性本不曾中断。一个触目的细节,例证Deligne在我离去之后其取舍的一场急遽“转向”:恰恰是Deligne本人(他分明深知在三角范畴框架里发展$\chi$-进上同调形式主义的重要性)向Jouanolou提供了一条技术关键想法,供一个有待研究的、成形的$\chi$-adic三角范畴的定义之用;那条想法在论文中得到发展。(对此可参见我1969年论Deligne工作的“报告”,第8节。)

(5月30日)关于Jouanolou的工作,另见“共同继承人……”一注,n\textsuperscript{o}~91。

\notehead{85$_{2}$}{}

一份有意味的“巧合”:恰恰是在同一场SGA~5研讨班里,这一整批人学得了那条证明原理——它既用于证明平展上同调里的双对偶定理(在拥有奇点消解的情形),也用于诸$R^if_*$无真性假设的\textbf{有限性}\dmark{}\origpage{350}定理,对$R\mathcal{H}om$、$Lf^!$亦然。(这些有限性定理,同样被从SGA~5的已刊版本里藏掖走、并进了SGA~4½;Illusie甚至不认为有必要在其导言里点明——我是写这些文字时才发觉的!)Zoghman没有听过那场研讨班的福气(他得到的是“好的参考文献”的替代品),是在另一处我使用过该办法的地方(为$\mathbb{C}$上光滑概形的De~Rham定理)学得它的。他其实也能在“好的参考文献”里学到:在那里,我的证明在解析框架里被照抄,以确立我的SGA~5学生们从此乐于称作“Verdier对偶”的东西(那条对偶,在我结识他的荣幸之前,我就已经知道了)。分明,一切都环环相扣!\textbf{同一份证明}(连同陈述,一并抄自我)被Verdier用作一种对偶的父权凭证——那种对偶,他除了在这场被肢解、被交付轻蔑的SGA~5研讨班里,无处学得;而它又被用来\textbf{对付}Mebkhout,成为(正因其“明摆着”)心照不宣的借口与手段,毫不脸红地巧取一项重要发现的信誉。

(5月30日)我记得,我头一回使用Hironaka式奇点消解、并领会消解作为证明工具的非凡力量,是为了对Grauert--Remmert定理做“三勺汤”式的证明——该定理描述某类复解析空间的有限覆盖上的复解析结构;以及在$\mathbb{C}$上有限型概形的情形里的类似陈述。(并不排除:恰恰在那次,原理是Serre向我提示的。)后一结果,是证明平展上同调与寻常上同调之比较定理的主要配料(其余,仰仗$Rf_!$形式主义归约;再稍加消解,从$Rf_!$过渡到$Rf_*$……)。

\subsection{欺瞒}

\notehead{85$'$}{}

(6月3日)实际上,我得知:他们本无须为这份父权提出问题,因为Berthelot与Illusie都是从Mebkhout嘴里学得好上帝定理的——前者在1982年2月,后者早在1979年(Mebkhout论文答辩之年)。虽然两人都不曾参加那场研讨会,他们却是那场研讨会之欺瞒的同谋,因为\dmark{}\origpage{351}他们不可能不知道那场围绕Mebkhout对好上帝定理尤其之父权的藏掖。我还可以想象:他们与研讨会的全体与会者一起,争先恐后地做了那场集体欺瞒的头一批受骗者——那场欺瞒,由他们的朋友Verdier与Deligne悉心操办(我的五位上同调学家学生中,有四位俨然是它的同谋)。至少就Illusie而言:去年夏天Mebkhout来我家之后,与Illusie通电话时,他对Mebkhout显然满不在乎,令我吃惊——他十分惊讶(几乎为他老师难过:他本指望老师有更好的判断力……)看到我在代数簇上同调理论的重启中,把头等角色派给Mebkhout。一些力量可观的共识,早已决定把Mebkhout归入无名小卒之列;我的朋友Illusie与这份三重矛盾活得逍遥自在,不问一个为什么:好上帝定理及其哲学的头等角色;围绕这些事情之父系的藏掖(他本人也在大帮人之中参与其中);以及他对Mebkhout的器量与角色的轻视——尽管他分明知道:Mebkhout是这些东西从未被点名的作者,而这些东西革新了数学的一个领域——在那个领域里,Illusie本人俨然权威。

我在这里重新找到那种健全意识与健全判断的完全堵塞——甚至在一桩表面上如对科学问题的评判般不涉私人的事情里;这种堵塞我已有不止一次机会影射,而每一次都重新令我错愕。而我在Illusie(想必还有许多别人)对Mebkhout——我的“遗孤学生”——的关系里确认的这份矛盾,肯定不是别的,正是他与我的关系中那桩更关键的矛盾的众多效果之一。正是这份矛盾——在他身上尤其,也在我的其他学生身上——在本书这支由我昔日学生们组成的葬礼行列的各注所进行的反思中,越来越清楚地显形……

\subsection{亡故者}

\notehead{86}{}

(5月11日)一如常有的情形,我是带着几分迟疑开始这场新反思的,主题“SGA~5--SGA~4½--反常性”——它似乎已被检视再检视到腻:“这会给一位早已听腻的读者留下糟糕的印象;再进这样的细节毫无雅致可言,SGA~5呀SGA~4½呀,这些都是过去了的事,不值得再饶舌……”

\dmark{}幸好我没有被这类熟悉的调门吓退——它想以“分明不值得”、“由它去吧”为口实,阻止我探一件事的底(至少就我当下所能探到的深度)。\origpage{352}倘若我曾发现过我认为有用而重要的东西,那从来是在这样的时刻:我懂得不听那份冒充“理性”、甚至“体面”的声音,而顺从我身上那份不体面的欲望——去看看甚至被认为“没意思”、寒碜,乃至粗劣、不体面的东西。我一生中不记得有哪一次,我曾后悔把某样东西看得更近一点——对着那些恨不得阻止我的积习反射。在《收获与播种》里,这些抑制反射比别的场合更强,因为这场反思注定要公之于众,这就立刻带来分寸(当我牵连第三方)与简洁(顾念读者)的约束。但我最终并不觉得,这些约束曾有任何一刻阻止我触及我想触及的东西、或把它深入到我感到欲望的地步。在曾一度显得棘手的临界情形里,我径直前行,带着这样的把握:必要时,我总还留着不把那场不客气的反思所“流出”的东西收进《收获与播种》的余裕。这些“边缘个案”,只在我犹豫是否牵连他人时出现;牵连我自己,则从未犹豫。但即使在第一种情形,事情竟是(这对我是一桩意外):我从不曾动用过那份“余裕”:《收获与播种》的文本,是我反思的全本——至少,是这场反思中找到通往书写之路的那一部分。

我感到:随着前一条注\fnote{88}{(*)}{即同日的“团结”一注,n\textsuperscript{o}~85。}那场简短的反思,局面大大澄清了。我是说:一个曾经故意弄得混乱的局面的某个本质侧面——我刚以“主题”的三重名目(SGA~5--SGA~4½--反常性)唤起过它——以全光显现在我眼前:即一种“团结”、一种“默契”;它们此前只被模糊地知觉过。这绝不意味着我自以为探明并理解了一个复杂局面的全部机关与来龙去脉——它至少直接而分明显然地牵连七个人:Zoghman~Mebkhout(在某种意义上充当某一局面的“显影剂”)、我的五位前上同调学家学生\dmark{}\origpage{353},以及我自己。我甚至不敢自称:对“SGA~5等等”这一局面里、在我自己身上起过作用的全部机关与动机——那场“倒霉的研讨班”距今已近二十年!——我都有所知觉。但我感到,比起昨天(或甚至今天早晨),我的状态好得多,足以理解并安放我希望将从这个或那个主要当事人那里到来的回声。

摆在我面前的主要问题(它想必在反思的另一个阶段已在场,如今带着新的力量重现)是(依我看)这个:经由我的学生们(差不多全员)进行的这场葬礼所发生的一切,是一件全然\textbf{非典型}的事,系于我的个人与我的独特命运的某些特点(如近十五年前我退出数学舞台,及围绕它的种种情势,等等)?抑或相反,是一件“再自然不过”的事,出于机缘的简单巧合——按“机会造就窃贼”的原理?我犹豫相信后者,尽管此刻辨不出、甚至隐约瞥不见:我这个人身上的哪个特殊方面,竟有这种本领,在我昔日学生们中间造就一份如此完美、如此一致的\textbf{赞同},去埋葬“老师”,并埋葬那些自称出自他、或著作分明带着他的印记(却并不因而是“他们的”)的人。是我这个人周围那圈“父亲”的光环吗——我已有机会谈过?抑或,我的离去这一件事本身,对他们的每一位置疑?此刻,我全然说不出来,因为缺乏那双会看的眼睛……也许到来的月份会就此教给我点什么\fnote{89}{(*)}{(5月30日)这个方向上的反思,见“掘墓人——或全体教众”一注,n\textsuperscript{o}~97。}。

最近三个星期里,我不止一次想到另一桩奇怪的“巧合”:葬礼“在其全部光辉中”的发现(经历LN~900--SGA~4½--SGA~5--反常层会议四幕,然后回到SGA~5与SGA~4½)——这项发现恰在此时完成:我刚结束一场对我的数学家过去、我与学生们之关系的深入反思。也就是说:恰在我刚刚——尽我官能之所能、并在当时已知事实的限度内,照它们经由常常雾蒙蒙的回忆所还原的样子——就这段过去“与自己算清账”的时刻\dmark{}\origpage{354}。换个说法:那恰是我终于\textbf{准备好}得知这件事、并从中获益的时刻。

“偶然”把事情安排得如此之好,以至禅观中连一个中断都没有。那场反思——它以对那些(依我所感)最重要的观念之命运的一场简短回顾开始\fnote{90}{(*)}{见3月31日的“我的孤儿们”与“拒受遗产——或矛盾的代价”两注(n\textsuperscript{o}~46、47)。}(那场反思还停留在某种模糊之中,只有某种基调执拗地浮出……)——这场反思在4月19日那个星期四,全然自然地继续下去。诚然,那还是带着“无耻”——借用方才的词,它也恰当地描摹了我当时感到的一桩事——的印象所激起的情绪:读那部“值得纪念的文集”LN~900之时。

在这“同一场”反思的新起点上,主发动机是“老板”——我的自尊、我的体面之感受了伤;写下我的情绪,我便在某种程度上从中解脱。分明是“自我”、“老板”在随后十天里领舞——那些日子,没有笑容也没有笑声,只有密不透风的严肃。无疑我必须经过那里、经过那十天的弯路,反思才回到它所离开的中心——我自己的身上。我仍记得那次回归的如释重负——像出了隧道重见天光!就在那时,我找回了笑与微笑,仿佛从未分离。那是4月29日。次日30日,当月最后一天,我欢欢喜喜地为反思的这一最终阶段画上句点。

那肯定也是我完全准备好接收下一个“包裹”的时刻——这一回由我的朋友Zoghman惠寄——“研讨会”包裹,两天后到达。到今天,我消化那份包裹的实质已进入第十天。而在这一阶段,尽管我为这场没完没了反弹的反弹而咬着嚼子,笑容却一天也不曾离我而去。而今天,我真的相信(诚然,已是第无数回!),是终于画句点的日子。

五天前,我就有过这同样的抵达终点之感,觉得剩下的只是总务活:这里那里补几条脚注,把涂改过多的页誊清\dmark{}\origpage{355}(那每一回都是思路尚欠清晰的征兆,它要求经由这看似机械的功夫各就其位;而文本经过誊清,总是带着一副新面目出来……)。那是在我刚写完如今那则“我的朋友们”注(n\textsuperscript{o}~79)之时;它自然而然地续成了几段“终曲”。但我终于把这些终曲与注的开头分开了。因为事实证明:那桩大名鼎鼎的总务活爆了:那些不打行距敲出的“脚注”,变成了真正的(而\textbf{非}脚注的)漂亮条注,颇具规模,不得不隔行重新打出,然后想法安置在此处或彼处。又过了几天,我才向明摆着的事实低头:继名为“研讨会”的行列之后,另一个行列正在成形,要加入送葬的队伍——而压轴的行列,将不是我心中已定的那场研讨会,而将由\textbf{学生}领队。而且就在今天,当第一个行列——已缩成一条注——刚添上第二条(“一种不义与无力的感觉”)时,我也知道了将由谁领队:是“遗孤学生”。于是,这支由一个小写的“学生”(遗孤的,合乎其寒微的身份)开道、又由一个大写的“学生”(这回毫不寒微)压阵的队伍,我看,终于齐备了!

现在,我想,是时候了——在第一次“假到站”之后——回到最后一曲《自深处》(De Profundis)的和声;它们今天比五天前更贴切。以下是我当时记下的;它们同样表达我此刻的心情。

(5月31日)最终,这又是又一次“假到站”——“终曲”这一回还是提前了!二十天过去,其间“总务活”不断爆成对被忽视的这一侧面或那一侧面的反思的重启。又有六条注加入本该压阵的“学生”行列。灵车在“学生”身后现身,载着四口棺材,掘墓人随行。分明还缺些东西,好给一支似乎谁也没在送葬的队伍以实体与意味。

经验使我变得谨慎;我等着事件到来,此刻不冒险预言:队伍是否终于齐备,或者是否还有一个被遗忘的行列,会在最后\dmark{}\origpage{356}一分钟钻进来,不肯错过那终极的仪式\fnote{91}{(*)}{(6月12日)谨慎确有其必要:一个新的行列“我的学生们”已从最初称为“学生”的行列里分离出来——后者改称“学生——别名老板”。}。

\subsection{屠杀}

\notehead{87}{}

(5月12日)\fnote{92}{(**)}{本注接续前一天的“团结”一注(n\textsuperscript{o}~85)的反思。}为了多少算上同调学家的读者的教益,也尤其为了我自己,我想把这场对一个辉煌研讨班的按部就班的劫掠——它落在我的两位前上同调学家学生手里,而在其余人的善意的目光之下\fnote{93}{(***)}{反思的后续将显明:这些“其余人”中的一位,曾为这场战役有效地出手,代人行事。}——的细节巡视一遍:正是在那场研讨班里,他们比全世界早十二年、从工匠本人手里,学得了使他们成名的手艺的基础与精微之处。

我的两场口头报告,从未以任何形态交到公众手里。其一是收官报告,论未决问题与猜想——它“不幸没有成文”,鉴于其分量之轻;而屠杀版的导言作者也确实认为,连提一提所论的是\textbf{哪些}未决问题与猜想,都是多余的。他又何必费那份心呢——那不过是些问题(人人尽可随意自设!)与猜想(连证明都没有!)(注87$_{1}$)。其二是开篇报告,它把研讨班一举安放进一个更宽广的语境(拓扑的、复解析的、代数的),并巡视了Euler--Poincaré、Lefschetz、Nielsen--Wecken各型公式——其中一些,构成研讨班的主要应用之一。导言作者用“……也如同……”一笔带过这场报告的消失,那种措辞对一份\textbf{漫不经心}的心境说得很多;在当时——研讨班的作者已从流通中消失七年之际——这种心境分明是理所当然。

有整整一组报告,是我做的:关于与一个闭链相配的同调类与上同调类的形式主义(上同调情形取正则环绕概形)\fnote{94}{(****)}{细节见注n\textsuperscript{o}~82“好的参考文献”。}。它们被公平地瓜分了:上同调归Deligne,同调归Verdier——后者还是稍微越界到上同调;作为回报,则不吝以那大名鼎鼎的“权复形”向Deligne道个小小礼\dmark{}\origpage{357}\fnote{95}{(*)}{见注83“玩笑——或权复形”。}。(更不必说他还把$R\mathcal{H}om$的有限性定理与双对偶定理据为己有——一字不差地抄自研讨班;反正,狮子的份额归Deligne,那是理所当然……)导言作者认为连提一提同调的报告也无须。确实无须:前一年,他的朋友Verdier已揽下补足那份缺失的“好的参考文献”的差事(对一场研讨班、对我,均不置一词)。

有关于$R^if_*$(f非真)运算之有限性定理的口头报告,以及作为推论,关于$R\mathcal{H}om_*$与$Lf^!$运算的报告。关键定理以Hironaka式奇点消解的技术证明(因而只在拥有消解的情形有效)。我所使用的这些论证,自研讨班以降已成通例(见注85$_{2}$)。Deligne得以在另一些更便利于应用、如今在大多数应用里得到验证的假设下,证明这些有限性定理,以及双对偶定理。人们本可以指望:他会要求把这些改进纳入那场他有幸学得平展上同调、学得其后来全部著作之基础的观念与技术的研讨班。但这一情势,却被拿来当作“理由”,把研讨班的这一部分截去。至于双对偶定理,从此在Illusie的笔下(在概形框架里)成了“Deligne双对偶定理”(报告I的导言)。这也算公道:解析情形里,Verdier前一年就早已把父权判给了自己(甚至无须费心另找一份证明)。

有一场报告,发展一条“泛Künneth公式”,它曾由Illusie成文。此前没有人想到过提炼这一类陈述;其灵感在于这样的直觉:“一般地”——即在基的一般点附近——一个相对概形表现得如一个“局部平凡纤维”(在拓扑语境里)。经由一条与前面提到的证明相仿的优雅证明,Deligne得以消除我所做的奇点消解假设。判给了——报告删除,以对同一位Illusie在那场号称“在前”的研讨班SGA~4½里的报告的一条引证“取而代之”。

\dmark{}\origpage{358}有一组报告,论非交换迹的形式主义;它是作为手段发展的,用以在从未被处理过的情形里显豁Lefschetz--Verdier公式的局部项。据说,这些报告最终由Bucur成文,其手稿“在一次搬家中天赐般地丢失了”——这简直成了通俗笑剧!\fnote{96}{(*)}{想必正是这一情势,当场启发了Deligne对SGA~5的那条光辉批评:说Lefschetz--Verdier公式(别忘了——它“还停留在猜想阶段”!!!)的局部项在那里竟连计算都没有!(关于这条批评的离奇——在知情的读者看来,它与Verdier前一年那条大名鼎鼎的“权复形”难分伯仲——见“白板”一注,n\textsuperscript{o}~67。如此一来,成风的是Verdier!)}而在Illusie执笔的SGA~5导言里,这些报告居然变成了“Grothendieck的\textbf{交换}迹理论,以光辉方式推广了Stallings的理论”(而它们分明是\textbf{非}交换的!)。这个口误\fnote{97}{(**)}{即把“非交换”迹理论之父权派给我的那个口误——那顶王冠,倒也没人等我来戴。它竟然一直保留到已刊版本里,越发引人注目:Illusie是我学生中工作上最一丝不苟、直到最后细节的一位,想必。}只能出自一位秘书欠佳(或太好……)的灵感;她想必与我朋友Ionel~Bucur的搬家公司过从甚密。(“光辉地”一词是我的笔插进去的,为的是更好地还原这个同样天赐的口误所万无一失地暗示的思想。)

我无从抱怨,既然Illusie把重做的活揽在了自己身上(而且,照他的说法,还是个“更精致的”版本,因为这按层论的方式调了味——不过我记得,Illusie,在我那个年代,你做过的“精致”创新,可比这一手多……)。想必他费了不少时日——如果我没记错,当年我把这台机器调试好,花过几个星期;说不定我的手稿也在同一场天赐的搬家里丢了,天知道哪位被我的滔滔不绝淹没了的亲爱听众,至少记下过能读懂的笔记……

有一件事值得注意,我此前不曾留意:他没有把这场报告插回它原定的报告XI的位置(那想必也是它在口头研讨班里的位置),而宁肯在那个位置留下一个大窟窿,把他的报告做成一场伪托的报告,称为“局部项的计算”。这标题倒分明对应我记得自己在口头研讨班里做过的东西——奇怪。但就在他为这场报告所写导言的第1行,作者赶紧为我们\dmark{}\origpage{359}澄清误会:“这份报告,成文于1977年1月,\textbf{不对应于研讨班的任何一场口头报告}”。随后接以Lefschetz--Verdier公式(这名字分明对我说过什么;而我记得自己分明纵横发展过一套非交换迹理论,正是为了在某些情形里计算“局部项”……),继以一条Langlands公式,以及一份1967年的Artin--Verdier证明(那毕竟是口头研讨班“终曲”之后的一年;它不可能不影响到这两位作者——其中至少一位,或许两位都听过研讨班)。到这一页的末尾,我们才顺带得知——与开头的宣告相反——还有“这份报告的第二部分,性质技术得多”(这种语言我在哪里读到过……),它(请欣赏这个措辞)“\textbf{受Grothendieck为建立曲线上某些上同调对应的Lefschetz公式所用的方法的启发}”,附一处对同一场研讨班报告XII的引证,尤其对不可或缺的SGA~4½的引证;分明,为了这么点东西,没有理由把它插回大窟窿的位置——方才那份“更精致的版本”想必已尽善尽美。Illusie与Deligne还肯以“启发”之源的礼遇引我,可谓客气;而他们朋友Verdier前一年的榜样,分明早已表明:这份顾忌,完全用不着了。

我回到Illusie为顶着SGA~5之名的那一卷所写的导言。我们在那里再度得知——一如Deligne在其SGA~4½导言里已宣告的——研讨班之所以终于出版,全靠\textbf{了他的朋友}:

\begin{quote}
“我感谢 P.~Deligne说服我,在报告III的新版里,写出Lefschetz--Verdier公式的一份证明,\textbf{从而排除了这份研讨班出版的障碍之一}”。
\end{quote}

我们又一次深陷闹剧——那位顺从的Illusie,在SGA~4½的导言里照单全收!倘若研讨班十多年不曾出版,那是因为(难就难在想得到)没有一个人(在1977年救场的Deligne之前)想到过:为那条(不无道理地)称为“Lefschetz--Verdier”的公式写一份证明,也许是个好主意——而那份公式的父权,除了他那位\dmark{}\origpage{360}形影不离的朋友、我的前学生Verdier亲自,\textbf{至少从}1964年起(注87$_{2}$)就骄傲地佩戴着;也就是说,至少早在我那场研讨班结束之前两年——它当时只待好心的志愿者,便能交到人人手中!

最后,作为研讨班的又一场(最后的?)残损:Serre就“(Serre--)Swan模”所做的漂亮报告消失了——那份报告题为“Brauer理论导引”。所幸Serre,眼见事态的走向,明智地把那份报告收进了他的书《有限群的线性表示》(Hermann,1971),把它交到数学公众手里。(注87$_{3}$)

这一回,我想,这幅画我巡视完了。这幅画,关乎一场研讨班的命运——我曾把自己的最好部分倾注其中(注88)\fnote{98}{(*)}{关于“自己的最好部分”这一说法的意味,见随后的“遗骸……”“……与身体”两注,n\textsuperscript{o}~88、89。其中第一条,把SGA~5研讨班——连同与之密不可分的SGA~4——安放为我“彻底完成”的著作部分的主干。};二十年后重逢,它已面目全非,被那些曾是它独家受益人的人屠杀——或者至少,被其中三人,并获得所有其他与会者的首肯。

我并不后悔再一次费心把渐渐强加于我注意力的东西追到底。我曾在对我的一位前学生的关系的长篇回顾之末,确认这场“事物的回归”\fnote{99}{(**)}{见同名注(4月30日,n\textsuperscript{o}~73)。},当时已隐约预感到:他不是唯一一个“兴冲冲为我送葬”的人——而只是到了此刻,我才得知它的气息、它的“气味”(借用当时我的一场梦里出现的说法)——一场\textbf{暴力}的气息。这份气息,既被那种呈献高度技术性实体的(表面超然、不动声色的)话语所掩藏,又被它所泄露\fnote{100}{(***)}{尤其指伴随SGA~5(Illusie执笔)与SGA~4½(Deligne执笔)的那些导言性文字里的话语。}。这场暴力所瞄准的,经由一副任人宰割的“遗骸”,恰恰是那位曾是“老师”、“父亲”的人本身——而与此同时,“学生们”久已占据了他那令人艳羡的位置,不曾遇到任何抵抗;并且\dmark{}\origpage{361}久已在他们中间选出新的“父亲”,注定取代旧者、统治他们。

我感到这份气息,然而它于我仍是异己的、未被理解的。要“理解”它,想必这份气息得在我身上活着,或曾在我身上活过。但四年前,我头一回感到并掂量了我生活中一桩我从未想过、总以为理所当然的事:即我童年时代对父亲的认同,\textbf{不曾}带有冲突的印记——在我童年的任何时刻,\textbf{我既不怕我的父亲,也不嫉妒他},同时对他怀着毫无保留的爱。这层关系——或许是给我的一生打下最深刻印记的关系(在这场四年前的禅观之前,我甚至不曾觉察它)——在我的童年,如同对另一个既是强力的又是慈爱的自我的关系——这层关系,不曾打上分裂与冲突的印记。倘若在我那常常撕裂的一生里,对我身上所怀之力量的认识始终活着;倘若在我那绝不免于恐惧的一生里,我既不惧怕任何一个人,也不惧怕任何一件事——我把它归功于这份寒微的情势;直到我年过五十,它仍不为人知。这份情势是一件无价的特权,因为对自身之内创造之力的亲密认识——它\textbf{就是}那股力量——正使力量得按其本性自由表达:经由创造,经由一种创造的生活。

而这份特权——它使我免于冲突最深的那几道印记之一——此刻却像一副枷锁,像我生活经验里的一个“\textbf{空缺}”。一个难以填补的空缺,在那里,许多别人拥有情感、意象、联想的丰富织体,向他们提供道路(只要他们好奇肯走),既深刻理解他人、也深刻理解自己;面对那些境况——我靠反复与旁证才能勉强把握,却仍像一个异乡人——我身上求知的欲望,忍饥挨饿。

\notehead{87$_{1}$}{}

(5月31日)这场收官报告——与开篇报告一起,想必是最有趣、最实的一份——显然并非对所有人都白费,如我读到MacPherson的文章“Chern classes for singular algebraic varieties”(Annals of Math. (2) 100, 1974, 第423--432页;1973年4月收到)时所看到的。我在其中,以“Deligne--Grothendieck猜想”之名,重新找到我在这场报告里、于概形框架中引入的主要猜想之一\dmark{}\origpage{362}。MacPherson在复数域上代数簇的超越框架里重述它,Chow环换作同调群。Deligne1966年在我那场报告里学得这条猜想\fnote{101}{(*)}{(6月6日)诚然,形态略有不同;见本注今日所写的后续。\par (1985年3月)Deligne本人提供的细节,见“把一个个 i 点清”一注,n\textsuperscript{o}~164(II~1)。}——那恰是他在那场研讨班里初次露面、开始熟悉概形语言与上同调技术的一年(见“格格不入者”一注,n\textsuperscript{o}~67$'$)——几年之后,把我算进猜想的命名里,想必就不再合宜了……

(6月6日)我借此机会在此写明:我在研讨班里、于概形框架中所陈述的猜想是什么;当时我肯定还指出了它在复解析框架(乃至刚解析框架)中的明摆着变形。我把它构想为一个“Riemann--Roch”型的定理,只是取离散系数而非凝聚系数。(Zoghman~Mebkhout还对我说过:他的D-模观点应当使人可以把这两个Riemann--Roch定理都看作包含在同一条晶状Riemann--Roch定理之中;后者因此在特征零,代表我在数学里引入的两条Riemann--Roch定理——一条在1957年,另一条在1966年——的自然综合。)取定一个系数环$\Lambda$(不必交换,但为简单设为诺特的;且为平展上同调之需,其挠性与所论概形的特征互素……)。对一个概形$X$,以
$$K.(X,\Lambda)$$
记由$\Lambda$-模的平展可构造层构成的Grothendieck群。利用函子$Rf_!$,该群函子地依赖于$X$(对诺特概形$X$与分离有限型的概形态射)。对正则$X$,我假设存在一个典范的群同态,扮演凝聚RR定理中“Chern特征”的角色\dmark{}\origpage{363}:
\begin{equation}\tag{15.1}\label{eq:discrete-chern-map}
ch_X: K.(X,\Lambda)\to A(X)\otimes_{\mathbb{Z}} K.(\Lambda),
\end{equation}
其中$A(X)$是$X$的Chow环,$K.(\Lambda)$是由有限型$\Lambda$-模构成的Grothendieck群。这个同态,应当由“离散Riemann--Roch公式”对正则概形间一个\textbf{真}态射$f:X\to Y$的成立而唯一确定;该公式写作与凝聚Riemann--Roch公式同样的样子,只是把Todd“乘子”换成相对全Chern类:
\begin{equation}\tag{15.2}\label{eq:discrete-rr-formula}
ch_Y(f_!(x)) = f_*(ch_X(x)\, c(f)),
\end{equation}
其中$c(f)\in A(X)$是$f$的全Chern类。不难看出:在一个拥有Hironaka强形式奇点消解的语境里,RR公式确实唯一确定诸$ch_X$。

当然,这预设我们处在Chow环有定义的语境里。(我不知道有谁哪怕尝试过,为并非有限型于一个域上的正则概形写一套Chow环理论。)否则,也可以在凝聚语境里惯常的“Grothendieck”环$K^{\circ}(X)$所配的分次环里工作,按惯常方式过滤(见SGA~6)。还可以把$A(X)$换成偶数次的$\ell$-进上同调环,即诸$H^{2i}(X,\underline{Z}_\ell(i))$的直和。这有一个缺点:引入一个人为参数$\ell$,并给出较“纯数字”的公式;而Chow环的魅力恰在其连续的结构——过渡到上同调时被摧毁。

早在$X$是代数闭域上光滑代数曲线的情形,$ch_X$的计算就牵涉Artin--Serre--Swan型精致的局部不变量。这就是说,那条一般的猜想是一条深刻的猜想;对其的追求,系于对这些不变量的高维类似物的理解。

评注。同样以$K^.(X,\Lambda)$记由Tor维数有限的$\Lambda$-平展层可构造复形构成的“Grothendieck环”(当$\Lambda$交换时,该环作用于$K.(X,\Lambda)$……),人们同样应有一个同态
$$ch_X: K^.(X,\Lambda)\to A(X)\otimes_{\mathbb{Z}} K^.(\Lambda),$$
它同样(mutatis mutandis)引出同一条Riemann--Roch(RR)公式。

\dmark{}现在设$\mathrm{cons}(X)$为$X$上可构造整值函数的环。人们以或多或少许套套逻辑的方式定义典范同态
$$K.(X,\Lambda)\to \mathrm{Cons}(X)\otimes_{\mathbb{Z}} K.(\Lambda),\qquad K^.(X,\Lambda)\to\mathrm{Cons}(X)\otimes_{\mathbb{Z}} K^.(\Lambda).$$
倘若现在我们只限于\textbf{特征零}的概形,那么(利用具真支撑的Euler--Poincaré特征)可以看出:群$\mathrm{Cons}(X)$对诺特概形的有限型态射是共变函子(此外,作为环函子,它是反变的——这与特征无关),而前面的套套逻辑态射是函子的。(这对应于一个“众所周知”的事实——我想,在SGA~5口头研讨班里并不曾证明——即\textbf{特征零}时,对代数概形$X$上一个局部常值的$\Lambda$-模层$\mathcal{F}$,其经由$f_!:K.(X,\Lambda)\to K.(e,\Lambda)\simeq K.(\Lambda)$的像等于$d\chi(X)$,其中$d$是$\mathcal{F}$的秩,$e=\mathrm{Spec}(k)$,$k$为假定代数闭的基域……)这立刻提示:Chern同态(1)与(1$'$)应当可以从套套逻辑同态(2)与(2$'$)出发、复合一个“普适”Chern同态(不依赖任何系数环$\Lambda$)
$$ch_X: \mathrm{Cons}(X)\to A(X)$$
而导出;于是,RR公式的两个“带$\Lambda$系数”的版本,便在形式上显现为包含于一条在可构造函数层面上的RR公式之中——后者总以同样的形态书写。

当我们对固定基域(重新允许任意特征)上的概形工作,或更一般地对固定的\textbf{正则}基底概形$S$(比如$S=\mathrm{Spec}(\mathbb{Z})$)上的概形工作时,与惯常写法(自1957年以来熟悉的凝聚框架)最贴合的RR公式形态,经由引入乘积
$$ch_X(x)\, c(X/S) = c_{X/S}(x)$$
而得到(其中$x$在$K.(X,\Lambda)$或$K^.(X,\Lambda)$中皆可);人们可以称之为$x$相对于基$S$的\dmark{}\origpage{365}\textbf{Chern类}。当$x$是$K^.(X,\Lambda)$的单位元、即取值$\Lambda$的常值层的类时,我们得到$X$相对于$S$的相对全Chern类在典范同态$A(X)\to A(X)\otimes K^.(\Lambda)$下的像。这么规定之后,RR公式等价于:对$S$上(有限型的)变动的正则概形$X$、固定$S$,这些相对Chern类的作法
$$c_{X/S}: K.(X,\Lambda)\to A(X)\otimes K.(\Lambda)$$
对真态射是函子的,(5$'$)的变形亦然。特征零时,这化归为相应映射
$$c_{X/S}: \mathrm{Cons}(X)\to A(X)$$
对真态射的函子性。正是在“绝对Chern类”映射(6)的存在与唯一这一形态下——在$S=\mathrm{Spec}(\mathbb{C})$的情形——该猜想呈现于MacPherson的工作之中;相关条件(这里与特征零的一般情形一样)是:a)(6)对真态射的函子性;b)$c_{X/S}(1)=c(X/S)$(在这里,即“绝对”全Chern类)。相对于我最初的猜想,MacPherson所呈献并证明的形态,有两点不同。其一是“减”:他不置于Chow环、而置于整系数上同调环——更确切地说,经由超越途径定义的整系数同调群。其二是“加”——也许正是在这里,Deligne对我最初的猜想做出了一份贡献(除非这份贡献出于MacPherson本人\fnote{102}{(*)}{(1985年3月)确实如此;见前一条脚注所引的注n\textsuperscript{o}~164。}):即对映射(6)的存在与唯一,无须把自己限于正则概形$X$,只要把$A(X)$换成整系数同调群。由此很可能:在一般情形亦然,只要以诺特概形$X$的\textbf{Chow群}(一般不再是环)充当$A(X)$(或更好:$A.(X)$)。或者换个说法:虽然$ch_X(x)$的启发式定义($x$在$K.(X,\Lambda)$或$K^.(X,\Lambda)$中)本质上地利用了环绕概形正则的假设,可一旦把它乘上“乘子”$c(X/S)$(当概形$X$有限型于固定正则概形$S$之上),所得乘积(4)似乎无须对$X$的任何正则性假设仍保有意义\dmark{}\origpage{366}——作为张量积
$$A.(X)\otimes K.(\Lambda) \quad\text{或}\quad A.(X)\otimes K^.(\Lambda)$$
中的元素,其中$A.(X)$记$X$的Chow群。MacPherson的证明精神(它不使用奇点消解)提示:(5$'$)的同态有可能有一个显式的“演算式”构造:“将就着”$X$的奇点本身,连同系数层$\mathcal{F}$(其类为$x$)的奇点,在$X$上“收取”一个取值于$K^.(\Lambda)$的闭链。这也合乎我1957年随凝聚Riemann--Roch定理引入的观念的精神:我在其中尤其做自交计算,小心不让所论的闭链“动”。一个显明的初步归约(经由把$X$嵌入一个$S$-概形)是$X$为正则概形$S$的闭子概形的情形……

此外,应当可能发展一条\textbf{奇异的}(凝聚)Riemann--Roch定理——这个想法我早已熟悉,说不出从何时起;只是我从不曾认真检验过它。某种程度上,正是这个想法(除去与“上同调、同调、卡积”形式主义的类比)曾在SGA~6里(1966/67)使我系统地引入$K.(X)$、$K^.(X)$与$A.(X)$、$A^.(X)$,而不满足于只用$K^.(X)$工作。我不记得1966年在SGA~5研讨班里是否也想过这类东西、是否在口头报告里暗示过。由于我的手写笔记已散失(也许在一次搬家里?),我恐怕永远无从知道了……

(6月7日)翻阅MacPherson的文章,有一件事打动了我:“Riemann--Roch”一词在其中未曾出现——这也正是我没有立刻认出我1966年在SGA~5研讨班里所做那条猜想的缘由;那条猜想,对我来说(现在仍是),是一个“Riemann--Roch”型的定理。写文章之际,MacPherson似乎甚至不曾觉察这层显然的亲缘。我料想缘由在于:Deligne——我离去之后,他以自己乐意的形态使这条猜想流通——曾尽力“抹去”它与Riemann--Roch--Grothendieck定理的显然亲缘。我自以为感到他如此行事的动机。一边,这削弱了这条猜想与我个人的联系,并使它目前流通的那个名称“Deligne--Grothendieck猜想”更可信。(附记:我不知道它在概形情形是否也在流通;倘若在,我很想知道顶着什么名称。)但我看更深的缘由,在于他那个执拗的念头:尽一切可能否认并摧毁我的著作、我的数学视野的根基处的统一性\fnote{103}{(*)}{可与“遗骸”一注(n\textsuperscript{o}~88)里关于“SGA~4½之役”之深层意味的评语对读:那场战役同样意在,经由外来文本SGA~4½对SGA~4与SGA~5这两个不可分割的部分——它们正发展着我的著作——的“强行塞入”,把我的著作围绕平展上同调的深层统一,炸成一堆无定形的“技术枝蔓”。}。这是一个触目的例子,说明在一个手段出众的数学家身上,一个与一切数学动机全然无关的执念,竟能使那种我称之为数学的“健全本能”的东西昏暗(乃至全然堵塞)。这份本能不可能不察觉一个“同一条”Riemann--Roch定理的“连续”与“离散”两个陈述之间的类比——我当年在口头报告里当然也强调过。如我昨天所写,这份亲缘想必不久将由一个成形的陈述(Zoghman~Mebkhout所猜)确认——至少在复解析情形——使人得以从一条共同的陈述导出两者。显然,在Deligne对着Riemann--Roch定理的那份“掘墓”心境里\fnote{104}{(**)}{这些心境——恰对着Riemann--Roch--Grothendieck定理——在“葬礼悼词”里表现得格外清楚;见“葬礼悼词(1)——或恭维”一注,n\textsuperscript{o}~104。},他根本不可能发现那个在解析框架里联结两者的唯一陈述,更无从在一般概形框架里提出类似陈述的问题。一如他在这样的心境里,无从提炼D-模这个丰产的观点——它以过于自然的方式,出自他正要埋葬的观念——甚至在多年之间,无从承认Mebkhout那项在他失败之处成功的丰产工作。

\notehead{87$_{2}$}{}

(5月31日)那是1957年——我的Bourbaki报告论L函数有理性之年;我在其中启发式地使用Verdier的结果(???)(尤其该个案里局部项的预期形态),而不等Illusie十三年后——应Deligne之邀——才肯去证明它。而且我觉得,当年Verdier把他那条超\dmark{}\origpage{368}一般的公式拿给我看——它来得如一场意外——时,他是用“六种运算”形式主义在几行里证明它的:这类公式,(几乎)写出来就是证明了它!倘若有什么“困难”,至多在核验一两处相容性的层面上\fnote{105}{(*)}{(6月6日)此外似乎:经由双对偶定理(其间已晋升为“Deligne定理”),Lefschetz--Verdier公式的初始证明曾依赖一条奇点消解假设;Deligne得以在一域上有限型概形的情形摆脱它。这正是浑水摸鱼的大好机会,给人这样的印象:SGA~5竟从属于那个“在先(原文如此)”的“研讨班SGA~4½”(而且它确实先出版了!)。}。况且,无论Illusie还是Deligne,都一清二楚:我在研讨班里为各种显式迹公式所给的证明\textbf{是完整的};它们绝不依赖Verdier那条一般公式——后者只扮演了“扳机”的角色,怂恿人们在尽可能一般的情形里显豁并证明迹公式。两人在这里的不诚,昭然若揭。对Deligne,我写“白板”一注(n\textsuperscript{o}~67)时就已清楚——但对不知情的读者,当然不清楚;同样不清楚的,还有那种放弃使用自己健全官能的知情读者。

(6月6日)至于Illusie,他完全入了其朋友的戏,试图搅浑水,摆出一副超技术的口头研讨班的样子——仿佛它连全部结果的完整证明都给不出,尤其迹公式。然而那些公式在65/66年间\textbf{确实}得到证明(而且是头一回),恰恰是在那里,他与Deligne有幸学得它们,以及与之相随的一整套精微技术\fnote{106}{(**)}{在顶着SGA~5之名出版的那一卷的导言第二段里,Illusie把围绕平展上同调Lefschetz公式的三场报告III、III~B、XII当作“研讨班的心脏”;而人们看到,在报告III~B的导言里,他小心注明(与事实相反)“这份报告不对应于研讨班的任何一场口头报告”;在报告III与III~B的导言里,他尽力给人这样的印象:它们从属于SGA~4½,而报告III被呈献为“猜想性的”!!实际上,SGA~5研讨班的全体,在技术上独立于报告III(Lefschetz--Verdier公式)——后者扮演启发性动机的角色;而报告III~B不是别的,正是Bucur那次搬家所造成的“窟窿”(报告XI);那窟窿,成了这场进一步肢解的可喜口实。为使人相信那场“技术枝蔓”研讨班的版本(其朋友Deligne的授意),Illusie小心砍掉了开篇报告——我在其中为将在这场研讨班里发展的各大主题画了一幅初步的全景;迹公式在其中只占一小角(因其朝Weil猜想方向的算术含义而取得特殊重要性)。关于这些“大主题”的概观,见后面的子注n\textsuperscript{o}~87$_{4}$(原书误作87$_{5}$)。}。

\dmark{}这使我想起:当然,我曾费心在研讨班里证明Lefschetz--Verdier公式\origpage{369}——那是最起码的,也是我打算发展的局部与整体对偶形式主义的一次格外触目的应用。这几天问题来到我心头:既然分明有十来份报告的成文工作,在我亲爱的学生们手里陷于困境,Deligne与Illusie要为其SGA~5出版的“障碍(原文如此)”命名,本有挑选的余地——为什么偏偏选中他们的好哥们Verdier的定理——他当时正把父权当作应得之物佩戴着,一如他佩戴导出范畴与三角范畴的父权,而那些他也从不曾费心成文(或至少交到公众手里)。这里有一种\textbf{荒谬之中的挑衅}(或我的前上同调学家学生群体里的某种集体犬儒——我认为他们在那场屠杀战役里全体同谋),它使我想起Verdier前一年光辉地发明的“权复形”(见同名注,n\textsuperscript{o}~83),或(在不义的音区里)Deligne给本该叫“Mebkhout的层”的那些层所取的名字“反常”(见“反常性”一注,n\textsuperscript{o}~76)。我在这样的发明里感到,是对整个数学共同体的支配与轻蔑之举——同时是一份\textbf{赌注};它分明一直赢到亡者不期现身之时——他俨然像一群沉睡者中间唯一醒着的人……

\notehead{87$_{3}$}{}

(6月5日)一场屠杀盘点之后,人们自会掂量Illusie在其顶着SGA~5之名的卷子的导言第2行里的这份自白:

\begin{quote}
“相对于最初的版本,唯一重要的变更涉及报告II[泛Künneth公式],未予复制;以及报告III[Lefschetz--Verdier公式],已全部重写,并增补一个编号III~B的附录\fnote{107}{(*)}{它被呈献为“研讨班的心脏”的一部分!(见前一条脚注。)}。除若干\dmark{}\origpage{370}细节修改与脚注增补之外,其余报告一仍\textbf{其旧}”(着重号是我加的)。
\end{quote}

这里,Illusie又一次甘当其那位罄竹难书的朋友另一句俏皮话的应声虫,即:SGA~4½的存在“不久将使人得以\textbf{照原样}出版SGA~5”(见“白板”一注,n\textsuperscript{o}~67)——而Illusie在其报告与导言里,处处尽力为这一冒名张目——即:SGA~5(他与朋友学艺之地)竟依赖于那部海盗文集SGA~4½,由其后十二年间拾得或劫掠的破烂与断枝拼成——每页转弯处都不吝对SGA~4½的大量引证……

最后的字,照例归Deligne:一个月前(5月3日)他给我写信,回复一份简短的信息之请(见“葬礼”一注,n\textsuperscript{o}~70的开头):

\begin{quote}
“总之,倘若SGA~4½这份文本问世时,你已七年不做数学[?!],那只是相应于SGA~5出版[?]的漫长迟延,\textbf{它太不完整,无法照原样有益地出版}。

希望这些解释令你满意。”
\end{quote}

它们即便不曾“令我满意”,至少令我大开眼界……

\notehead{87$_{4}$}{}

(6月6日)或许是时候指出:口头研讨班里发展过的主要主题有哪些;已刊文本,只有靠旁证才能对其形成概念。

I)对偶理论的局部侧面;其本质技术配料是(一如凝聚情形)双对偶定理(补以一条“上同调纯粹性”定理)。我的印象是:这后一条定理的几何意味——作为一条局部Poincaré对偶定理——我当年在口头研讨班里分明解释过,却被我的学生们从此全然忘却\fnote{108}{(*)}{核实之后:这一几何诠释,至少在Illusie的成文里得到了保留。}。

II)迹公式,包括比惯常迹公式(其两边为整数,或一般为某系数环——如$\underline{Z}/n\underline{Z}$或$\ell$-进环$\mathbb{Z}_\ell$、乃至$\mathbb{Q}_\ell$——的元素)更精微的“非交换”迹公式;它们置于一作用在所论概形上的\textbf{有限}群的代数之中,系数取于一个合宜的环(如前一括号里所设想者)。这一推广来得十分自然:因为即便在惯常型的Lefschetz公式里,对“扭转的”系数层,人们也被引至以一个Galois覆盖(一般分歧)替换原概形,以“解扭”系数,而Galois群作用于其上。如此,“Nielsen--Wecken”型的公式,便自然地进入概形语境。

III)Euler--Poincaré公式。一边,借助Serre--Swan模,对代数曲线的一条“绝对”公式作了周详研究(推广了驯分歧系数的情形;该情形给出较朴素的Ogg--Chafarévitch--Grothendieck公式)。另一边,有未发表的、深刻的“离散”Riemann--Roch型猜想;其中一条,七年后以一种杂糅的版本重现,顶着“Deligne--Grothendieck猜想”之名,由MacPherson经超越途径证明(见注n\textsuperscript{o}~87$_{1}$)。

我对这两大主题(Lefschetz公式、Euler--Poincaré公式)之间深刻关系所能做的评语,也一并散失,未留痕迹。(一如我的习惯,我把全部手写笔记留给了“志愿者们-执笔者(原文如此)”;口头研讨班的手写笔记,我本来有一整套——尽管有些简略——如今片纸无存。)

IV)与一个闭链相配的同调类与上同调类形式主义的详论,自然出自一般对偶形式主义与那条关键想法:在所论闭链上做“具支撑”的上同调,利用上同调纯粹性定理。

V)有限性定理(包括泛型有限性定理)与任意支撑上同调的泛Künneth定理。

研讨班还发展了一套从挠系数过渡到$\ell$-进系数的技术(报告V与VI)。那是研讨班最技术的部分;研讨班一般用挠系数工作,事后再“取极限”以导出相应的$\ell$-进结果。这一观点是一个临时的\dmark{}\origpage{372}权宜之计,且待Jouanolou的论文(至今未出版!)直接在$\ell$-进框架里给出所需的形式主义。

我不把对几个经典概形的一些计算、以及Chern类的上同调理论,算进“主要主题”;Illusie在其导言里把它们捧为研讨班“最有趣的”之一。由于日程已满,我在口头研讨班里不曾认为有必要在这些计算与这一构造上流连:只需把十年前我为Riemann--Roch定理、在Chow环的语境里给出的论证,几乎一字不差地搬来即可。另一方面,明摆着应当把它们纳入成文的研讨班,以便给平展上同调的使用者一份便利的参考。Jouanolou揽下了这项工作(报告VIII);他本该把它看作一项美差——既服务于数学共同体,又为自己学到一些日后必用的基本技术——却把它当作苦役:其成文一拖数年\fnote{109}{(*)}{(6月12日)翻阅那份报告,我进而确信:Jouanolou与我的其他上同调学家学生之间,有完美的默契。}。想必他的论文也如出一辙:至今仍是一份幽灵参考,一如Verdier的那份……“极限过渡”那一部分也不该算作研讨班的“主要主题”之一,因为它不与某个特定的几何想法相连。毋宁说,它反映平展上同调语境(区别于超越语境)特有的一项技术复杂:平展上同调的主要定理首先关涉挠系数(与剩余特征互素);要有对应于特征零系数环的理论——Weil猜想所必需——必须对$\mathbb{Z}/\ell^n\mathbb{Z}$型系数环取极限,以获得“$\ell$-进”结果。

把这一切说清之后,口头研讨班五个主要主题里,唯一一个似乎以完整形态出现在已刊文本里的,是主题I。主题IV与V干脆消失,被SGA~4½吸收,附带的收益是可以大量引证它、给人这样的印象:SGA~5依赖一份自居在先的Deligne文本。主题II与III以残缺形态出现在已刊卷里,并始终维持对SGA~4½文本的同一冒名(那份文本,实际上整个出自母研讨班SGA~4、SGA~5)。

\subsection{遗骸……}

\notehead{88}{}

\dmark{}(5月16日)SGA~4与SGA~5这两场前后相续的研讨班——对我而言,它们如\textbf{同一场}“研讨班”——从虚无中同时发展出:topos的\textbf{语言}——强大的综合与发现的工具;以及平展上同调——一件打磨停当、效率完美的\textbf{工具};其本质形式性质,彼时甚至比寻常空间的上同调理论理解得还透\fnote{110}{(*)}{即便把自己限于最贴近“簇”的空间,如可三角剖分的空间。}\origpage{373}。这一整体,代表我在数学里、就彻底完成的工作而言所能做出的最深刻、最革新的贡献。同时——它并无意于此,虽然一切始终以事情之明摆着的自然进行——这项工作代表我作为数学家在其著作中所完成的规模最大的技术壮举\fnote{111}{(**)}{若干困难或意外的结果,是由别人(Artin、Verdier、Giraud、Deligne)得到的;工作的某些部分是与别人合作完成的。这丝毫不减损(至少在我心里)我对这项工作在我全部著作中之地位的评判的分量。况且,我想在《主题概要》的一篇附录里,更周详地回到这一点,在分明已属必要之处把点补在i上。}。这两场研讨班,对我是密不可分的。它们在其统一之中,同时代表\textbf{视野}与\textbf{工具}——topos,与一套完整的平展上同调形式主义。

视野至今仍遭拒斥,而工具近二十年来深刻革新了代数几何其最令我心醉的一面——“算术”的一面:经由一种直觉、一套“几何”性质的概念与技术行装所把握的算术。

肯定不只是为了暗示其上同调“摘要”先于SGA~5部分,Deligne才给它戴上SGA~4½这个欺眼之名——反正要戴,大可以叫SGA~3½!在“SGA~4½之役”里,我感到一种意图:把那部他的全部著作皆出于斯的著作(那部他始终无法摆脱的著作!)——其分明而深层的统一,在SGA~4与(真正的)SGA~5两场研讨班的整体里昭然可见——呈献为一副\textbf{分裂}的面目(一如他自己也是分裂的……),被一份外来的、轻蔑的文本的强行塞入\textbf{一分为二};那份文本,妄想显出是活的心脏、思想与\dmark{}\origpage{374}视野的精髓——而它在那视野里不曾有过任何份额\fnote{112}{(*)}{这份思想——就其主导观念与本质结果而言——早在年轻的Deligne登台、来与我接触学代数几何与上同调技术(1965至1969年)之前,就已完全成熟。(5月30日)对此可参见“格格不入者”一注,n\textsuperscript{o}~67$'$。};而环绕它的两个“断片”,则成了近乎怪诞的附录,一堆对那部自居为中心与本质的著作——出自Deligne的笔下,而我寒微的身量则被客气地接纳(在彻底下葬之前)跻身“合作者”之列——的“枝蔓”与“技术补充”\fnote{113}{(**)}{见“绿灯”与“颠倒”两注,n\textsuperscript{o}~68、68$'$。}。

“偶然”把事情安排得妙极。这副“任人宰割的遗骸”——那场被“执笔者们”一直弃置不顾的“倒霉研讨班”,我离去时落在我的上同调学家学生们手里、听凭其处置的——那可不是老师著作中\textbf{随便}哪一部分!既不是SGA~1与SGA~2(我在其中独自发展、尚不自知那些即将成为未来主作“起飞”的两件不可或缺的技术辅助工具的东西),也不是SGA~3(我的贡献主要是为使概形的“四面八方”技术纯熟而做的无穷音阶与琶音——有时颇费力气),也不是SGA~6(系统发展我围绕Riemann--Roch定理与相交形式主义的十年旧观念),乃至SGA~7(按一场反思的内在逻辑,它从对中心工具的持有、对上同调的掌控中流出)。分明是我著作的\textbf{主干部分}——其成文工作尚未完成(而且是经他们的手……)——我把它,至少部分地,留在了我的上同调学家学生们手里。他们选择屠杀的,正是著作的这一主干;他们据为己有的碎块,忘了那使其获得意味、美与创造之德的那份统一(注90)。

同样不是偶然:他们手执七拼八凑的工具、否认使其从虚无中诞生的精神与视野,竟无一人能在革新之著作重生之处——逆着他们的冷漠与轻蔑——辨认出它。也无从解释:六年之后,当那件新工具终于被Deligne把握,他们竟全体一致地埋葬了那位在孤独中创造它的人——Zoghman~Mebkhout\dmark{}\origpage{375},被否认的老师之遗孤学生!同样不是偶然:在Deligne最初的冲劲回落之后(那冲劲,在几年间把他带向一套新Hodge理论的强力启动,与Weil猜想的证明),尽管他手段惊人、我的上同调学家学生们手段出众,我今天仍确认这场“阴郁停滞”——在一个无比丰饶、一切似乎尚待去做的领域里。这无须惊讶:近十五年来,灵感的主要来源与某些“大问题”\fnote{114}{(*)}{这份“灵感的主要来源”,当然是“Motif瑜伽”。它只在Deligne一人身上起过作用——他把它留在怀里独享“收益”,且以一种被剥除其大半力量的狭隘形态,拒斥这部瑜伽的某些本质方面。在由它激发、却遭忽视或被悄悄抹黑的“大问题”之中,我如今(虽不折不扣是局外人)看到:标准猜想;以及面向一切惯常类型的系数——或多或少贴近“Motif”本身(它们对后者扮演“普适”系数的角色:生出所有其他系数的那种)——的“六种运算”形式主义的发展。可与“我的孤儿们”一注(n\textsuperscript{o}~46)里的评语对读。},即便它们在场、人们每走一步都与它们遭逢,却被小心绕开、藏掖起来,如那位十五年来人们不断要埋葬的人的信使一般。

\subsection{……与身体}

\notehead{89}{}

(5月17日)那份曾在我身上活着的念头、那种事物的视野——我曾以为已把它传递出去——我看作一个活的身体:健全而和谐,带着活物更新的力量、构想与生育的力量。瞧,这个活的身体变成了一副\textbf{遗骸},在人们中间被瓜分——这条腿或这半边身子,被妥善制成标本,在某一家充当战利品;另一条,被剁碎,在另一家充当击头的棍棒或回旋镖;还有一条,天知道,原样下到自家的汤锅里(我们已经不在乎这个了!)——其余的,统统可以烂在垃圾堆里……

这就是——形象化的措辞诚然,但在我看来,恰当地表达事情的某种实在——终于向我显现的图景。那根棍棒,至多会在这里或那里敲裂一两颗头颅\fnote{115}{(**)}{(5月31日)而且它还满可以用来证明某条“困难尽人皆知的”定理!};但这些四散的碎块——战利品也罢、棍棒也罢、家常汤也罢——将永远不具那个活的身体那份如此简单、如此分明的力量:那创造一个新生命的、爱的拥抱的力量……

(5月18日)这个活的身体、这副碎块随风四散的“遗骸”的意象,想必是在刚过去的一星期里在我身上渐渐成形的。\dmark{}它在我的打字机笔下呈现的那副滑稽形态,丝毫不意味着这意象是半点\textbf{虚构}、一场顺着话头即兴而来的略带阴森的滑稽表演。这意象表达一份\textbf{实在};当它经书面表述取得物质形态的那一刻,我被深深地触动。\origpage{376}这份实在,想必早在我“离去”以来的十四年里、甚至更早,已在此处彼处以断片的方式被我获知。一些信息的断片,起初由一份心不在焉、倾注别处的注意力,登记在一个表浅的层面——但它们全都指向同一方向,想必已在更深的层面,聚成某种意象——一幅未曾表述的意象,我不曾费心去认识它,那时我有更要紧的事。这份意象,在3月底以来持续进行的反思中、即这六七个星期里,大大丰富并明确了。更确切地说:零散的信息要素,终于经一份全然在场的、有意识的留心的检视,渐渐聚成\textbf{另一幅}意象——在检视与探查的思维这个较表浅的层面上;经由一项工作,它看似独立于第一幅意象在更深层的在场。这项有意识的工作,六天前在“屠杀”之发生的骤然视野中达到顶点——那时我感到了一份\textbf{暴力}的“气息”、“气味”;我想,这是整场反思里的头一回\fnote{116}{(*)}{(6月12日)近些年里,我曾有过这样的时刻:感到我的某些前学生对另一些与我“同葬者”怀有暴力的意图;但从不曾感到一种出自集体意志(在这里结成五人)、指向我这个人、经由我的著作而行的暴力。}。也是在这样的时刻,一个健全和谐之躯分明遭“屠杀”的感觉,想必已在临近表面的层次出现;而那幅更深的弥散意象,想必也已开始浮出水面,为那幅成形中的意象带来一层血肉——一种单凭思维无力给出的“气味”。

这份“血肉”的侧面,昨夜又在一场梦里显现——正是受这场梦的推动,我此刻回到昨天写下的文字。梦里,我身体的几处地方被割了相当深的口子。起初是嘴唇与口腔里的割口,大量出血;我对着镜子用大量清水漱口,漱口水被血染得通红。然后是腹部的伤口,也大量出血,尤其其中一处,血一鼓一鼓地涌出,仿佛是动脉(做梦者无心顾及解剖学的写实)。我甚至想到:倘若照这样流下去,我兴许会躺倒在地的;我用手按住伤口,蜷起身子止血——血确实止住了奔流,渐渐结成一个血块与一大片痂。后来,我小心地揭开那片痂:细微的愈合已经开始。我的一个手指也被割了,上面缠着厚厚一团绷带……

我无意对这个梦做更精微、更周详的描述,也不打算在这里(或别处)深探它。这个梦“照原样”已经以惊人的力量向我显示:昨天我谈的那个“身体”——我在写它时,把它看作与我分离的东西,像一个我或许孕育、生出、如今已走向世界走自己路的——这个身体,至今仍是我这个人亲切的一部分:那是\textbf{我的}身体,由血肉与一种生命的力量构成,这力量使它经得起深的创伤、并自我再生。而我的身体,想必也是世上与我系得最深、最不可解的东西……

做梦者没有让我停留于“屠杀”与瓜分遗骸的意象。那份意象要还原的,是我曾强烈知觉到的、\textbf{他人}身上的一份意图与心境的实在,而不是我自己经历这场侵凌、这场残损——经由一件我仍与之紧紧相连的东西——的方式。我与之相系到何等地步,做梦者刚刚让我隐约瞥见。这与我在“事物的回归(或一脚踩进菜盘)”一注(n\textsuperscript{o}~73)的反思里(以较弱的力度)所知觉的东西相合:在那里,我试着稍稍勾勒“构想了一件事的人与这件事之间的深深纽带”的那份感觉——它出现在那天的反思里。在那场4月30日的反思(仅仅三周之前)之前,在我的一生里,我曾装作无视这条纽带,或至少淡化它,顺着时兴老生常谈那条划好的坡道。操心一件已脱手的作品的命运——尤其操心我们的名字是否还多少附着于它——被感受为一种小家子气、一种悭吝;然而,人人却觉得自然:一个我们养大的(并且自以为爱过的)亲生孩子,倘若选择否认他出生时所得的名字,我们会深受触动。

\subsection{继承人}

\notehead{90}{}

\dmark{}(5月18日)我不知道六十年代里,可曾有哪个学生(Deligne除外)感到了这份本质的统一——超越他与我一起进行的有限工作。也许某些人模糊地感到过;而这份知觉,在我离去之后的年月一过,便永无挽回地失落了。可以肯定的却是:从1965年我们初次接触起,Deligne就预感到了这份活的统一。\origpage{378}正是这份对一个宏阔蓝图之统一旨趣的精微知觉,无疑是他身上那份强烈兴趣的主要兴奋剂——对着我有可传授、可传递的一切。这份兴趣,在1965到1969年这四年持续的数学接触里,从不曾衰减\fnote{117}{(*)}{这一时期包括五年,其中一年(1966),我的朋友在比利时服兵役。}。它赋予我们之间的数学交流以我谈过的那份卓越的品质;与别的数学家朋友,我只是在寥寥的瞬间知晓过它。正是这份对本质之物的知觉,以及它在他身上激起的那份炽热的兴趣,使他得以仿佛游戏般学得我能教他的一切:既有技术手段(锌丝概形的技术,Riemann--Roch与相交的瑜伽,上同调形式主义,平展上同调,topos的语言),也有使其统一的总体\textbf{视野},最后是\textbf{Motif瑜伽}——它当时是这份视野的主要果实,是我到那时为止所发现的最强的灵感之源。

清楚的是:Deligne是我的学生中直到今天仍唯一的一个,曾在某个时刻(我想自1968年起)把我可传递的一切,就其本质的统一与其手段的多样性,完全消化并据为己有\fnote{118}{(**)}{当我谈“全部”时,应理解为:就视野与手段中的本质事物而言。这当然不意味着,不曾有一些未发表的想法与结果,我从未想过向他谈起。相反,我不认为1965--69年间的任何一场数学反思,是我没有“趁热”向我的朋友谈过的——始终怀着乐趣,也始终有所获益。}。正是这份情势——想必人人都感到了——使他俨然我著作天造地设的“合法继承人”。显然,这份遗产既不使他累赘,也不限制他——它不是一副重担,而给他添翼;我是说:它以其活力滋养着他生而有之的那些“翅膀”\dmark{}\origpage{379}——一如别的视野、别的遗产(诚然没那么个人……)也将滋养它们……

这份他曾在成长与起飞的关键年月里赖以为养的遗产,以及使其具有美与创造之德的那份统一——他曾那么懂得去感觉、它已成为他自身一部分的东西——我的朋友此后\fnote{119}{(*)}{奇怪的是,这份分裂,想必在我们相识的头一年就已在场(它已经表现为对SGA~5研讨班的一份暧昧态度——那是他与概形、与Grothendieck风格的上同调技术、与平展上同调的初次接触),最迟在1968年、以毫不含糊的形态(见“驱逐”一注,n\textsuperscript{o}~63)——也就是数学交流完美无缺、他的数学思想的起飞在我看来尚未打上冲突印记的时刻。他当时(“顺带地”)做出过许多有趣的贡献(我在SGA~4的导言里极乐于把它们捧出来),题材却是他在我离去之后竭力要埋葬的。}予以否认,不懈地藏掖遗产,否认并摧毁那份作为其灵魂的创造性的统一。在我的学生中间,他是第一个做出榜样的人:一面据有工具、据有“碎块”,一面拼命肢解它们所自出的统一、那个活的身体。他自己的创造冲劲,由此被他身上这份深层的分裂所阻滞、吸收、最终肢解;这分裂驱使他否认并摧毁恰恰构成其力量、滋养其冲劲的东西。

我看这份分裂经由三个互相支援、密不可分的效果表达。其一是能量涣散的效果:在否认、肢解、取代、藏掖的努力中四散。其二在于对某些观念与手段的拒斥——它们对本是他选作中心主题的那个主题的“自然”发展本质不可少\fnote{120}{(**)}{这份拒斥,尤其表现为:导出范畴与三角范畴被埋葬(直到1981年),六种变换的形式主义被埋葬(直到今天),topos的语言被埋葬(亦然);以及对同调代数与同伦代数那套宏阔奠基纲领的一种“以轻蔑来封锁”——那纲领,我现在(二十年后)试图以《追寻栈》给出一个草图,而他当然同样不曾不感到其需要。最后,即便他从Motif瑜伽(被埋葬直到1982年)汲取灵感,这部瑜伽仍被削去其一部分力量:它与六种变换的形式主义——其一个本质的形式侧面——脱了钩。依我看,这一侧面,也被严格地逐出Hodge--Deligne理论之外。}。其三,是执著于那个首当其冲的主题:在其中,要取代、排挤一位每走一步都在场的老师、并无休止地抹去他——恰恰是这个主题,被注入最强烈的那份矛盾,它主宰其数学家的一生\dmark{}\origpage{380}。

我第一手所知的东西,加上一份从不曾骗过我的素朴本能或嗅觉,使我清楚:倘若Deligne不曾被这份其工作内部的深层矛盾撕裂,今天的数学便不会是如今这副样子\fnote{121}{(*)}{写“今天的数学”这些文字时,我想的不只是我们今天对数学事物或多或少的认识。背景里还有对一个特定\textbf{精神}的念头——在数学家的世界里,尤其在被人们(不带讥讽或嘲弄的口吻)称为“上流”数学界的东西里:那个“定调子”、裁定何为“重要”乃至“正当”、何为不然的精神;那个控制信息手段、并在很大程度上控制前程的精神。也许我夸大了身处旗手之位的一个人对给定环境在给定时代的“时代精神”所能有的影响。Deligne的影响,依我看,堪比(无论好坏)二十年前Weil在接待我的那个环境里所显得具有的影响——那个环境,我曾二十年认同于它。(5月31日)可与“掘墓人——或全体教众”一注(n\textsuperscript{o}~97)的(互补)反思对读。}:它本会在其若干本质的部分,经历宏阔的革新——一如我曾亲身充当其主要工具的那一场革新——恰恰是同一位Deligne拼命对抗、使之改道的那一场!\fnote{122}{(**)}{(6月16日)我深信:单凭我引入数学的主导观念正常发展这一点——沿着六十年代已获得的冲劲(那冲劲,被“电锯效应”猛然切断,下两条注将谈到它……)——今天的数学,在我离去十五年后,本会在其若干本质的部分,不同于如今的样子……}

同样毫无疑问:他天造地设,本该成为一所强力几何学派的灵魂——延续那所曾围绕我形成的学派:一所汲取其母体活力、以及接我班者之创造力的学派。但这所曾围绕我形成的学派、这座曾环绕那些紧张的受业年月的滋养母体,在我离去后的翌日便已解体。所以如此,恰恰是因为在分明接我班的那个人身上\fnote{123}{(***)}{这份\textbf{既成事实的接班},以毫不含糊的具体征兆表达:他在IHES接了我的班(我是在他入所的次年离开的——见“驱逐”一注,n\textsuperscript{o}~63);并且他接手了代数簇上同调这一中心主题,用的正是我为这一目的、在十五年间(1955至1970年)发展的手段。},找不到那个也将是一场共同冒险所结合的群体的灵魂的人——为一项其规模超出每个人的手段的任务。

我的印象是:我离去之后,我的每一位学生都退回自己的角落——工作自然有的是,数学里从不缺工作——但这个“角落”不再嵌入一个整体,这份“工作”不再被一股潮流、一个更宏阔的旨趣所承载。肯定,自我离去起——甚至更早——我的大部分学生或前学生的目光,都转向那位天造地设的“继承者”:他们中间最光彩、也离我最近的一位。在这个敏感的时刻,我的朋友想必平生头一回感到那种对他人的权力突然握在自己手里:那种对一个学派的命运生杀予夺的权力——那学派,他出身于斯;与他同学四年的朋友们,想必期待他保证其延续。局面整个在他手里,是他来定调子……他也确实定了——经由摧毁遗产;而首先摧毁的,是那份信任与那份期待\fnote{124}{(*)}{(5月26日)在反思的进程中,我还察觉对我的默示继承人的另一份全然不同的“期待”——这一回不出自我的学生们,而出自“全体教众”——对此可参见“掘墓人——或全体教众”一注(n\textsuperscript{o}~97)的结尾。我几乎不怀疑:这两份方向相反的期待——一份系于一个非常特定的时刻,另一份贯穿葬礼的十四年——都是实在的。更有甚者,我倾向于认为:在我的不止一位昔日学生身上,两份期待想必同时在场:一面是期待在他们的最光彩者身上,也找到那保证一所学派、一部著作——他们在其中有位置、有份额——之延续的人;一面是期待看到(倘若可能)那个人的任何痕迹被抹去——他的离去,曾在划好的坡道的安稳之中,以如此的力量,突然向他们发问……}——那些曾与他同学于同一位老师的人,本不会不带给他的……

当然,许多人为Deligne的著作所折服,并非没有理由。但我也清楚:这部著作,在其令人惊叹的最初冲劲(以Weil猜想的证明告终)之外,离使出“全副本领”还很远。诚然,它见证一种少有的技术驾驭与从容,把他置于“最好者”之列。但它没有我年少时在他身上知觉到的那份寒微的德性——革新的德性。他身上曾怀有的这份德性、那份幼童的鲜活或天真,久已深深埋没、遭否认。我本要写:凭着这份“德性”与其不凡的天赋\dmark{}\origpage{382},也凭着他的天赋之舒展所曾受惠的那些特殊情势,Deligne本被召唤去“支配”我们时代的数学,一如Riemann、一如Hilbert,各支配其时代的数学。一些扎根于日常语言的积习思路,向我提示了“支配”这一意象;但它给出的,是对实在的一幅失真的把握。这些伟人当然曾完完全全“把握”“消化”“据为己有”其时代的已知数学;这无疑也给了他们一份对技术手段的出众驾驭。但倘若他们名副其实地显得“伟大”,那不是凭其技术的壮举——从一块顽石上“硬凿”艰难的证明。而是凭各自在数学的几个重要部分所带来的革新,经由简单而丰产的“想法”——也就是说:因为曾把目光投向那些简单而本质的事物——那些在他们之前无人肯垂顾的事物。那种\textbf{看见}简单而本质之事物的孩童般的能力——无论它们多么寒微、多么为众人所不屑——革新的力量、每个人身上创造的力量,正寓于其\textbf{中}。这份力量,在我认识的那个年轻人身上——一个无人识得、谦卑而热忱地爱着数学的人——以一种罕见的程度在场。岁月流转,这份寒微的“力量”似乎已从那位受人仰慕与畏惧、无碍地享有其威望、以及威望赋予他的(有时是独断的)对他人之权力的数学家身上,消失了。

这份在我朋友身上对一件极精微、极活泼、为众人所忽视而有创造之力的东西的\textbf{窒息},自我离去以来我已感到过多次;近年来尤甚。但需要这些星期的发现,以及我自3月底以来持续的反思(循着《收获与播种》的冲劲),才开始在其全部幅度上感到这份窒息在我朋友生活中的毁灭性效果——也在许多我曾亲近认识的其他人身上。不只对我“1970以后”的学生(及类似者)——他们领教过他的恶意(某些情形里,也许是无意识的),它对着每一个人施展,沉重地压在三人身上;而且——我现在隐约看到——也在我“1970以前”的学生中间:经由一道\textbf{旨趣上的连续性}的毁灭,以及整体感、统一感的毁灭——那份本赋予其工作以比一堆署着他们名字的抽印本的堆积更深、更广的意味的东西(注91)\fnote{125}{(*)}{(6月16日)这第二个侧面,只是在葬礼的反思中才向我显形。倘若我有幸见到一位显赫的数学家动用“令人气馁的权力”,那恰恰是在不久前还俨然我天造地设继承人的那位身上。写“令人气馁的权力”一节时,我想到他有过许多(在反思转回我自己之前),却还不曾有一丝疑心(至少在意识的层面上):这份权力,竟在那些人中间找到施展的机会——对他们,他本该(一如当年对我)扮演完美数学家的典范……}。

最近七年里,不止一次;最近这些星期、这些天里,又是不止一次,我在一种于某个层面上被感受为巨大的\textbf{糟蹋}的东西面前,感到一阵悲伤——当人身上与别人身上最珍贵的东西,仿佛刻意地被挥霍或窒息。然而,我终究也懂得:这样的“糟蹋”是人的境况的一个基调,它以这样或那样的形态,无处不在——在最寒微与最显赫的人们的生活里,一如在各族各邦的生活里。这份“糟蹋”——它不是别的,正是冲突、分裂在每个人的生活中的行动——是一种我刚刚才开始探其边的丰富而深刻的实体,一份属于我去“吃”去消化的食粮。由此,这份糟蹋,以及我在每一步遇到的别的糟蹋、以及每在路转弯处来到我面前、常常不合时宜的任何事物——这份糟蹋与这些不合时宜之物,身上都带着一份\textbf{益处}。倘若禅观有意义、有革新之力,那正是在这样的限度内:它使我得以领受那份(经我积习的折光)显为“为害”之物的益处,使我得以\textbf{以}那看似为毁灭而造的东西为食。

以自己的经历为食,任它更新自己而不总是绕开它——这就是充分承担自己的生活。我身上有这份力量;每一刻,用它或弃之于闲置,都由我。对我的朋友皮埃尔,对曾是我学生的每一个人,亦然——同我一样自由,以这份“糟蹋”为食;这份糟蹋,我在这场漫长禅观的最后几天里,正巡视完毕。对读到这些为读者而写的文字的人,亦然。

\subsection{共同继承人}

\notehead{91}{}

\dmark{}(5月19日)从四面八方传到我这里的、关于我昔日学生们的回音,稀薄之至。\origpage{384}我离去之后,几乎没有人肯给我一点活着的音讯——哪怕寄一份抽印本\fnote{126}{(*)}{(5月31日)对此可参见“沉默”一注之后的注n\textsuperscript{o}~84\textsuperscript{1}。}。然而,把我得到的这点东西聚拢起来,我可以形成一个概念——诚然极其粗略。也许在随后的几个月里它会明确起来,倘若这场反思促使他们中某些人现身。

我已有机会确认Deligne著作在我离去之后的深刻断裂,尽管在若干方面,他违背本意地显出接班人的样子,因而显出某种连续性。我曾感到这道断裂,必定深深波及我其余所有学生的工作。这正是我想稍稍接近地勾勒的印象。

这些学生中唯一一个其工作显然(至少乍看)延续了他与我一起所做的工作的,似乎是Berthelot\fnote{127}{(**)}{从Verdier在我离去后几年间所追随的对偶主题——在一个与我发展它的语境相近的解析语境——看,有一种如Berthelot情形的连续性印象。但我觉得,那多少是一种“例行公事的连续性”;而我主要寻找其征兆(或征兆之缺席)的,是一种创造性的连续——在未知中延续一股最初的冲劲……}。他也是长年来唯一不断给我寄许多抽印本的——甚至也许全部抽印本。它们都处在晶体上同调这个艰难的题目里;其系统的启动,是他论文的对象。然而我觉得,一如我其他“上同调学家”(交换的)学生,其著作同样带着对我引入的某些主要观念的冷淡:导出范畴(及Verdier提炼的三角范畴),六种运算的形式主义,topos(注91\textsuperscript{1})。如Zoghman~Mebkhout本人所说,他自己的著作——题材与Berthelot的如此贴近(注91\textsuperscript{2})——正处在这条线上:这些观念,连同Sato学派的观念。倘若它们不曾被我的上同调学家学生们——Deligne与Verdier领衔——否认,那么很可能早在七十年代初,Mebkhout的晶体理论(他1975年才着手发展,而且顶着这同一批学生的漠然)就已抵达六种\dmark{}\origpage{385}运算形式主义的完全成熟——而它至今仍未达到\fnote{128}{(*)}{(6月7日)贸然下这一评判,我曾有一番踌躇:它可能被解读为贬低Mebkhout理论的原创性。那与我的想法绝不相符;更何况,我对我的每一位上同调学家学生的手段都有极佳的评价(当他们不被外在于健全数学意识的成见所堵塞时)。我的朋友Zoghman本人,消除了我可能有的这份顾忌:他自己深信,“正常地”说,早在七十年代初,本该是我的学生们来发展他的理论。在某个层面上,他们自己想必是最先深信这一点的:那理论的本该的作者,是他们,或Deligne——而世风的普遍败坏再一助阵,便足够他们表现得俨然自己(或Deligne)确实就是作者!对此可参见“研讨会”与“欺瞒”两注,n\textsuperscript{o}~75$'$与85$'$。}。

我倒是记得曾向Verdier谈起那个使我着迷的问题:离散可构造系数与连续系数的关系;而毫无引起他兴趣的样子。这想必随后使Deligne着了迷,因为他用一整年(1969年)主持一场研讨班,以建立一部词典;那词典想必不曾使他满意,因为他随后把它弃之不顾。(见“当值无名之辈与好上帝定理”一注,n\textsuperscript{o}~48$'$。)况且此后,他被自己的下葬综合征“堵”到这种地步:直到1980年10月,他都未曾知觉Mebkhout工作的重要性——而当他终于觉察时,是以人们所知道的那副掘墓心境(见注n\textsuperscript{o}~75至76)。

据我所知,Verdier自答辩以来的著作,其实质仅限于在解析语境(它有时呈现额外的技术困难)里重做我在概形-凝聚框架里所做的事,不引入任何新想法。这甚至相当离奇:凭他分明发展出的反射、凭他消息如此灵通,他竟没有靠自己摇动手柄而撞上Mebkhout的理论——而且竟至少不曾认出:他的“学生”正在做些依我看有趣的事情,而且是他所错过的(一如Deligne所错过的)。

老实说,尽管我对离散系数与连续系数之关系的问题着迷,我却当真不曾疑心Mebkhout的晶体理论——它将在我离去后的十年间绽放。相反,有一个宏阔的主题,出自我五十年代(1955--1960)交换与非交换\dmark{}\origpage{386}上同调的反思;它仅仅被着手过(在“交换”的侧面,即以加性范畴的语言):在Verdier的工作里——六十年代初启动,答辩之后被弃置(见注n\textsuperscript{o}~81)。非交换的侧面,更晚在Giraud的论文里着手:他发展一种几何语言,以topos上的1-叠,表述维数$\leq 2$的非交换上同调。自六十年代下半叶起,这两道起跑线的不足就已显而易见:一边,“三角范畴”观念(Verdier提炼的)不足以表达与一个导出范畴相系的丰富的结构(该观念注定被\textbf{求导者}这个丰富得多的观念取代);另一边,需要为任意维数的非交换上同调发展一套几何语言,以topos上的n-叠与$\infty$-\textbf{叠}表述。人们(或我)感到需要这两条进路的综合;它将为同调代数与同伦代数充当共同的概念基础。这样的工作,也与Illusie的论文工作直接连续——两个侧面在其中都有代表。

经由求导者观念(非交换与交换框架皆有效),Bousfield--Kan论同伦极限的奠基工作(Lecture~Notes n\textsuperscript{o}~304,1972年出版),也处在这条弥散的纲领的脉络里——至少自1967年起,它只待人手来发展。去年1月——我尚不知道一个月后我将投身《追寻栈》——我把关于同伦型的“积分”的反思(同伦论者熟悉其名:“同伦(归纳)极限”)呈给Illusie;那时我还全然不知道Bousfield与Kan的工作的存在,也不知道这类运算已被别人考察过。显明的是,Illusie同样不知——尽管他被认为自1970年我“逝世”以来,始终泡在同调-同伦的水域里!由此可见,他与某些自然地属于奠基反思、属于他六十年代亲身推进的那条路线的实在,似乎失联到了何等地步\fnote{129}{(*)}{同伦型的“积分”这一观念,是在分层结构的剖视的语境里、于1981年底重新强加于我的。}。他想必给自己挖了个小窝,不再常出来……

\dmark{}\origpage{387}随着对topos观念本身、对全部“范畴的无意义”的轻蔑,毫不奇怪:Giraud如今对他头一项大工作主题已全然冷淡。诚然,Deligne两年前开棺出土Motif之际,曾装作忽然发现非交换上同调武库——gerbe、lien等——的趣味,仿佛是他刚刚引入它们,连同Motif与MotifGalois群\fnote{130}{(*)}{见“一个梦的回忆……——或Motif的诞生”一注,n\textsuperscript{o}~51。}。这种马戏,多半无力重燃他自己拼命扑灭的火焰……去年2月,我曾把一封二十页的信的抄件寄给Giraud;它成了《追寻栈》的开卷第一章。那是一场绝无技术性的反思;在其中,我得以“一跃而过”当年曾使Giraud(以及许多别人)却步的“炼狱”——即“非严格”n-范畴(我如今称之为n-叠)的观念:它始终停留在启发的层面,却分明是基本的。那就是《追寻栈》的起点。去年12月,在Contou-Carrère答辩之际,我们相遇(双方心境全然友好);我从Giraud那里得知:他连读这封信的好奇心都不曾有!我得到的印象是:他对这类东西已经一笔勾销。在那里可能有一片丰饶的实体——在他久已放弃的方向上——这个念头,似乎连掠过他心头都不曾。我试图让他明白——恐怕未能成功——那里有一项多汁而宏阔的工作,近二十年来等着有人做;而我到了暮年才终于着手,以至少在总体上勾出一个草图——听事物本身的口授——那片丰饶的实体;我这个“亡者”仍强烈地感到它,而我的学生们早已把它忘了。

Jouanolou同样放弃了一条研究方向——他刚以论文起手的方向。这一方向,成了由恰恰向他提供其选题之关键技术钥匙的那个人所倡导的时尚的轻蔑对象。三年前“反常”研讨会上围绕三角范畴的“抢滩”之际,同一位Deligne忽然(面不改色地)装作发现眼前那项宏大的奠基工作——其缺失忽然处处可感,而他正是十年来头一个劝阻它的人\dmark{}\origpage{388}——这样一项工作的需要,对我来说,早在1963/64年平展上同调起步时就已分明;对Deligne,也丝毫不差,自他开始听说$\ell$-进上同调与三角范畴之时——即次年他登上我的研讨班之时。问题在于:在$\mathbb{Z}_\ell$上(比如说,一个基底概形之上)构造“可构造三角范畴”、并在这一框架里发展“六种运算”的形式主义(这一步,我记得在Jouanolou的论文里已告完成)之外,还要做一项类似的工作:把基底环$\mathbb{Z}_\ell$换成一个多少任意的诺特$\mathbb{Z}_\ell$-代数,比如$\mathbb{Q}_\ell$或$\mathbb{Q}_\ell$的一个(代数?)扩张。这属于那些二十年来时机已熟、却始终等着有人去做的事情——且待吹在它们身上的轻蔑之风平息……

Raynaud夫人的工作的自然延续(平展上同调中的若干弱Lefschetz定理,以1-叠的语言),本会置于一个严格禁忌的$\infty$-叠语境里——不必提它!Sinh~夫人的工作亦然——1968年开始,1975年才完成——其自然延续,本会是一个所谓“monomiale”范畴的包裹$\infty$-范畴观念,或这类范畴的三角变形\fnote{131}{(*)}{见“一个梦的回忆……——或Motif的诞生”一注,n\textsuperscript{o}~51。}——不必想它!另一个是把她的工作转译成topos上叠的语言——多可怕!至于Monique~Hakim,她也不幸把论文做在一个题目上——自我不合时宜的离去以来,这题目如今在边缘上显得有点可笑:局部带环的topos上的相对概形,您说说看!她关于这一主题的小书,刊于《Grundlehren》(Springer),想必每年卖三四册——难怪我在那家名声不好,他们对接受我可能推荐的文本不再热心。对我而言,那是为一切“绝对”的“簇”观念(代数的、解析的等等)向一般“基底”做“相对化”的第一步试验——其需要,在我看来是明摆着的(注91\textsuperscript{3})。人们会说:直至今日,没有它也过得很好。但同样真确的是:人在这世上两百多万年,没有做数学也过得很好。总之,Monique~Hakim——她做论文的动机与我为她提议该题的动机并不相同——肯定绝无念头与那样一个主题保持哪怕一点接触;那主题,一旦脱离有利共识的语境、或脱离一股顶风冒雨坚持一种执拗而可靠视野的思想,便不再能对她有任何意味。

至于Neantro~Saavedra~Rivano,他似乎已全然从流通中消失——在世界数学家人名册(以及一切官方的名册)上,我甚至找不到他名字的踪迹。可以肯定的是:他那过于范畴家的论文题目,在那些裁定何为严肃、何为不然的先生们那里,想必不会有好名声。依我看,这篇论文最自然的延续,不多不少,正是那幅“Motif的宏阔图景”——对这个学生较为谦逊的抱负而言,那主题分明嫌大。然而他到底得了那份意外的荣幸:不过两年前,亲眼见自己的论文被其中一位大人物从头到尾重做。(对此可参见“下葬——或新父亲”与“白板”两注,n\textsuperscript{o}~52与67.)

最终,在我的十二位“1970年以前”的学生中,我唯一不太清楚其工作中是否有一道或多或少的苛烈或深刻的\textbf{断裂}——相对于他们在我的接触中所追随的工作——的,是Michel~Demazure与Michel~Raynaud(注91\textsuperscript{4})。我所知道的只是:他们继续做数学,而且(就其出众的手段而言,想必在意料之中)属于我方才所称的“上流”数学界。

前面这段从有时极其单薄的材料出发的简短反思,当然在很大程度上是假设性的、极其粗略的。我希望被提到的人肯原谅我或许粗疏的判断之误;倘若他们肯就此向我示意,我将乐于更正。在这里,我再次意识到:每个人的情形,肯定不同于所有其他人,代表一个远比一个如我般遥远的人所能合理把握的实在复杂得多的实在,更遑论以几行文字表达。所有这些保留说尽,我却仍感到这场反思并非无用——至少对我而言:以若干具体事实,稍稍安放昨天显现的那份仍属弥散的印象(它想必在未成形的层面,多年以来已在场):即一道\textbf{断裂}——它在我离去后的翌日降临于我的许多学生身上,并且在个人的层面上,映照一所“学派”的突然消失——一夜之间;在那道学派受业岁月的关键年月里,他们必定感到自己是它的一部分\dmark{}\origpage{390}。

\notehead{91\textsuperscript{1}}{}

(5月22日)我刚读到 P.~Berthelot的一篇综述文章,论CIRM吕米尼研讨会(1982年9月6--10日)“p-adic分析及其应用”;题为“刚几何与特征p代数簇的上同调”(24页),勾勒综合Dwork--Monsky--Washnitzer上同调与晶体上同调的主要想法。晶体上同调(受Monsky--Washnitzer启发)的出发点观念(连同其名字本身),以及经由引入由刚解析空间构成的site使之完备的观念——这些我在六十年代引入的想法,已成了所有从事这一主题者的日课,以Berthelot为首;其论文正在于发展并充实这些出发观念中的某些。这并不妨碍:我的名字,无论正文还是参考文献,都严格缺席。这样,第四位殡仪人学生清晰在案。下一个轮到谁?

(6月7日)有一件事值得注意:自我引入晶体上同调的起步观念十五年有余、自Berthelot确立该理论对真光滑概形确是“对的”理论的论文十有余年之后,人们仍未抵达我所称的晶体上同调的“掌控”局面——堪比为平展上同调在SGA~4与SGA~5研讨班里发展的那一种。对一个含对偶现象的上同调形式主义的(第一阶)“掌控”,我指的不多不少,正是对一套六种运算形式主义的完全占有。虽然我还不够“在局中”、无以赏识晶体语境特有的困难,但我不会惊讶:这一相对停滞的主因,在于Berthelot及其他人对这套形式主义观念本身的冷淡;这使他们(一如Deligne对其仍处童年状态的Hodge理论)忽视那头一个本质的“台阶”——通往一个全然“成年”的上同调形式主义所必达的台阶\dmark{}\origpage{391}。想必正是同一类心境,也使他不曾察觉Mebkhout的观点对他自己研究的趣味。

附记:当我在这里谈“晶体上同调”,而语境是放弃真性的假设时(对一个“全然成年”的形式主义,这是必需的),不言而喻:人们所用的晶状site,其对象是并非纯无穷小、而是“合宜的”(带除幂的)拓扑代数型的“增厚”(带除幂)。对晶状site的这一原初形态之推广的需要(对我来说,原初形态只是“对的”晶体理论的第一次逼近),我一开始就很清楚;而Berthelot是(连同出发观念一起)从我这里学得它的。这一关联的书面影射,见《主题概要》第5条。

\notehead{91\textsuperscript{2}}{}

相当离奇的是:除我之外,似乎没有人觉察,未具名的Mebkhout的理论是晶体理论的一个本质的新侧面。我——一个近十五年来已全然“脱钩”于上同调的人——却觉察到了,一旦Mebkhout去年费心向我多少解释了他所做的事情。而当我在Illusie面前(当作理所当然地)提起这一点时,他的神气,仿佛把这看作把两件分明毫无关系的东西(D-模与晶体)做了一点“边缘离奇”的拉扯。然而我第一手知道:他有数学家的嗅觉;我的其他学生(在这个案里即上同调学家们,从Deligne起)也有——但我确认:在某些局面里,这嗅觉对他们不再有任何用处……我越想,越觉得离奇:在这样的气候里,Mebkhout到底还是做成了他的工作,没让他的数学嗅觉被其长辈们的全然不理解——他们高踞其上……——缴了械。

\notehead{91\textsuperscript{3}}{}

主要是从五十年代末我在Cartan研讨班上关于复解析空间理论之基础、以及“Teichmüller式水平模簇”之精确几何诠释的报告起,我才懂得了对迄今所用的“簇”的通行观念(代数的、实或复解析的、微分的——或此后,“温和拓扑”中的变形)做双重推广的重要性。其一是放宽定义,以容许任意的“奇异性”与“标量函数”结构层中的幂零元——以我的概形奠基工作为模型。\dmark{}\origpage{392}另一项推广,是向合宜的局部带环topos之上的“相对化”(诸“绝对”观念,经由取一个点的topos为基而得到)。这一概念工作,成熟已逾二十五年,在Monique~Hakim的论文里着手,至今仍等着有人接手。一个格外有趣的个案,是相对刚解析空间的观念:它使人得以把寻常的复解析空间与变剩余特征的局部域上的刚解析空间,看作同一个相对刚解析空间的“纤维”;一如相对概形观念(它终于入了俗)使人得以把定义在不同特征域上的代数簇相互联结。

\notehead{91\textsuperscript{4}}{}

Demazure的论文工作,一如Raynaud的,本质上使用他们在与我的接触中学得的一套纯熟的概形技术;但各自工作的本质观念,不属于“Grothendieck式”的行装——这使他们的工作有别于我第一期其他学生的工作。很可能这一情势的后果,是其著作中的一份连续性,免于“老师下葬综合征”效果造成的断裂。这未必意味着这一综合征不曾以别的方式触及这一位或那一位。三年前,Raynaud对Contou-Carrère相对局部Jacobi簇工作的态度,曾使我吃惊。所宣布的结果深刻、艰难,而且美不胜收,远超对“众所周知”之事的简单推广;与Cartan典型曲线理论有出乎意料的联系;有漂亮的显式公式——全然在Raynaud的射程之内(也在我的射程之内)。他接待的那份冷淡,想必对Contou-Carrère的战略撤退起了决定性的分量:他放弃了那个他曾毫无保留地投入、而看来只会给他带来麻烦的题目……\fnote{132}{(*)}{细节见“棺材3——或有点太相对的Jacobi簇”一注(n\textsuperscript{o}~95)的子注n\textsuperscript{o}~95\textsuperscript{1}。}。我那封向他表示我(痛心地)惊讶于他对这些结果之美无动于衷的信,没有得到回音。

\subsection{……与电锯}

\notehead{92}{}

\dmark{}将近四年前,当我来到这一带安家时,离我家不远有一片漂亮的樱桃园。\origpage{393}我散步时常绕过去看看。我喜欢看那些茁壮的樱桃树——正当盛年,树干粗壮,仿佛自古就与这片土地长成一体;野草在树下自由滋蔓。它们想必没见过化肥,也没见过农药;樱桃时节,我就到那里去摘有味道的樱桃。树,想必有二三十棵。

有一天我再去,看见所有的树干都在人高的地方被锯断,树冠瘫倒在树干旁的地上,茬口朝天——一幅大屠杀的景象。用一把好电锯,想必顷刻而成,至多个把钟头。这样的场面我从未见过——人们锯树,一般总肯弯下腰,贴着地面锯。樱桃滞销,这我同意;这片园子想必也产不出几吨,这没问题——但这些树茬诉说的,可不是滞销与收成……

昨天,我又有了这同样的感觉:一根强劲的树干,根系深厚,树汁丰沛,枝条繁壮、延续着它的冲劲——被齐人高、干净利落地锯断,仿佛纯粹为了取乐。正是费心把主枝一根一根看过、看清每一根都被锯断,才终于使我看清发生了什么。那本该循着一股冲劲、循着一种根系深处内在的必然而舒展的东西,被一刀切断,不加手续,好让它在众人眼前被指定为笑柄。

这使我想起Zoghman谈起的那个“误会”:它想必发生在我的学生们(Deligne除外)与我之间。分明的是:冲劲与视野,都不曾从我传到我的任何一位学生身上(Deligne除外——他分明“除外”,名副其实!)。每一位都消化了一套技术行装——对他所选的题目做出色的工作有用(甚至不可或缺),日后兴许还用得上。是否有别的什么东西的起头——超越于此——我说不好。即便有起头,它面对电锯也绝无机会:电锯顷刻把它刈平……

我知道,只要还有人做数学\dmark{}\origpage{394}——而且除非彻底放弃人们做了两个多千年的那一类数学——他们早晚会不由自主地让我眼前这些僵卧的树枝一一复活。其中一些,已经被我的电锯朋友据为己有;倘若上天借他余年,很可能他还会照样处置另一些,甚至全部。但大多数,已全然不合他的风格。不过,也许他终究会厌倦于不断取代别人——那肯定极其累人,而且无利可图——而满足于做他自己(这已经不错了)。
